\section{Geometric steps I: Hodge--Tate, \'etale and crystalline steps}\label{sec:pf:steps1}

% ---------------------------------------------------------------------------
% local macros (prefix stepsI)
\providecommand{\stepsIBcris}{B_{\mathrm{cris}}}
\providecommand{\stepsIBdR}{B_{\mathrm{dR}}}
\providecommand{\stepsIBst}{B_{\mathrm{st}}}
\providecommand{\stepsIAcris}{A_{\mathrm{cris}}}
\providecommand{\stepsIDcris}{D_{\mathrm{cris}}}
\providecommand{\stepsIDst}{D_{\mathrm{st}}}
\providecommand{\stepsIDdR}{D_{\mathrm{dR}}}
\providecommand{\stepsIH}{\mathrm{H}}
\providecommand{\stepsIHQ}{\mathrm{H}_{\mathbb{Q}}}
\providecommand{\stepsIHl}{\mathrm{H}_{\ell}}
\providecommand{\stepsIHtot}{\mathrm{H}_{\mathrm{tot}}}
\providecommand{\stepsIHdg}{\mathrm{Hdg}}
\providecommand{\stepsIPiece}{\mathrm{Piece}}
\providecommand{\stepsIhop}{\mathrm{hop}}
\providecommand{\stepsICl}{\mathbb{C}_\ell}
\providecommand{\stepsIHk}{P^{\mathrm{Hk}}}
\providecommand{\stepsIrt}{P^{\mathrm{rt}}}
\providecommand{\stepsINw}{\mathrm{N}w}
\providecommand{\stepsIOw}{\mathcal{O}_{(w)}}
\providecommand{\stepsISpc}{\mathsf{Spc}}
\providecommand{\stepsISch}{\mathsf{Sch}}
\providecommand{\stepsISh}{\mathsf{Sh}}
\providecommand{\stepsIBGL}{\mathrm{BGL}}
\providecommand{\stepsIPGL}{\mathrm{P}}
\providecommand{\stepsItw}{\mathbb{Z}/q(i)}
\providecommand{\stepsIsch}{\mathrm{sch}}
\providecommand{\stepsIsp}{\mathrm{sp}}
\providecommand{\stepsIchar}{\chi_{\mathrm{cyc}}}
% ---------------------------------------------------------------------------

This section proves eight results of the formal development \cite{LeanSource2026}. Three of them
are \emph{step lemmas} about $\ell$-adic Galois representations cut out by Hecke eigensystems from
the \'etale cohomology of smooth projective varieties over number fields; together with the other
step lemmas they are assembled in \lref{UpperTheory.Gal.Int.Geo.steps_holds} (which states that all
step lemmas hold) into the geometric realization theorem (\cref{thm:geoone}). They are:
\begin{itemize}[nosep]
\item $\sfname{CrysStep}$ (\cref{pf:steps1:crys}): at a place $w\mid\ell$ where the tower has a
  smooth projective integral model, the representation is crystalline;
\item $\sfname{JumpStep}$ (\cref{pf:steps1:jump}): at such a place, every degree $b$ in which the
  Hecke eigensystem occurs in Hodge cohomology $\stepsIH^c(Z,\Omega^b)$ is a jump of the Hodge
  filtration (a Hodge--Tate weight) of the representation;
\item $\sfname{PhiStep}$ (\cref{pf:steps1:phi}): if moreover the model carries correspondences
  satisfying a pointwise congruence relation, the eigenvalues of a geometric Frobenius on the
  Weil--Deligne representation are roots of the Hecke polynomial.
\end{itemize}
The other five results are foundational inputs that the development proves rather than assumes:
the proper base change theorem (\cref{pf:steps1:E2}), the projective bundle formula for vector
bundles over simplicial presheaves (\cref{pf:steps1:D1}), the finiteness of \'etale cohomology of
proper varieties, derived from the other properties of \'etale cohomology used by the development
(\cref{pf:steps1:E4}), the vanishing of $\stepsIH^1$ of the generic fibre of a strictly henselian
local ring of the relative affine line over an absolutely integrally closed base
(\cref{pf:steps1:genfib}), and the existence of the comparison pairings used in the formal proof of
the crystalline comparison theorem (\cref{pf:steps1:pseam}).

Several of these results are classical theorems: proper base change \cite{ArtinGrothendieckVerdier1972},
the projective bundle formula \cite{Grothendieck1958,Grothendieck1977}, finiteness
\cite{ArtinGrothendieckVerdier1972,Deligne1977}, the crystalline and Hodge--Tate comparison theorems
\cite{Faltings1988,Faltings1989,FontaineMessing1987,Tsuji1999,BhattMorrowScholze2018}. For these we
state precisely the form in which the development proves them, deduce that form from the literature,
and indicate the route of the formal proof in remarks. What is specific to the development is
converted: the deduction of the three step lemmas from the comparison theorems (in particular the
Hodge--Tate spectral step \cref{pf:steps1:loce}, which transports Hecke eigenclasses in Hodge
cohomology to Hodge--Tate eigenvectors of the Galois representation), the reduction of the generic
fibre statement to Kummer theory and to the local acyclicity of the relative affine line, the origin
case of the triviality of torsors (\cref{pf:steps1:origin}), the algebraic proof that top powers of
hyperplane classes do not vanish (\cref{pf:steps1:powK}), the injectivity and bijectivity of the
Hodge--Tate pairing of the development (\cref{pf:steps1:htinj,pf:steps1:htbij}), the crystalline
congruence relation (\cref{pf:steps1:finall,pf:steps1:top,pf:steps1:wzero}), and the precise shape of
the intermediate statements. Long technical lemmas internal to the development are cited by their
formal labels and deferred.

\emph{Organisation.} \Cref{sec:pf:steps1:conv} fixes conventions: \'etale cohomology as defined in
the development, Hecke towers and their integral models, and the local conditions at places above
$\ell$. \Cref{sec:pf:steps1:etale} treats the \'etale results, \cref{sec:pf:steps1:cris} the
crystalline comparison inputs and $\sfname{CrysStep}$, \cref{sec:pf:steps1:jump} $\sfname{JumpStep}$, including
the Hodge--Tate comparison of the development (\cref{sec:pf:steps1:htdev}), and \cref{sec:pf:steps1:phi}
$\sfname{PhiStep}$ with the crystalline congruence relation.

% ===========================================================================
\subsection{Conventions}\label{sec:pf:steps1:conv}

\subsubsection{\'Etale cohomology in the development}\label{sec:pf:steps1:conv:et}

\begin{definition}[Spaces and cohomology]\label{pf:steps1:def:coh}
Let $\stepsISch$ be the category of schemes (in a fixed universe) with the big \'etale topology, and
let $\stepsISpc$ be the category of simplicial presheaves of sets on $\stepsISch$, called
\emph{spaces}. A scheme $X$ defines the constant simplicial presheaf $\stepsIsch(X)$ on the
representable presheaf $h_X$, and a ring $A$ defines $\stepsIsp(A)=\stepsIsch(\Spec A)$. For a
commutative ring $R$ let $\stepsISh_R$ be the category of \'etale sheaves of $R$-modules on
$\stepsISch$. For a space $E$ let $C_R(E)$ be the chain complex of the simplicial sheaf $R[E]$ of
free $R$-modules on $E$, and for $F\in\stepsISh_R$ fix an injective resolution $I(F)$. Then
\[
 \stepsIH^n(E,F)=\Hom_{K(\stepsISh_R)}\bigl(C_R(E),\,I(F)[n]\bigr)\qquad(n\in\Z)
\]
is the group of homotopy classes of degree-$n$ cochain maps; it is contravariant in $E$ (pull-back
$f^*$) and covariant in $F$. A map $f\colon E\to E'$ is an \emph{$F$-equivalence} if
$f^*\colon\stepsIH^n(E',F)\to\stepsIH^n(E,F)$ is bijective for all $n$.
\formal{GlobalTheory.EtC.Spc} \formal{GlobalTheory.EtC.H} \formal{GlobalTheory.EtC.Eqv}
\end{definition}

\begin{definition}[Tate twists]\label{pf:steps1:def:tw}
Let $q\ge1$ and $i\in\Z$. For a ring $A$ let $\mathrm{Prim}_q(A)$ be the set of roots of the
cyclotomic polynomial $\Phi_q$ in $A$ if $q\in A^\times$, and $\emptyset$ otherwise. The sheaf
$\stepsItw$ is the \'etale sheafification of the presheaf
$U\mapsto(\Z/q)[\mathrm{Prim}_q(\Gamma(U,\OO_U))]/\langle[\zeta^a]-a^i[\zeta]:a\in(\Z/q)^\times\rangle$,
and $\stepsIH^n(E,\stepsItw)$ denotes the cohomology with coefficients in it, as a sheaf of
$\Z/q$-modules. A space $E$ is \emph{$q$-invertible} if $q\in\Gamma(U,\OO_U)^\times$ for every scheme
$U$ with $E_0(U)\ne\emptyset$. The development defines cup products
$\stepsIH^n(E,\Z/q(i))\times\stepsIH^m(E,\Z/q(j))\to\stepsIH^{n+m}(E,\Z/q(i+j))$, a unit class, and, for
a line bundle $L\colon E\to\stepsIBGL_1$ (\cref{pf:steps1:def:bgl}), a first Chern class
$c_1(L)\in\stepsIH^2(E,\Z/q(1))$, chosen by Hilbert's $\varepsilon$ among the classes represented by
cocycles whose components in simplicial degrees $0,1,2$ are given by Kummer classes of $q$-th roots of
the transition functions of $L$; $c_1(L)^k$ denotes the $k$-th cup power.
\formal{GlobalTheory.EtC.tw} \formal{GlobalTheory.EtC.Ht} \formal{GlobalTheory.EtC.Inv}
\formal{GlobalTheory.EtC.cup} \formal{GlobalTheory.EtC.c1}
\end{definition}

\begin{lemma}[Comparison with classical \'etale cohomology]\label{pf:steps1:compare}
Let $X$ be a scheme and $F\in\stepsISh_{\Z/q}$. Then $\stepsIH^n(\stepsIsch(X),F)$ is canonically
isomorphic to the \'etale cohomology $\stepsIH^n(X_{\mathrm{et}},F|_{X_{\mathrm{et}}})$ of the
restriction of $F$ to the small \'etale site of $X$, compatibly with pull-backs. If
$q\in\Gamma(X,\OO_X)^\times$, then $\stepsItw|_{X_{\mathrm{et}}}\cong\mu_q^{\otimes i}$.
\end{lemma}

\begin{proof}
The complex $C(\stepsIsch(X))$ is the chain complex of a constant simplicial sheaf, hence
quasi-isomorphic to $(\Z/q)[h_X]$ placed in degree $0$; so $\stepsIH^n(\stepsIsch(X),F)
=\mathrm{Ext}^n((\Z/q)[h_X],F)$, which is the cohomology of the localized site
$(\stepsISch/X)_{\mathrm{et}}$. Restriction from the big to the small \'etale site of $X$ preserves
injectives and induces isomorphisms on cohomology \cite[Tag~0757, Lemma~0758]{Stacks2024}. For the
second statement: \'etale locally on $X$ there is a primitive $q$-th root of unity $\zeta$; then on each
connected component $\mathrm{Prim}_q=\{\zeta^a:a\in(\Z/q)^\times\}$, the quotient is free of rank one on
$[\zeta]$ with $[\zeta^a]=a^i[\zeta]$, and $[\zeta]\mapsto\zeta^{\otimes i}$ (for $i<0$: the dual basis
vector) defines an isomorphism with $\mu_q^{\otimes i}$ independent of $\zeta$; these isomorphisms glue.
\end{proof}

By \cref{pf:steps1:compare}, statements about $\stepsIH^n(\stepsIsch(X),\Z/q(i))$ with $q$
invertible on $X$ are statements about the usual groups $\stepsIH^n(X_{\mathrm{et}},\mu_q^{\otimes i})$,
and this is how the classical theorems below are applied. We write $\stepsIH^n(X,\Z/q(i))$ for
$\stepsIH^n(\stepsIsch(X),\Z/q(i))$.

\begin{definition}[$\ell$-adic cohomology]\label{pf:steps1:def:ladic}
For a prime $\ell$ and a space $E$, $\stepsIHl^n(E)(i)$ is the $\Zl$-module of compatible families
$(x_a)_{a\ge1}$, $x_a\in\stepsIH^n(E,\Z/\ell^a(i))$, under the reduction maps (the naive inverse
limit), and $\stepsIHQ^n(E)(i)=\stepsIHl^n(E)(i)\otimes_{\Zl}\Ql$. For a $k$-scheme $Z$ and a field
extension $\Omega/k$, $\stepsIHtot(Z_\Omega)=\bigoplus_n\stepsIHQ^n(Z_\Omega)(0)$, with the action of
$\Aut_k(\Omega)$ by pull-back along the induced automorphisms of $Z_\Omega$.
\formal{GlobalTheory.EtC.Hl} \formal{GlobalTheory.EtC.HQ}
\end{definition}

\begin{definition}[Models]\label{pf:steps1:def:model}
Let $K\subset\Qlbar$ be a finite extension of $\Ql$ with ring of integers $\OO_K$. A \emph{model} $M$
is a scheme $M.X$ with a morphism $M.X\to\Spec\OO_K$; $M_K$ is its generic fibre,
$V_n(M)=\stepsIHQ^n((M_K)_{\Qlbar})(0)$ with its action of $G_K=\Gal(\Qlbar/K)$, and
$T_n(M)=\stepsIHl^n((M_K)_{\Qlbar})(0)$. The model is \emph{good} if it is projective over $\OO_K$ and
strictly semistable with respect to some uniformizer, and \emph{smooth good} if it is good and smooth
over $\OO_K$. It is \emph{connected of dimension $d$} if $M_K$ is connected and smooth of relative
dimension $d$. For a basis $b$ of $V_n(M)$, $\mathrm{mat}_n(b)\colon G_K\to\GL_a(\Ql)$ is the matrix of
the Galois action. \formal{CrystallineSemistableComparison.Model}
\formal{CrystallineSemistableComparison.Model.mat}
\end{definition}

\subsubsection{Hecke towers and their integral models}\label{sec:pf:steps1:conv:towers}

\begin{definition}[Hecke towers, pieces]\label{pf:steps1:def:tower}
A \emph{Hecke tower} $H$ over a field $k$ consists of an index set $\Gamma$, a $k$-scheme
$q\colon Z\to\Spec k$, and for each $g\in\Gamma$ a scheme $Y_g$ with $s_g,t_g\colon Y_g\to Z$ and
$q s_g=qt_g$. It is \emph{nice} if $q$ is smooth and projective and every $s_g$ is finite \'etale with a
constant number of lifts of geometric points. For a field embedding $\tau\colon k\to K$,
$\mathrm{bc}_\tau H$ is the base change. For an extension $\Omega/k$, the Hecke operator
$\stepsIhop(g)$ on $\stepsIHtot(Z_\Omega)$ is \emph{some} linear map which is an \emph{operator of the
correspondence} $(s_g,t_g)_\Omega$, i.e.\ for which there are a finite \'etale Galois cover $c\colon W\to
Z_\Omega$ and sections $\sigma_1,\dots,\sigma_J\colon W\to Y_{g,\Omega}$ splitting $s_g$ over $c$ with
$c^*\stepsIhop(g)(x)=\sum_j(t_g\sigma_j)^*x$; it is chosen by Hilbert's $\varepsilon$, and is $0$ if there
is none. For a field $L\supseteq\Ql$ and $\theta\colon\Gamma\to L$, the \emph{$\theta$-piece} is the joint
eigenspace
\[
 \stepsIPiece_\theta(H,\Omega,L)=\bigcap_{g\in\Gamma}\ker\bigl(\stepsIhop(g)_L-\theta(g)\bigr)
 \subseteq L\otimes_{\Ql}\stepsIHtot(Z_\Omega).
\]
A homomorphism $\rho_0\colon\Aut_k(\Omega)\to\GL_d(L)$ is a \emph{matrix of the $\theta$-piece} if some
basis $b_1,\dots,b_d$ of the piece satisfies $\sigma^*b_j=\sum_i\rho_0(\sigma)_{ij}b_i$ (the action being
$L$-linear). For $\rho_0,\rho\colon\Gamma'\to\GL_d(L)$, $\rho$ is \emph{graded of} $\rho_0$ if for some
$c\colon\{1,\dots,d\}\to\N$ every $\rho_0(g)$ is block upper triangular for $c$ (entries $(i,j)$ with
$c(j)<c(i)$ vanish) and $\rho(g)_{ij}=\rho_0(g)_{ij}$ if $c(i)=c(j)$, $=0$ otherwise. For a number field
$E$ with algebraic closure $\Omega$, $G_E=\Aut_E(\Omega)$, and $\rho\colon G_E\to\GL_d(\Qlbar)$, $\rho$ is
\emph{attached to the $\theta$-piece} if it is graded of a matrix of the $\theta$-piece.
\formal{GlobalTheory.GalV.HTower} \formal{UpperTheory.Gal.Int.Geo.Nice}
\formal{UpperTheory.Gal.Int.Geo.bcAlong} \formal{GlobalTheory.GalV.hop} \formal{GlobalTheory.GalV.IsCorrOp}
\formal{UpperTheory.Gal.Int.Geo.piece} \formal{UpperTheory.Gal.Int.Geo.IsMat}
\formal{UpperTheory.Gal.Int.Geo.IsGr} \formal{UpperTheory.Gal.Int.Geo.IsPiece}
\end{definition}

\begin{definition}[Hodge eigenclasses]\label{pf:steps1:def:hodge}
For $q\colon Z\to\Spec k$ let $\stepsIH^b(Z,\Omega^a)$ be the sheaf cohomology of the sheaf of relative
$a$-forms, with pull-backs and, along finite \'etale maps, traces. For a tower $H$ over $k$ and
$\theta\colon\Gamma\to k$, a class $x\in\stepsIH^b(Z,\Omega^a)$ is a \emph{$\theta$-eigenclass} if every
$s_g$ is finite \'etale and $\mathrm{tr}_{s_g}(t_g^*x)=\theta(g)x$ for all $g$; $\theta$ \emph{occurs in Hodge
degree $a$} for $H$ if some $\stepsIH^b(Z,\Omega^a)$ contains a nonzero $\theta$-eigenclass.
\formal{GlobalTheory.HodgeCoh.hodgeH} \formal{UpperTheory.Gal.Int.Geo.IsHEig}
\formal{UpperTheory.Gal.Int.Geo.HOcc} \formal{UpperTheory.Gal.Int.Geo.HodgeAt}
\end{definition}

\begin{definition}[Covers]\label{pf:steps1:def:cover}
Let $H$ be a tower over $k$, $\theta\colon\Gamma\to L$ and $O\to k$ a ring map. A \emph{cover} $C$ of
$(H,\theta)$ over $O$ consists of a smooth projective $O$-scheme $u\colon Y\to\Spec O$ of some relative
dimension, a morphism $\pi\colon Z\to Y$ over $\Spec k\to\Spec O$ such that the induced
$\pi_k\colon Z\to Y_k$ is finite \'etale, and a finite group $\Delta$ acting on $Z$ by automorphisms
$a_\delta$ over $Y_k$ such that $\coprod_{\delta}Z\to Z\times_{Y_k}Z$ is an isomorphism and each $a_\delta$ is
$t_g$ for some $g$ with $s_g$ an isomorphism and $\theta(g)=1$. $C$ is \emph{onto} if $\pi_k$ is
surjective. For $s,t\colon T\to Y$ and $m\in\N$, $g$ \emph{lies under $(s,t)$ with multiplicity $m$} if
$m>0$ and there is $\kappa\colon Y_g\to T$ with $s\kappa=\pi s_g$, $t\kappa=\pi t_g$ and
$Y_g\to Z\times_YT$ finite \'etale of degree $m$. $C$ is \emph{good} if every $g$ lies under some finite
\'etale correspondence $(s,t)$ of constant degree on $Y$ over $O$. For a finite place $w$ of a number field
$E$, $\stepsIOw$ is the localization of $\OO_E$ at $w$ (not its completion) and $\stepsINw=\#\OO_E/w$.
\formal{UpperTheory.Gal.Int.Geo.Cover} \formal{UpperTheory.Gal.Int.Geo.Cover.Onto}
\formal{UpperTheory.Gal.Int.Geo.Cover.Lies} \formal{UpperTheory.Gal.Int.Geo.Cover.Good}
\formal{UpperTheory.Gal.Int.Geo.GoodAt}
\end{definition}

\begin{definition}[Congruence data]\label{pf:steps1:def:cong}
For a prime $p$ and a scheme $Y$, a \emph{$p$-correspondence} $T$ on $Y$ is a pair
$\alpha,\beta\colon T_Z\to Y$ with $\alpha$ finite, \'etale over $Y[1/p]$, and with exactly $d_T$ lifts of
every geometric point of $Y[1/p]$. For $\varpi\in\Gamma(Y,\OO_Y)$ and $p$-correspondences
$T_0,\dots,T_n$, the \emph{pointwise congruence relation for $(q,n,\varpi)$} holds if for every valuation
ring $V$ with algebraically closed fraction field and $p\ne0$ in $V$, every $y\colon\Spec V\to Y$ whose
closed point maps to a generic point of $\{\varpi=0\}$, every endomorphism $\phi$ of the residue field
$\kappa_V$ with $\phi(a)=a^q$, and all enumerations $(z_{j,i})_i$ of the lifts of $y$ along $\alpha_{T_j}$,
the identity of multisets of $\kappa_V$-points of $Y$
\[
 \sum_{j\ \mathrm{even}}q^{\binom j2}\bigl\{\phi^{n-j}(\bar v_{j,i})\bigr\}_i
 =\sum_{j\ \mathrm{odd}}q^{\binom j2}\bigl\{\phi^{n-j}(\bar v_{j,i})\bigr\}_i
\]
holds, where $\bar v_{j,i}$ is the reduction of $\beta_{T_j}\circ z_{j,i}$; this is the congruence relation
\[
 \textstyle\sum_j(-1)^jq^{\binom j2}\,\Fr^{n-j}\circ T_j=0
\]
read pointwise at generic points of the special fibre. A
cover $C$ over $O$ has a \emph{congruence datum} $(p,f,n,\vartheta)$ if there are an irreducible
$\varpi\in O$, $p$-correspondences $T_0,\dots,T_n$ on $Y$ over $O$ satisfying the pointwise congruence
relation for $(p^f,n,u^*\varpi)$, indices $\gamma_j\in\Gamma$ and integers $m_j$ such that $\gamma_j$ lies
under $(\alpha_{T_j},\beta_{T_j})$ with multiplicity $m_j$ and $\theta(\gamma_j)=m_j\vartheta_j$. We say that
$(H,\theta)$ has a \emph{congruence datum at $w$ with parameters $(n,\vartheta)$} if some onto cover over
$\stepsIOw$ has a congruence datum $(p,f,n,\vartheta)$ with $p\in w$ prime and $\stepsINw=p^f$.
\formal{UpperTheory.Gal.crystallineComparison.Corr} \formal{UpperTheory.Gal.crystallineComparison.PtId}
\formal{UpperTheory.Gal.Int.Geo.Cover.Cong} \formal{UpperTheory.Gal.Int.Geo.CongAt}
\end{definition}

\begin{definition}[Hecke polynomials]\label{pf:steps1:def:heckepol}
For $q,n\in\N$, $\vartheta\colon\N\to R$ and a multiset $S$,
$\stepsIHk_{q,n,\vartheta}(X)=\sum_{j=0}^n(-1)^jq^{\binom j2}\vartheta_jX^{n-j}$ and
$\stepsIrt_S(X)=\prod_{s\in S}(X-s)$. For $\iota\colon\Qlbar\xrightarrow\sim\C$, a place $w$ and a multiset
$R\subset\C$, \emph{$R$ is the root multiset of the Hecke polynomial of $\vartheta$ at $w$} if
$\vartheta_0\ne0$ and $\iota(\stepsIHk_{\stepsINw,n,\vartheta})=\iota(\vartheta_0)\stepsIrt_R$.
\formal{UpperTheory.Gal.crystallineComparison.heckePol} \formal{UpperTheory.Gal.crystallineComparison.rootPol}
\formal{UpperTheory.Gal.Int.Geo.HasRoots}
\end{definition}

\subsubsection{Local conditions above $\ell$}\label{sec:pf:steps1:conv:loc}

\begin{definition}[Period packages and bridges]\label{pf:steps1:def:package}
A \emph{period package} $P$ for $\ell$ consists of rings $P.B$ and $P.B_{\mathrm{dR}}$ with actions of
$G_{\Ql}$, a Frobenius $\varphi$ on $P.B$, an embedding $P.B\to P.B_{\mathrm{dR}}$, a filtration on
$P.B_{\mathrm{dR}}$ and Fontaine's period data; it is \emph{tied} if there are ring isomorphisms
$P.B\cong\stepsIBcris$ and $P.B_{\mathrm{dR}}\cong\stepsIBdR$ with the period rings constructed in the
development, compatible with the structure maps from $\Zl$. For a number field $E$ with algebraic
closure $\Omega$, a subgroup $G_w\subseteq G_E$ and a topological ring $O$, a \emph{bridge} $B_w$ consists of
a finite extension $K_w\subset\Qlbar$ of $\Ql$, an embedding $j\colon\Omega\to\Qlbar$ with $j(E)\subset K_w$
dense, a topological isomorphism $e\colon G_w\cong G_{K_w}$ compatible with $j$, and a finite extension
$LL\supseteq K_w$ containing all conjugates of $K_w$. If $G_w$ is the decomposition group $D_A$ of a
valuation ring $A\subset\Omega$ over $w$, $B_w$ \emph{induces} $A$ if $A=\{x:|j(x)|\le1\}$; then $w\mid\ell$.
We write $\tau_w\colon E\to K_w$ for the embedding induced by $j$.
\formal{Formalism.BasicNotions.LocPrime} \formal{Formalism.BasicNotions.LocTied}
\formal{Formalism.BasicNotions.LocBridge} \formal{GlobalTheory.Pst.Induces}
\end{definition}

\begin{definition}[Crystalline representations, jumps, Weil--Deligne]\label{pf:steps1:def:crys}
Let $r\colon G_w\to\GL_a(\Qlbar)$. Then $r$ is \emph{crystalline for $(P,B_w)$} if $r$ is the extension of
scalars of a representation $r'$ with values in $\GL_a(LL')$ for a finite extension $LL'\supseteq LL$ (a
\emph{model} of $r$), and the $LL'$-linear representation $V$ of $G_{K_w}$ defined by $r'\circ e^{-1}$
satisfies $\dim_{LL'}\stepsIDcris(V)=[K_{w,0}:\Ql]\dim_{LL'}V$, where $\stepsIDcris(V)=(V\otimes_{\Ql}P.B)^{G_{K_w}}$
and $K_{w,0}$ is the maximal unramified subextension. For $\tau\colon E\to\Qlbar$, a multiset $J\subset\Z$ is
a \emph{jump multiset of $r$ at $\tau$} if for some model $V$ is de Rham and, for an embedding
$\tau'\colon K_w\to LL'$ inducing $\tau$, $J$ is the multiset of jumps of the Hodge filtration on the
$\tau'$-part of $\stepsIDdR(V)$. Finally $\mathrm{WD}(r)$ is the Frobenius-semisimplified Weil--Deligne
representation over $\C$ (via $\iota$) of the Weil group $W_A\subseteq D_A$ attached to $r$ by Fontaine's
construction if $r$ is potentially semistable, and the trivial one otherwise; $\deg\colon W_A\to\Z$ is the
degree (elements of degree $-1$ are geometric Frobenius elements).
\formal{GlobalTheory.Pst.IsCrysAt} \formal{GlobalTheory.Pst.IsJumpsAt} \formal{GlobalTheory.Pst.wdAt}
\end{definition}

\begin{definition}[Conditions at $w$]\label{pf:steps1:def:conds}
Let $\rho\colon G_E\to\GL_d(\Qlbar)$ and $w$ a finite place. Write ``$\forall^{\mathrm t}$'' for ``for every
tied period package $P$, every valuation ring $A\subset\Omega$ over $w$ and every bridge $B_w$ for $D_A$
inducing $A$''. Then: $\rho$ is \emph{crystalline at $w$} if $\forall^{\mathrm t}$, $\rho|_{D_A}$ is crystalline for
$(P,B_w)$; $\rho$ satisfies the \emph{Frobenius condition with roots $R$ at $w$} (for a multiset $R\subset\C$)
if $\forall^{\mathrm t}$, $\rho|_{D_A}$ is crystalline for $(P,B_w)$ and there are $g\in W_A$ with $\deg g=-1$ and
$m\in\N$ such that the characteristic polynomial of $g$ on $\mathrm{WD}(\rho|_{D_A})$ divides $(\stepsIrt_R)^m$;
and, for a predicate $Q$ on $\Hom(E,\C)\times\Z$, $\rho$ satisfies the \emph{jump condition at $w$ for $Q$} if
$\forall^{\mathrm t}$, for all $\tau\colon E\to\Qlbar$, finite $LL'$ and $\tau'\colon K_w\to LL'$ inducing $\tau$, there is
a jump multiset $J$ of $\rho|_{D_A}$ at $\tau$ with $j\in J$ whenever $Q(\iota\tau,j)$.
\formal{UpperTheory.Gal.Int.Div.CrysAt} \formal{UpperTheory.Gal.Int.Geo.PhiK} \formal{UpperTheory.Gal.Int.Geo.JOcc}
\end{definition}

\begin{definition}[The three step propositions]\label{pf:steps1:def:steps}
Let $E$ be a number field with algebraic closure $\Omega$, $\ell$ a prime, $\iota\colon\Qlbar\cong\C$, $H$ a nice
tower over $E$, $\theta\colon\Gamma\to\Qlbar$, $\rho\colon G_E\to\GL_d(\Qlbar)$ continuous and attached to the
$\theta$-piece, and $w\mid\ell$ a finite place. The propositions assert, for all such data:
\begin{enumerate}[label=\textup{(\roman*)},nosep]
\item $\sfname{CrysStep}$: if $(H,\theta)$ has an onto cover over $\stepsIOw$, then $\rho$ is crystalline at $w$;
\item $\sfname{JumpStep}$: if $(H,\theta)$ has an onto good cover over $\stepsIOw$, then $\rho$ satisfies the jump
  condition at $w$ for the predicate ``$j=b\ge0$ and $\iota\theta$ occurs in Hodge degree $b$ for
  $\mathrm{bc}_{\iota\tau}H$'';
\item $\sfname{PhiStep}$: if $(H,\theta)$ has a congruence datum at $w$ with parameters $(n,\vartheta)$ and $R$ is the
  root multiset of the Hecke polynomial of $\vartheta$ at $w$, then $\rho$ satisfies the Frobenius condition with
  roots $R$ at $w$.
\end{enumerate}
\formal{UpperTheory.Gal.Int.Geo.CrysStep} \formal{UpperTheory.Gal.Int.Geo.JumpStep}
\formal{UpperTheory.Gal.Int.Geo.PhiStep}
\end{definition}

% ===========================================================================
\subsection{\'Etale foundations}\label{sec:pf:steps1:etale}

\begin{definition}[The propositions A1--E4]\label{pf:steps1:def:axioms}
The development organises the properties of \'etale cohomology it uses as the following list of
propositions about the groups of \cref{pf:steps1:def:coh}. Throughout, $q\ge1$, $i\in\Z$, and $F$ denotes an
arbitrary sheaf of modules.
\begin{description}[nosep,leftmargin=2.6em,style=sameline]
\item[A1] Quasi-isomorphisms of complexes of sheaves induce bijections on $\Hom_{K}(-,I(F)[n])$.
\item[A2] For $f\colon Y\to X$ smooth and surjective, the augmented \v{C}ech nerve
  $\check C(f)\to\stepsIsch(X)$ is an $F$-equivalence.
\item[A3] A map of spaces which is an $F$-equivalence in each simplicial degree is an $F$-equivalence.
\item[A4] (Mayer--Vietoris) If $p\colon E\to E'$ maps a cover of $E$ by two subobjects (whose free sheaves
  are jointly surjective in each degree) to such a cover of $E'$, and restricts to $F$-equivalences on the
  pieces and on their intersection, then $p$ is an $F$-equivalence.
\item[A5] Every class in $\stepsIH^n(\stepsIsch(X),F)$, $n>0$, vanishes on some \'etale surjective $Y\to X$.
\item[A7] Short exact sequences of sheaves give long exact cohomology sequences.
\item[A8] Restriction of scalars along a ring map induces bijections on cohomology.
\item[A9] $\stepsIH^n(E,F)=0$ for $n<0$, and $\stepsIH^0(\stepsIsch(X),F)=F(X)$ naturally.
\item[A10] A map of sheaves bijective on sections over all $U$ with $E_0(U)\neq\emptyset$ induces
  bijections on $\stepsIH^n(E,-)$.
\item[A11] The reduction maps $\Z/Q(i)\to\Z/q(i)$, $q\mid Q$, are transitive and compatible with cup
  products, units and $c_1$.
\item[B1] For a filtered colimit $A=\varinjlim A_j$ of rings with $q\in A^\times$,
  $\stepsIH^n(\stepsIsp A,\stepsItw)=\varinjlim_j\stepsIH^n(\stepsIsp A_j,\stepsItw)$.
\item[B2] If $A$ is strictly henselian local, $\stepsIH^n(\stepsIsp A,F)=0$ for $n>0$.
\item[B3] If $f\colon Y\to X$ is a quasi-compact quasi-separated morphism over a quasi-compact
  quasi-separated $X\to S$ with $q$ invertible on $S$, and the base change of $f$ to $\Spec A$ is a
  $\stepsItw$-equivalence for every strictly henselian local ring $A$ over $S$, then $f$ is a
  $\stepsItw$-equivalence.
\item[B4] If $f\colon Y\to X$ is integral and surjective and $q$ is invertible on $X$, the augmented \v{C}ech
  nerve of $f$ is a $\stepsItw$-equivalence.
\item[B5] If $A$ is an absolutely integrally closed domain (every monic polynomial of positive degree has a
  root) with $q\in A^\times$, then $\stepsIH^n(\stepsIsp A,\stepsItw)=0$ for $n>0$.
\item[B7] For a field $K$ with $q\in K^\times$, $\stepsIH^n(\stepsIsp K,\stepsItw)$ is isomorphic to the continuous
  Galois cohomology of $\mu_q^{\otimes i}$.
\item[C1] If $k$ is algebraically closed and $K/k(t)$ algebraic with $q\in K^\times$, then
  $\stepsIH^n(\stepsIsp K,\stepsItw)=0$ for $n\ge2$.
\item[C2] If $B$ is a Dedekind domain of finite type over an algebraically closed field $k$ with $q\in k^\times$,
  then $\stepsIH^n(\stepsIsp B,\stepsItw)=0$ for $n\ge2$.
\item[C3] If $T$ is an absolutely integrally closed local domain with $q\in T^\times$ and $B$ is an \'etale
  $T[x]$-algebra, then $\Spec(\mathrm{Frac}(T)\otimes_TB)\to\Spec B$ is a $\stepsItw$-equivalence.
\item[C4] (Homotopy invariance) $\Spec A[x]\to\Spec A$ is a $\stepsItw$-equivalence if $q\in A^\times$.
\item[C5] If $A$ is strictly henselian with $q\in A^\times$, then $\stepsIH^n(\stepsIsp A[t,t^{-1}],\stepsItw)=0$
  for $n\ge2$ and $\stepsIH^1\cong\Z/q$.
\item[D1] The projective bundle formula (\cref{pf:steps1:D1}).
\item[D2] Existence of power-sum Chern classes (\cref{pf:steps1:D2}).
\item[D3] Existence of slant products with group homology classes (\cref{pf:steps1:D2}).
\item[E1] (Affine analogue of proper base change) For a henselian pair $(A,I)$ with $q\in A^\times$,
  $\Spec(A/I)\to\Spec A$ is a $\stepsItw$-equivalence.
\item[E2] Proper base change (\cref{pf:steps1:E2}).
\item[E3] If $T$ is an absolutely integrally closed local domain with $q\in T^\times$ and $X\to\Spec T$ is
  smooth, then the generic fibre $X_{\mathrm{Frac}(T)}\to X$ is a $\stepsItw$-equivalence.
\item[E4] Finiteness (\cref{pf:steps1:E4}).
\end{description}
\formal{GlobalTheory.EtC.A1}\,\dots\,\formal{GlobalTheory.EtC.E4}
\end{definition}

With the exception of D1, D2, D3, E2 and E4, these propositions are proved in other parts of the
development. E1 is \lref{UpperTheory.Et.affine_proper_base_change}, which states E1 (it is Gabber's
affine analogue of the proper base change theorem \cite{Gabber1994}), and E3 is
\lref{UpperTheory.Et.etale_base_change_of_duality}, which states E3; \cref{pf:steps1:genfib} below is a
key input of the proof of E3.

\subsubsection{Proper base change}

\begin{theorem}[Proper base change; E2]\label{pf:steps1:E2}
Let $q\ge1$, let $A$ be a strictly henselian local ring with $q\in A^\times$, let $p\colon X\to\Spec A$ be
proper, and let $i\in\Z$. Then the inclusion of the closed fibre $X_0=X\times_{\Spec A}\Spec(A/\mathfrak m_A)\to X$
is a $\stepsItw$-equivalence: restriction $\stepsIH^n(X,\Z/q(i))\to\stepsIH^n(X_0,\Z/q(i))$ is bijective for
every $n\in\Z$. In particular E2 follows from the propositions A1--A5, A7--A10, B1--B5 and E1.
\formallabel{UpperTheory.Et.proper_base_change}
\formallabel{UpperTheory.Et.E2.main}
\end{theorem}

\begin{proof}
For $n<0$ both groups vanish (A9). For $n\ge0$, by \cref{pf:steps1:compare} the map is the restriction map
$\stepsIH^n(X_{\mathrm{et}},\mu_q^{\otimes i})\to\stepsIH^n(X_{0,\mathrm{et}},\mu_q^{\otimes i})$ of classical \'etale
cohomology, which is bijective by the proper base change theorem in the form
\cite[Exp.~XII, Cor.~5.5(iii)]{ArtinGrothendieckVerdier1972}: for $S$ the spectrum of a henselian local ring,
$f\colon X\to S$ proper and $F$ a torsion abelian sheaf, $\stepsIH^n(X,F)\to\stepsIH^n(X_0,F|_{X_0})$ is bijective for
all $n$. Apply it to the torsion sheaf $\mu_q^{\otimes i}$. Since the conclusion holds unconditionally, the implication
in the last sentence holds as well.
\end{proof}

\begin{remark}[The formal proof of E2]\label{pf:steps1:E2rem}
The development does not quote the classical theorem. Its proof of the implication in \cref{pf:steps1:E2} derives E2
from A1--A5, A7--A10, B1--B5 and E1: it reduces to projective morphisms by Chow's lemma, proves the statement for
towers built from finite morphisms and relative affine lines, using the cohomology of the punctured affine line over
a strictly henselian base, and compares with the closed fibre. The hypothesis $q\in A^\times$ is not needed in
\cite[Exp.~XII]{ArtinGrothendieckVerdier1972}; it is part of the formal statement.
\end{remark}

\subsubsection{The projective bundle formula}

\begin{definition}[Vector bundles on spaces]\label{pf:steps1:def:bgl}
For $N\ge0$ let $\stepsIBGL_N$ be the space $U\mapsto$ nerve of the one-object category
$\GL_N(\Gamma(U,\OO_U))$; a map $V\colon E\to\stepsIBGL_N$ is a \emph{rank-$N$ vector bundle on $E$} (its
components on $1$-simplices are transition matrices), and $V$ is \emph{trivial} if all transition matrices
are $1$. Let $\stepsIPGL_N$ be the space $U\mapsto$ nerve of the category whose objects are the unimodular
vectors $v\in\Gamma(U,\OO_U)^N$ (generating the unit ideal) and whose morphisms $v\to v'$ are the pairs
$(g,\lambda)\in\GL_N\times\mathbb G_m$ with $gv=\lambda v'$; it carries the maps
$\mathrm{pr}\colon\stepsIPGL_N\to\stepsIBGL_N$, $(g,\lambda)\mapsto g$, and $\mathrm{line}\colon\stepsIPGL_N\to\stepsIBGL_1$,
$(g,\lambda)\mapsto\lambda$. For $V\colon E\to\stepsIBGL_N$ the \emph{projective bundle} is
$\mathbb P(V)=E\times_{\stepsIBGL_N}\stepsIPGL_N$, with projection $\pi_V$ to $E$ and tautological line bundle
$\OO_V=\mathrm{line}\circ\mathrm{pr}_2$.
\formal{GlobalTheory.EtC.glCat} \formal{GlobalTheory.EtC.PGL} \formal{GlobalTheory.EtC.PGL.pr}
\formal{GlobalTheory.EtC.PGL.line}
\end{definition}

For a space $E$, a vector bundle $V$ and $n,i\in\Z$ consider the map
\begin{equation}\label{pf:steps1:eq:pbf}
 \bigoplus_{j=0}^{N-1}\stepsIH^{n-2j}(E,\Z/q(i-j))\longrightarrow\stepsIH^n(\mathbb P(V),\Z/q(i)),\qquad
 (a_j)_j\longmapsto\sum_{j=0}^{N-1}\pi_V^*(a_j)\cup c_1(\OO_V)^j .
\end{equation}

\begin{lemma}[Discrete spaces carry only trivial bundles]\label{pf:steps1:disc}
If $E$ is constant as a simplicial object, every vector bundle $V\colon E\to\stepsIBGL_N$ is trivial.
\formallabel{GlobalTheory.EtC.disc_eq_trivV}
\end{lemma}

\begin{proof}
All $1$-simplices of a constant simplicial object are degenerate, and a map of simplicial objects sends
degenerate simplices to degenerate simplices, i.e.\ to identity morphisms of $\GL_N$.
\end{proof}

\begin{lemma}[Projective spaces over schemes]\label{pf:steps1:pbfsch}
Let $U$ be a scheme with $q\in\Gamma(U,\OO_U)^\times$. For every $V\colon\stepsIsch(U)\to\stepsIBGL_N$ the map
\eqref{pf:steps1:eq:pbf} is bijective for all $n,i$.
\formallabel{GlobalTheory.T.etaleCohomology_projectiveBundle_formula}
\formallabel{UpperTheory.Cb.etaleCohomology_projectiveBundle_formula}
\end{lemma}

\begin{proof}
By \cref{pf:steps1:disc} $V$ is trivial, so $\mathbb P(V)$ is the nerve of the action groupoid of
$\mathbb G_m$ on the scheme $\mathbb A^N_U\smallsetminus0$ of unimodular vectors, i.e.\ the \v{C}ech nerve of the
smooth surjective $\mathbb G_m$-torsor $\mathbb A^N_U\smallsetminus0\to\mathbb P^{N-1}_U$; by A2 (a theorem of the
development, \lref{GlobalTheory.T.etC_a2}, stating A2) its cohomology is that of $\mathbb P^{N-1}_U$, and
$\OO_V$ is the tautological line bundle. By \cref{pf:steps1:compare} the groups and maps in \eqref{pf:steps1:eq:pbf} are
then those of classical \'etale cohomology with coefficients $\mu_q^{\otimes i}$, and the map is bijective by the
projective bundle theorem \cite[Exp.~VII, Th.~2.2.1]{Grothendieck1977}: for a locally free $\OO_U$-module $\mathcal E$ of
rank $N$ and $q$ invertible on $U$, $\bigoplus_{j<N}\stepsIH^{n-2j}(U,\mu_q^{\otimes(i-j)})\to
\stepsIH^n(\mathbb P(\mathcal E),\mu_q^{\otimes i})$, $(a_j)\mapsto\sum_j\pi^*a_j\cup\xi^j$, is an isomorphism, $\xi$ being
the first Chern class of the tautological line bundle. The sign convention for $c_1$ does not matter, since the powers of
$-\xi$ form a basis as well. The development proves the lemma by a self-contained computation instead.
\end{proof}

\begin{lemma}[Discrete spaces]\label{pf:steps1:pbfdisc}
Let $E$ be the constant space on a presheaf $F$, and assume that $E$ is $q$-invertible. Then
\eqref{pf:steps1:eq:pbf} is bijective for every vector bundle $V$ on $E$ and all $n,i$.
\formallabel{GlobalTheory.EtC.pbf_disc}
\formallabel{GlobalTheory.T.projectiveBundleFormula_discrete_space}
\end{lemma}

\begin{proof}
By \cref{pf:steps1:disc} $V$ is trivial. There are a bounded below complex $J$ of injective sheaves of
$\Z/q$-modules and cochain maps $\psi_j\colon I(\Z/q(i-j))[-2j]\to J$ such that, for every such $F$, $J$
computes the cohomology of $\mathbb P(V)$ and \eqref{pf:steps1:eq:pbf} is, up to a bijection, the map induced by
$\sum_j\psi_j$ (\lref{GlobalTheory.T.exists_injective_cochainModel_trivialProjectiveBundle}, deferred). By
\cref{pf:steps1:pbfsch} the map induced by $\sum_j\psi_j$ is bijective on the cohomology of $\stepsIsch(U)$ for
every scheme $U$ on which $q$ is invertible. A cochain map between bounded below complexes of injective
sheaves of $\Z/q$-modules which induces bijections on the cohomology of all such representable presheaves
induces bijections on the cohomology of every $q$-invertible constant space
(\lref{GlobalTheory.T.injectiveComplexMap_quasiIso_of_unit_schemes}, deferred). This proves the lemma.
\end{proof}

\begin{lemma}[Reduction to discrete spaces]\label{pf:steps1:pbfred}
If \eqref{pf:steps1:eq:pbf} is bijective for all vector bundles on all $q$-invertible constant spaces, then it is
bijective for all vector bundles on all $q$-invertible spaces.
\formallabel{GlobalTheory.PBFg.pbf_of_disc}
\formallabel{GlobalTheory.T.projectiveBundleFormula_of_discrete_spaces}
\formallabel{GlobalTheory.PBFc.pbf_of_model}
\end{lemma}

\begin{proof}
Let $E$ be $q$-invertible and $V$ a vector bundle on $E$. Up to an isomorphism inducing bijections on
cohomology, $\mathbb P(V)$ is the diagonal of a bisimplicial space whose $p$-th column is the trivial projective
bundle over the constant space $E_p$, and by the Eilenberg--Zilber (Alexander--Whitney) comparison its
cohomology is that of the total complex of the associated double complex. Choose cocycles $z_j$ representing the
$j$-th powers of the first Chern class of the tautological line bundle on a universal space; their pull-backs
represent $c_1(\OO_V)^j$ on $\mathbb P(V)$ and, on each column, the corresponding powers for the trivial projective
bundle. Cup product with the $z_j$ defines a morphism of column-filtered double complexes inducing
\eqref{pf:steps1:eq:pbf} on total cohomology and, on the $p$-th column, the map \eqref{pf:steps1:eq:pbf} of the
constant space $E_p$, which is bijective by hypothesis. A morphism of first-quadrant column-filtered double
complexes which is bijective on the cohomology of every column is bijective on total cohomology. The
construction of the bisimplicial model and of the $z_j$, and the filtration argument, are carried out in the
deferred lemmas \lref{GlobalTheory.T.twisted_cocycle_classes_pullback_cup},
\lref{GlobalTheory.T.twisted_cocycle_simplicial_levels_cup} and
\lref{GlobalTheory.ColFilt.bijective_total_of_terms}.
\end{proof}

\begin{theorem}[Projective bundle formula; D1]\label{pf:steps1:D1}
Let $q\ge1$, $E$ a $q$-invertible space, $N\ge0$ and $V\colon E\to\stepsIBGL_N$. For all $n,i\in\Z$ the map
\eqref{pf:steps1:eq:pbf} is bijective.
\formallabel{GlobalTheory.EtC.projectiveBundleFormula_vectorBundle_any_space}
\formallabel{GlobalTheory.EtC.projectiveBundleFormula_of_unitOnEtalePoints}
\end{theorem}

\begin{proof}
Combine \cref{pf:steps1:pbfdisc,pf:steps1:pbfred}.
\end{proof}

For schemes this is the classical projective bundle formula of \'etale cohomology
\cite{Grothendieck1958}, \cite[Exp.~VII]{Grothendieck1977}; the formal statement allows arbitrary spaces (for
instance \v{C}ech nerves of covers, on which a map to $\stepsIBGL_N$ is a descent datum for a vector bundle).

\begin{theorem}[Power-sum Chern classes and slant products; D2, D3]\label{pf:steps1:D2}
\textup{(D2)} For $q\ge1$ and $k\ge1$ there is a rule assigning to every vector bundle $V\colon E\to\stepsIBGL_N$ a class
$s_k(V)\in\stepsIH^{2k}(E,\Z/q(k))$ which is natural under pull-back along maps into $q$-invertible spaces, additive
under block sums on $q$-invertible spaces, and satisfies $s_k(L)=c_1(L)^k$ for line bundles $L$ on $q$-invertible
spaces.

\textup{(D3)} For every ring $R$, group $\Gamma$, scheme $X$ and $m\ge0$, every $m$-cycle $c$ of the group homology of
$\Gamma$ with coefficients in $R$ is represented by a ``slant'' cochain map lowering degrees by $m$ from the chains of
$\stepsIsch(X)$ to the chains of $\mathrm B\Gamma\times\stepsIsch(X)$, and the slant maps
$\mathrm{H}_m(\Gamma,R)\otimes\stepsIH^{m+n}(\mathrm B\Gamma\times\stepsIsch(X),F)\to\stepsIH^n(\stepsIsch(X),F)$ used by the
development are given by composition with such cochain maps.

In particular \textup{(D2)} follows from the propositions A1--A5, A7--A11, B1--B5, B7, C1--C5 and D1.
\formallabel{UpperTheory.Et.etale_duality_trace_isomorphism}
\formallabel{UpperTheory.Et.integerApproximationCount_holds}
\formallabel{UpperTheory.Et.D2.a0}
\end{theorem}

\begin{proof}
(D2) is Grothendieck's theory of Chern classes \cite{Grothendieck1958}, in its \'etale form
\cite[Exp.~VII, \S3]{Grothendieck1977}. For a rank-$N$ bundle $V$ on a $q$-invertible space $E$, the classes $c_j(V)\in
\stepsIH^{2j}(E,\Z/q(j))$ are defined by the relation $\sum_{j=0}^N\pi_V^*c_j(V)\cup c_1(\OO_V)^{N-j}=0$, $c_0=1$, in the
cohomology of $\mathbb P(V)$, which determines them by the projective bundle formula \cref{pf:steps1:D1}; they are natural
because \eqref{pf:steps1:eq:pbf} and $c_1$ are natural, and $c_1(L)$ is the first Chern class for $N=1$. The Whitney formula
$c(V\oplus W)=c(V)c(W)$ for block sums is proved by the splitting principle: pull-back to the flag bundle of $V$ (an
iterated projective bundle, to which \cref{pf:steps1:D1} applies at each stage) is injective, and there $V$ acquires a
filtration with line bundle quotients; the argument of \cite[\S\S3--4]{Grothendieck1958} goes through verbatim, using only the
projective bundle formula for these iterated projective bundles, the naturality and multiplicativity of cup products and
of $c_1$, and the Mayer--Vietoris property A4. Then $s_k$ is the $k$-th Newton polynomial in
the $c_j$, so it is additive under block sums, natural, and $s_k(L)=c_1(L)^k$. Since this holds unconditionally, the
implication in the last sentence of the theorem holds. The development proves (D3) from A1--D2
(\lref{UpperTheory.Et.D3.a0}, deferred), by an explicit construction of the slant cochain maps.
\end{proof}

\subsubsection{Finiteness}

\begin{theorem}[Finiteness from the axioms; E4]\label{pf:steps1:E4}
Assume the propositions A1--A5, A7--A11, B1--B5, B7, C1--C5, D1--D3 and E1--E3. Then E4 holds: for $q\ge1$, an
algebraically closed field $k$ with $q\in k^\times$, a proper $k$-scheme $X$ and $n,i\in\Z$, the group
$\stepsIH^n(X,\Z/q(i))$ is finite.
\formallabel{UpperTheory.Et.e4}
\formallabel{UpperTheory.Et.EtaleTopology.top}
\end{theorem}

\begin{proof}
The conclusion holds unconditionally: for $n<0$ the group is zero by A9, and for $n\ge0$, by
\cref{pf:steps1:compare}, it is $\stepsIH^n(X_{\mathrm{et}},\mu_q^{\otimes i})$, which is finite by the finiteness
theorem for proper schemes over a separably closed field \cite[Exp.~XIV, Cor.~1.2]{ArtinGrothendieckVerdier1972},
$\mu_q^{\otimes i}$ being a constructible torsion sheaf. Hence the implication holds. One of the two formal statements
attached to the theorem has, besides A1--E3, the \'etale descent of Picard groups as a further hypothesis; it is implied
by the other.
\end{proof}

\begin{remark}[The formal proof of E4]\label{pf:steps1:E4rem}
The formal proof genuinely derives E4 from the hypotheses, together with the \'etale descent of Picard groups,
\lref{GlobalTheory.T.etale_descent_picard_group}. It is an induction on the degree $n$ and on the Krull dimension of
$X$: zero-dimensional schemes are treated directly; an integral $X$
of positive dimension is related, by proper surjective maps and a fibration in curves over an algebraically closed
extension field, to schemes of smaller dimension and to curves; for curves the cohomology is controlled by Kummer
theory (finiteness of units modulo $q$-th powers and of the $q$-torsion of Picard groups), classes of degree $1$ and
$2$ on curves are killed by finite covers, and the comparison uses the descent spectral sequence of a proper
surjective map and the constancy of the cohomology of geometric fibres. Since the conclusion is classical, these
$133$ intermediate steps of the formal proof are not converted here.
\end{remark}

\begin{corollary}[Finiteness]\label{pf:steps1:E4cor}
E4 holds. \formallabel{UpperTheory.Et.etale_finiteness_of_duality}
\end{corollary}

\begin{proof}
Apply \cref{pf:steps1:E4} to the propositions A1--E3, which are theorems of the development: D1 is
\cref{pf:steps1:D1}, D2 and D3 are \cref{pf:steps1:D2}, E2 is \cref{pf:steps1:E2}, E1 and E3 are the results cited
above, and A1--A5, A7--A11, B1--B5, B7, C1--C5 are, in this order,
\lref{GlobalTheory.T.quasiIso_hypercohomology_bijective},
\lref{GlobalTheory.T.etC_a2}, \lref{GlobalTheory.T.levelwiseEquivalence_isEquivalence},
\lref{GlobalTheory.T.mayerVietoris_equivalence_of_cover},
\lref{GlobalTheory.T.positiveDegreeClasses_die_on_etaleCover},
\lref{GlobalTheory.T.etaleCohomology_coefficient_longExactSequence},
\lref{GlobalTheory.T.etaleCohomology_coefficientRing_invariance},
\lref{GlobalTheory.T.etaleCohomology_negative_vanishing_degree_zero},
\lref{GlobalTheory.T.etaleCohomology_bijective_of_support_bijective},
\lref{GlobalTheory.EtC.RedCoef.a11},
\lref{GlobalTheory.T.etaleCohomology_twisted_filtered_colimit_continuity},
\lref{GlobalTheory.T.strictlyHenselianLocalRing_acyclic_every_sheaf},
\lref{GlobalTheory.T.etC_b3}, \lref{GlobalTheory.T.integralSurjection_cechDescent_tateTwist},
\lref{GlobalTheory.T.absolutelyIntegrallyClosed_domain_tateTwistAcyclic},
\lref{GlobalTheory.T.etaleCohomology_field_eq_galoisCohomology},
\lref{GlobalTheory.T.etale_cohomology_vanishing_algebraic_over_ratFunc},
\lref{GlobalTheory.T.smooth_affine_curve_etale_cohomologicalDimension_le_one},
\lref{GlobalTheory.T.etale_cohomology_ind_etale_base_change},
\lref{GlobalTheory.T.homotopy_invariance_etale_alt} and
\lref{GlobalTheory.T.etale_cohomology_locally_acyclic}, each of which states the corresponding proposition of
\cref{pf:steps1:def:axioms}.
\end{proof}

\begin{lemma}[Finiteness of $\ell$-adic cohomology]\label{pf:steps1:vfin}
Let $\ell$ be a prime, $k$ an algebraically closed field with $\ell\in k^\times$, $X$ a proper $k$-scheme and
$n,i\in\Z$. Then $\stepsIHl^n(X)(i)$ is a finitely generated $\Zl$-module and $\stepsIHQ^n(X)(i)$ is a
finite-dimensional $\Ql$-vector space.
\formallabel{UpperTheory.Gal.G3.vFin_holds} \formallabel{LateralTheory.EWeil.AdicKerPf.AdicKer_of}
\end{lemma}

\begin{proof}
Write $M_a=\stepsIH^n(X,\Z/\ell^a(i))$; these groups are finite by \cref{pf:steps1:E4cor}, and
$M=\stepsIHl^n(X)(i)=\varprojlim_aM_a$. The short exact sequences
$0\to\Z/\ell^a(i)\xrightarrow{\ell}\Z/\ell^{a+1}(i)\to\Z/\ell(i)\to0$ give exact sequences
$M_a\xrightarrow{\ell_*}M_{a+1}\to M_1$ (A7; the Bockstein sequences of the development,
\lref{LateralTheory.Fin.AdicFin.etale_bockstein_sequence_adic_finiteness}, deferred), and the composite of
reduction $M_{a+1}\to M_a$ with $\ell_*$ is multiplication by $\ell$. If $x=(x_a)\in M$ has $x_1=0$, the finite
nonempty sets $Y_a=\{y\in M_a:\ell_*y=x_{a+1}\}$ form an inverse system under reduction, so they have a compatible
family $y=(y_a)$, and then $x_{a+1}=\ell_*(\mathrm{red}\,y_{a+1})=\ell y_{a+1}$, i.e.\ $x=\ell y$. Hence $M/\ell M$
embeds into the finite group $M_1$. Moreover $\bigcap_k\ell^kM=0$, since $\ell^k$ kills $M_a$ for $a\le k$. As $M$ is
a profinite $\Zl$-module with $M/\ell M$ finite, it is finitely generated (topological Nakayama); in the
development this last step is \lref{LateralTheory.EWeil.AdicFinPf.AdicFin_of}, which states that such limits
are finitely generated. Finite dimensionality of $\stepsIHQ^n=\stepsIHl^n\otimes\Ql$ follows.
\end{proof}

\subsubsection{The generic fibre of a strictly henselian local ring of the affine line}

A ring $T$ is \emph{absolutely integrally closed} if every monic polynomial of positive degree over $T$ has a root
in $T$. An absolutely integrally closed local domain is integrally closed, has an algebraically closed fraction
field, and is strictly henselian.

\begin{theorem}[$\stepsIH^1$ of generic fibres]\label{pf:steps1:genfib}
Let $q\ge1$, let $T$ be an absolutely integrally closed local domain with $q\in T^\times$, and let $A$ be a
$T[x]$-algebra (compatibly a $T$-algebra) which is a henselian local ring with separably closed residue field and in
which every element lies in the image of some $T[x]$-algebra map $E\to A$ from an \'etale $T[x]$-algebra $E$. Then
for every $i\in\Z$
\[
 \stepsIH^1\bigl(\stepsIsp(\mathrm{Frac}(T)\otimes_TA),\,\Z/q(i)\bigr)=0 .
\]
\formallabel{GlobalTheory.T.subsingleton_ht_one_genericFibre_of_sh_etale_cover}
\formallabel{GlobalTheory.LocAcycDegOne.tateTwistCohomology_genericFibre_henselLine_subsingleton}
\end{theorem}

The hypothesis on $A$ is satisfied by every strict henselisation of $T[x]$ at a prime ideal, and also by its
quotients. For a strict henselisation the theorem is the degree-one case of the local acyclicity of the smooth
morphism $\mathbb A^1_T\to\Spec T$ \cite[Exp.~XV, Th.~2.1]{ArtinGrothendieckVerdier1972}, read at the generic
geometric point of $\Spec T$. The proof of the development reduces the theorem to Kummer theory and to the
triviality of $\Z/q$-torsors (\cref{pf:steps1:torsor}). Fix $q,T,A$ as in the theorem, put $K=\mathrm{Frac}(T)$,
$S=T\smallsetminus\{0\}$ and $B=K\otimes_TA=S^{-1}A$.

\begin{lemma}[Roots of unity]\label{pf:steps1:zeta}
There is $\zeta\in T$ with $\Phi_q(\zeta)=0$; it satisfies $\zeta^q=1$, and $1-\zeta^k\in T^\times$ for $0<k<q$. In
particular every $T$-algebra contains a root of $\Phi_q$.
\formallabel{GlobalTheory.LocAcycDegOne.exists_cyclotomic_root}
\formallabel{GlobalTheory.exists_pow_eq_one_isUnit_one_sub_pow}
\end{lemma}

\begin{proof}
$\Phi_q$ is monic of degree $\varphi(q)>0$, so it has a root $\zeta\in T$. As $q\in T^\times$, $\zeta$ is a primitive
$q$-th root of unity in the domain $T$, so $X^q-1=\prod_{k=0}^{q-1}(X-\zeta^k)$; dividing by $X-1$ and evaluating at
$1$ gives $\prod_{k=1}^{q-1}(1-\zeta^k)=q\in T^\times$, whence each factor is a unit.
\end{proof}

\begin{lemma}[Presentation by a strict henselisation]\label{pf:steps1:shpres}
There are a henselian local ring $O$ with separably closed residue field which is a $T[x]$-algebra and the directed
union of $T[x]$-subalgebras $O_\alpha$, each isomorphic to the localization of an \'etale $T[x]$-algebra at a prime
ideal, and a surjective $T[x]$-algebra map $\pi\colon O\to A$.
\formallabel{GlobalTheory.T.surj_of_strictHenselisation_polynomial_alt}
\formallabel{GlobalTheory.surj_of_strictHenselisation_polynomial}
\end{lemma}

\begin{proof}
Let $\mathfrak P\subset T[x]$ be the preimage of the maximal ideal of $A$ and $O$ the strict henselisation of $T[x]$
at $\mathfrak P$ with respect to the separable closure of $\kappa(\mathfrak P)$ in the residue field of $A$; it is the
directed union of the local rings of \'etale $T[x]$-algebras at primes above $\mathfrak P$. Since $A$ is henselian
local with separably closed residue field, the universal property of the strict henselisation gives a local
$T[x]$-algebra map $\pi\colon O\to A$. Every $a\in A$ lies in the image of some $\phi\colon E\to A$ with $E$ \'etale over
$T[x]$; the prime $\phi^{-1}(\mathfrak m_A)$ makes $E$ an \'etale neighbourhood of $\mathfrak P$ whose residue field
embeds into that of $A$, so $\phi$ factors through $O$, and $a\in\pi(O)$.
\end{proof}

Fix $O$ and $\pi$ as in \cref{pf:steps1:shpres} and put $L=S^{-1}O$; then $\pi$ induces a surjection $L\to B$. Since
$T$ is a normal domain, each $O_\alpha$ is a normal domain, hence so is $O$.

\begin{lemma}[Primes of the generic fibre]\label{pf:steps1:primes}
Every nonzero prime ideal $\mathfrak q$ of $L$ is maximal, and every element of $L$ is congruent modulo $\mathfrak q$
to an element $c/d$ of $K$ ($c,d\in T$, $d\ne0$). In particular $L/\mathfrak q\cong K$.
\formallabel{GlobalTheory.T.exists_ne_zero_polynomial_mem_prime_localization}
\formallabel{GlobalTheory.T.prime_tensor_polynomial_isMaximal_of_mem_image}
\end{lemma}

\begin{proof}
$L$ is the directed union of the rings $S^{-1}O_\alpha$, which are localizations of \'etale $K[x]$-algebras at
primes; in particular $L$ is a domain. Let $\mathfrak q\ne0$ and $0\neq y\in\mathfrak q$, say $y\in S^{-1}O_\alpha$.
The prime $\mathfrak q_\alpha=\mathfrak q\cap S^{-1}O_\alpha$ is nonzero; since the fibres of an \'etale map are
discrete, the primes $0\subsetneq\mathfrak q_\alpha$ of the domain $S^{-1}O_\alpha$ cannot both lie over
$0\subset K[x]$, so $\mathfrak q\cap K[x]$ is a nonzero prime, i.e.\ $(x-a)$ for some $a\in K$ ($K$ is algebraically
closed). Then $L/\mathfrak q$ is a directed union of domains which are localizations of \'etale algebras over
$K[x]/(x-a)=K$, i.e.\ finite separable field extensions of $K$, hence equal to $K$. So $\mathfrak q$ is maximal with
residue field $K$, and every element is congruent to the image of an element of $K$.
\end{proof}

\begin{lemma}[Units]\label{pf:steps1:units}
Every unit of $B$ is a $q$-th power.
\formallabel{GlobalTheory.LocAcycDegOne.unitsL}
\formallabel{GlobalTheory.T.henselian_directed_union_subalgebras}
\formallabel{GlobalTheory.QUnits.tensorFractionRing_units_pow_surj}
\formallabel{GlobalTheory.LocAcycUnits.units_localization_pow_surj}
\formallabel{GlobalTheory.T.exists_root_of_unit_zero_dimensional}
\end{lemma}

\begin{proof}
First let $\ker(L\to B)=0$, so $B\cong L$. Every unit $u$ of $L$ has the form $u=(c/d)\,w$ with $c,d\in S$ and
$w\in O^\times$, by \lref{GlobalTheory.T.units_localization_strictHenselian_eq_const_mul_unit} (deferred), which
states exactly this: units of the generic fibre of $O$ are constants times units of $O$. The elements $c,d$ have
$q$-th roots in $T$ (roots of $X^q-c$ and $X^q-d$), and $w$ has a $q$-th root in $O$ because $O$ is henselian with
separably closed residue field and $X^q-w$ is separable modulo $\mathfrak m_O$ ($q$ and $w$ are units). Hence $u$ is a
$q$-th power.

Now let $J=\ker(L\to B)\ne0$. Every prime of $L$ containing $J$ is nonzero, hence maximal by
\cref{pf:steps1:primes}, so $B\cong L/J$ is zero-dimensional. Let $u\in B^\times$, lift it to $w_0\in L$, and let
$\mathfrak m$ be a maximal ideal of $L/J$, with preimage $\mathfrak q$ in $L$. By \cref{pf:steps1:primes},
$dw_0\equiv c\pmod{\mathfrak q}$ for some $c,d\in T$, $d\ne0$. Let $r\in T$ with $r^q=cd^{q-1}$; then
$(r/d)^q-w_0=c/d-w_0=-d^{-1}(dw_0-c)\in\mathfrak q$. So $u$ is a $q$-th power modulo every maximal ideal of the
zero-dimensional ring $B$, in which $q$ is invertible; hence it is a $q$-th power. Indeed, modulo the nilradical the
ring is von Neumann regular, where roots defined on the members of a finite cover of the spectrum by clopen sets can
be patched with orthogonal idempotents; the root then lifts along the nil ideal by Hensel's lemma, since $X^q-u$ has
invertible derivative at any root.
\end{proof}

\begin{lemma}[Kummer step]\label{pf:steps1:kummer}
Let $R$ be a ring, $q\ge1$ with $q\in R^\times$, and $\zeta\in R$ with $\zeta^q=1$ and $1-\zeta^k\in R^\times$ for
$0<k<q$. Suppose that every $\Z/q$-torsor over $R$ has a section, where a \emph{$\Z/q$-torsor} is a finite projective
faithfully flat $R$-algebra $D$ with an action of $\Z/q$ by $R$-algebra automorphisms such that
$D\otimes_RD\to\prod_{g\in\Z/q}D$, $x\otimes y\mapsto(x\,g(y))_g$, is bijective, and a section is an $R$-algebra map
$D\to R$. Then $\Pic(R)$ has no nontrivial $q$-torsion.
\formallabel{GlobalTheory.T.picard_torsion_trivial_of_torsors_have_sections}
\formallabel{GlobalTheory.KummerStep.kummer_step}
\formallabel{GlobalTheory.T.kummerAlgebra_finite_projective_faithfullyFlat}
\formallabel{GlobalTheory.T.kummerAlgebra_zmod_torsor_bijective}
\end{lemma}

\begin{proof}
Let $\mathcal L$ be an invertible $R$-module with $\mathcal L^{\otimes q}\cong R$, and let $\phi\colon\mathcal
L^{\otimes q}\to R$ be an isomorphism, viewed as a $q$-linear form; for every prime $\mathfrak p$ there is
$e\in\mathcal L$ with $\phi(e,\dots,e)\notin\mathfrak p$. Let $D=\bigoplus_{k=0}^{q-1}\mathcal L^{\otimes k}$ with the
multiplication given by the tensor product and $\phi$ (the Kummer algebra), and let the generator of $\Z/q$ act on
$\mathcal L^{\otimes k}$ by $\zeta^k$. Locally on $\Spec R$, where $\mathcal L=Re$ is free, $D\cong R[y]/(y^q-u)$ with
$u=\phi(e,\dots,e)\in R^\times$ and the generator acting by $y\mapsto\zeta y$; this is finite free of rank $q$ and
faithfully flat, and the torsor map is bijective because $y^q-u$ splits over $D$ as $\prod_k(Y-\zeta^ky)$ with
$\zeta^ky-\zeta^{k'}y$ a unit for $k\ne k'$ (here $q\in R^\times$ and $1-\zeta^k\in R^\times$ are used). Hence $D$ is a
$\Z/q$-torsor. A section $s\colon D\to R$ restricts to an $R$-linear map $\mathcal L\to R$ which is an isomorphism,
since locally $s(y)^q=u$ is a unit. So $[\mathcal L]=1$ in $\Pic(R)$.
\end{proof}

\begin{lemma}[Torsors over the generic fibre]\label{pf:steps1:torsor}
Let $T$ be an absolutely integrally closed domain, $\mathfrak P$ a prime of $T[x]$, $O$ a strict henselisation of
$T[x]$ at $\mathfrak P$, $q\ge1$ with $q\in T^\times$, and $L=(T\smallsetminus0)^{-1}O$. Then every $\Z/q$-torsor over
$L$ (in the sense of \cref{pf:steps1:kummer}) has a section.
\formallabel{GlobalTheory.LocAcyc.torsor_trivial}
\end{lemma}

\begin{proof}
A $\Z/q$-torsor over $L$ in this sense is a finite \'etale Galois covering of $\Spec L$ with group $\Z/q$. Let
$\mathfrak p=\mathfrak P\cap T$; the localization $T_{\mathfrak p}$ is again absolutely integrally closed, hence
strictly henselian, and $O$ is a $T_{\mathfrak p}$-algebra, so $\Spec L=\Spec(O\otimes_{T_{\mathfrak p}}K)$ is the
geometric fibre, over the generic point of $\Spec T_{\mathfrak p}$ (with algebraically closed residue field $K$), of
the strict localization of $\mathbb A^1_T$ at a geometric point over $\mathfrak P$. The smooth morphism
$\mathbb A^1_T\to\Spec T$ is locally $1$-aspherical for the set of primes invertible on $\mathbb A^1_T$
\cite[Exp.~XV, Th.~2.1]{ArtinGrothendieckVerdier1972}; as $q\in T^\times$, all prime factors of $q$ belong to this set,
so $\stepsIH^1(\Spec L,\Z/q)=0$: every such covering is trivial, and in particular has a section.
\end{proof}

\begin{remark}[The formal proof of \cref{pf:steps1:torsor}]\label{pf:steps1:torsorrem}
The development proves \cref{pf:steps1:torsor} without quoting \cite{ArtinGrothendieckVerdier1972}, along the classical
route of \cite[Exp.~XV, \S\S2--3]{ArtinGrothendieckVerdier1972}. The case where $\mathfrak P$ is the origin
$(\mathfrak m_T,x)$ is \cref{pf:steps1:origin} below; general primes are reduced to closed points of the closed fibre and
to the origin. We convert the proof of the origin case, which carries the argument; the reduction of the general case to
it is not needed here, \cref{pf:steps1:torsor} being classical.
\end{remark}

\begin{lemma}[Points over strictly henselian rings]\label{pf:steps1:ptsh}
Let $A$ be a domain, $R$ a local $A$-algebra, $j\colon R\to O$ a ring map to a strictly henselian local ring, $\beta\colon
O\to B$ and $\psi\colon L_R=(A\smallsetminus0)^{-1}R\to B$ ring maps with $\psi|_R=\beta j$, and $D$ a finite $L_R$-algebra
such that the integral closure $C$ of $R$ in $D$ is a nonzero finite \'etale $R$-algebra. Then there is a ring map
$\sigma\colon D\to B$ with $\sigma|_{L_R}=\psi$.
\formallabel{GlobalTheory.T.exists_point_closureFiniteEtale_strictHenselian}
\formallabel{GlobalTheory.KeyOriginNoeth.exists_point_of_closureFiniteEtale}
\end{lemma}

\begin{proof}
The $O$-algebra $C\otimes_RO$ is nonzero (it is faithfully flat over $O$, $C$ being finite \'etale and nonzero over the
local ring $R$) and finite \'etale over the strictly henselian ring $O$, hence a finite product of copies of $O$: by
\cite[Tag~04GL]{Stacks2024} it is a product of quotients of $O$, which are flat over $O$ and hence equal to $O$, and of an
algebra with no prime over the maximal ideal, which vanishes because it is finite over $O$. A projection gives an
$R$-algebra map $\gamma\colon C\to O$. Every $d\in D$ is integral over
$L_R$, so $ad\in C$ for some $a\in A\smallsetminus0$; hence $D=(A\smallsetminus0)^{-1}C$. The images of the elements of
$A\smallsetminus0$ under $\beta\gamma$ are their images under $\psi$, which are units of $B$; so $\beta\gamma$ extends to
$\sigma\colon D\to B$, and $\sigma|_{L_R}=\psi$ because both restrict to $\beta j$ on $R$.
\end{proof}

\begin{lemma}[The origin case]\label{pf:steps1:origin}
Let $T$ be an absolutely integrally closed local domain, $q\ge1$ with $q\in T^\times$, $\mathfrak P_0=\mathfrak m_T[x]+xT[x]$ the
origin of $T[x]$, $O$ a strict henselisation of $T[x]$ at $\mathfrak P_0$, and $L=(T\smallsetminus0)^{-1}O$. Then every
$\Z/q$-torsor over $L$ has a section.
\formallabel{GlobalTheory.KeyOriginCut.keyOrigin}
\formallabel{GlobalTheory.T.torsor_point_over_strictHenselisation_noetherianNormal}
\formallabel{GlobalTheory.KeyOriginNoeth.torsor_point_of_essFiniteType}
\end{lemma}

\begin{proof}
For a ring $R$ over $T[x]$ write $L(R)=(T\smallsetminus0)^{-1}R$, and similarly $L_A(R)=(A\smallsetminus0)^{-1}R$ for a subring
$A\subseteq T$. Let $D$ be a $\Z/q$-torsor over $L$.

\emph{Step 1: an \'etale neighbourhood.} The ring $O$ is the filtered colimit of the local rings $E_Q$ of \'etale
$T[x]$-algebras $E$ at primes $Q$ over $\mathfrak P_0$. Since $D$, its $\Z/q$-action and the inverse of the torsor
isomorphism are given by finitely many elements and relations, there are such $(E,Q)$ with a map $j\colon E_Q\to O$ over
$T[x]$ and a $\Z/q$-torsor $D_1$ over $L(E_Q)$ such that every ring map $D_1\to L$ over $L(j)$ yields a section of $D$
(\lref{GlobalTheory.T.strict_henselisation_prime_eq}, deferred: a routine limit argument).

\emph{Step 2: a noetherian model.} The ring $T$ is the directed union of its finitely generated subrings. Descending
$(E,Q,j,D_1)$ to one of them and replacing it by the localization at the contraction of $\mathfrak m_T$ of its integral
closure (which lies in $T$, as $T$ is integrally closed, and is finite over the subring, rings of finite type over $\Z$
being Nagata), one obtains a noetherian integrally closed local domain $A$, essentially of finite type over $\Z$, with a
local injection $A\to T$ and $q\in A^\times$, an \'etale $A[x]$-algebra $E_A$ with a prime $Q_A$ over the origin of $A[x]$,
a map $j_A\colon(E_A)_{Q_A}\to O$ compatible with $A[x]\to T[x]\to O$, a map $\psi_A\colon L_A((E_A)_{Q_A})\to L$ extending
$j_A$, and a torsor $D_A$ over $L_A((E_A)_{Q_A})$ such that every ring map $D_A\to L$ over $\psi_A$ yields a ring map
$D_1\to L$ over $L(j)$ (\lref{GlobalTheory.T.exists_noetherian_model_torsor_origin}, deferred).

\emph{Step 3: removing ramification.} After enlarging $A$ (to a noetherian integrally closed local domain $A'\subseteq T$
essentially of finite type over $\Z$) and passing to a smaller \'etale neighbourhood, the data become \emph{good}: the
neighbourhood $R'=(E_{A'})_{Q_{A'}}$ has a retraction $\rho'\colon R'\to A'$ over $A'$ with kernel generated by $x$, and the
integral closure $C'$ of $R'$ in the new torsor $D_{A'}$ is \'etale over $R'$ at every prime whose contraction to $A'$ has
height at most one, and admits a ring map $C'\to\mathrm{Frac}(A')$ compatible with $\rho'$; moreover there are maps
$\beta\colon L_A((E_A)_{Q_A})\to L_{A'}(R')$ and $h\colon D_A\to D_{A'}$ over $\beta$ with $\psi_{A'}\beta=\psi_A$
(\lref{GlobalTheory.T.strictHenselisation_enlarge_good_at_origin}, deferred). This is where Abhyankar's lemma enters:
the ramification of the torsor along the height-one primes of the base is tame, since $q$ is invertible, and is killed
by adjoining roots of local parameters \cite[Exp.~X, 3.6]{Grothendieck1971}, which exist in $T$ because $T$ is absolutely
integrally closed.

\emph{Step 4: purity.} For good data the integral closure $C'$ is a nonzero finite \'etale $R'$-algebra
(\lref{GlobalTheory.T.integralClosure_finiteEtale_originKeyLemma_essFiniteType}, deferred): \'etaleness in codimension
at most one over the base and along the zero section propagates to all of $\Spec R'$, by purity of the branch locus as in
\cite[Exp.~XV, \S2]{ArtinGrothendieckVerdier1972}.

\emph{Step 5: conclusion.} The ring $O$ is strictly henselian, so \cref{pf:steps1:ptsh}, applied to $A'$, $R'$, the map
$j_{A'}\colon R'\to O$, the localization $O\to L$, $\psi_{A'}$ and $D_{A'}$, gives $\sigma'\colon D_{A'}\to L$ over
$\psi_{A'}$. Then $\sigma'h\colon D_A\to L$ lies over $\psi_{A'}\beta=\psi_A$; Step 2 gives a ring map $D_1\to L$ over
$L(j)$, and Step 1 a section of $D$.
\end{proof}

\begin{lemma}[Picard groups]\label{pf:steps1:pic}
$\Pic(B)$ has no nontrivial $q$-torsion.
\formallabel{GlobalTheory.LocAcycDegOne.picL}
\formallabel{GlobalTheory.T.pic_generic_fibre_quotient_strictHenselisation_no_torsion_alt}
\formallabel{GlobalTheory.pic_generic_fibre_quotient_strictHenselisation_no_torsion}
\formallabel{GlobalTheory.pic_localization_strictHenselisation_no_torsion}
\formallabel{GlobalTheory.StrictHensUnion.isStrictHenselisation_of_directed_union}
\end{lemma}

\begin{proof}
If $\ker(L\to B)=0$, then $B\cong L$. The ring $O$ is a strict henselisation of $T[x]$ at some prime: it is henselian
with separably closed residue field, and the directed union of localizations of \'etale $T[x]$-algebras. Moreover
$q\in L^\times$, the image of $\zeta$ from \cref{pf:steps1:zeta} satisfies the hypothesis of \cref{pf:steps1:kummer},
and every $\Z/q$-torsor over $L$ has a section by \cref{pf:steps1:torsor}. So \cref{pf:steps1:kummer} applies.

If $\ker(L\to B)\ne0$, then all prime ideals of $B$ are maximal: two primes $\mathfrak q\subseteq\mathfrak q'$ of $L$
containing the kernel are nonzero, hence maximal by \cref{pf:steps1:primes}, hence equal. A ring all of whose prime
ideals are maximal has trivial Picard group: an invertible module is locally free of rank one, and over a
zero-dimensional ring a module which is free of rank one at every point is free, by patching local generators with
idempotents on the Boolean spectrum of the reduced ring and lifting along the nilradical.
\end{proof}

\begin{proof}[Proof of \cref{pf:steps1:genfib}]
The element $q$ is invertible in $B$, and $B$ contains a root of $\Phi_q$ (\cref{pf:steps1:zeta}). By
\lref{GlobalTheory.T.tateTwistCohomology_spec_twist_independent} (for a ring $A$ with $q\in A^\times$ containing a
root of $\Phi_q$, the groups $\stepsIH^n(\stepsIsp A,\Z/q(i))$ and $\stepsIH^n(\stepsIsp A,\Z/q(j))$ are isomorphic for
all $n,i,j$) it suffices to treat $i=1$. By \lref{GlobalTheory.T.kummer_etale_vanishing} (if $q\in B^\times$, every
unit of $B$ is a $q$-th power and $\Pic(B)$ has no nontrivial $q$-torsion, then $\stepsIH^1(\stepsIsp B,\Z/q(1))=0$; this
is the Kummer sequence), the theorem follows from \cref{pf:steps1:units,pf:steps1:pic}.
\end{proof}

% ===========================================================================
\subsection{The crystalline comparison theorem and $\sfname{CrysStep}$}\label{sec:pf:steps1:cris}

\subsubsection{The comparison theorem in the form used}

\begin{definition}[Admissibility]\label{pf:steps1:def:adm}
Let $\stepsIBdR$ be Fontaine's field of de Rham periods as constructed in the development and $t\in\stepsIBdR$ the
period of $\Z_\ell(1)$; let $B_{\mathrm{cr}}\subseteq\stepsIBdR$ be the subring generated by the image of
$\stepsIAcris$, $\ell^{-1}$ and $t^{-1}$ (this is Fontaine's ring $\stepsIBcris=\stepsIAcris[1/t]$
\cite{Fontaine1994}). For a subring $S\subseteq\stepsIBdR$, a map $r\colon G_K\to M_d(\Ql)$ is
\emph{$S$-admissible} if there are $P,Q\in M_d(S)$ with $PQ=1$ and $r(\sigma)\,\sigma(P)=P$ for all $\sigma\in G_K$,
i.e.\ the columns of $P$ are the coordinates of a $G_K$-invariant $S$-basis of $\Ql^d\otimes S$. For a period package
$P$, \emph{$P$-admissibility} is the analogous notion with $S$ the image of $P.B$ in $P.B_{\mathrm{dR}}$.
\formal{CrystallineSemistableComparison.Bcr} \formal{DeRhamPotentiallySemistable.DRPst.IsAdm}
\formal{UpperTheory.Gal.Seam.IsAdmP}
\end{definition}

\begin{theorem}[Comparison package for smooth good models]\label{pf:steps1:cst}
Let $K\subset\Qlbar$ be finite over $\Ql$, and let $\stepsIBdR^+\subset\stepsIBdR$ be the valuation ring.
\begin{enumerate}[label=\textup{(\roman*)},nosep]
\item \textup{(Crystalline comparison)} For every smooth good model $M$ over $\OO_K$, every $n$ and every basis $b$ of
  $V_n(M)$, the matrix representation $\mathrm{mat}_n(b)$ is $B_{\mathrm{cr}}$-admissible.
\item \textup{(de Rham comparison)} There are biadditive pairings $c_{M,n}\colon V_n(M)\times\stepsIH^n_{\mathrm{dR}}(M_K/K)^\vee
  \to\stepsIBdR$, for all models $M$ over $\OO_K$ and all $n$, $\Ql$-linear in the first and $K$-linear in the second
  variable, natural for morphisms between smooth good models, such that for every smooth good $M$: $V_n(M)$ and
  $\stepsIH^n_{\mathrm{dR}}(M_K/K)$ are finite-dimensional and $\mathrm{Fil}^{n+1}\stepsIH^n_{\mathrm{dR}}=0$;
  $c(\sigma v,\phi)=\sigma(c(v,\phi))$ for $\sigma\in G_K$; for a basis $b_1,\dots,b_a$ of $V_n(M)$ and $z\in\stepsIBdR^a$,
  $\sum_iz_ic(b_i,\phi)=0$ for all $\phi$ implies $z=0$; $t^rc(v,\phi)\in\stepsIBdR^+$ whenever $\phi$ kills
  $\mathrm{Fil}^{r+1}$; and if $t^r\sum_iz_ic(b_i,\phi)\in\stepsIBdR^+$ for all $r$ and all $\phi$ killing
  $\mathrm{Fil}^{r+1}$, then $z\in(\stepsIBdR^+)^a$.
\item \textup{(Frobenius compatibility)} Let $r$ be a \emph{residue datum} of $K$: a perfect field $k$ of characteristic
  $\ell$, a surjection $\rho\colon\OO_K\to k$ and a ring map $s$ from $k$ to the tilt of $\OO_{\stepsICl}$ with
  $|s(\rho(x))^\sharp-x|<1$ for $x\in\OO_K$; put $K_0=W(k)[1/\ell]$, and let $D_n(M)$ be the $K_0$-vector space with
  Frobenius $\varphi_D$ attached to $M$ and $r$ by the development ($K_0\otimes_{W(k)}$ the inverse limit of a system of
  degree-$n$ cohomology groups over the truncated Witt vectors of $k$). There are biadditive pairings
  $c_{M,n}\colon V_n(M)\times D_n(M)^\vee\to\stepsIBcris$, $\Ql$-linear in the first and $K_0$-linear in the second variable,
  natural for morphisms between smooth good models, such that for every smooth good $M$: $V_n(M)$ and $D_n(M)$ are
  finite-dimensional, $c$ is $G_K$-equivariant, the independence statement of \textup{(ii)} holds with $\stepsIBcris$ in
  place of $\stepsIBdR$, and $c(v,f)=\sum_s\varphi\bigl(c(v,e_s^*)\bigr)\,f(\varphi_De_s)$ for every $K_0$-basis $(e_s)$ of
  $D_n(M)$ with coordinate functionals $(e_s^*)$.
\end{enumerate}
\formallabel{UpperTheory.Ru.crystallineComparison.smoothModel_crystalline}
\formallabel{UpperTheory.Gal.Seam.cstSmooth}
\formallabel{CrystallineSemistableComparison.Close.cstSmooth_of_side}
\formallabel{CrystallineSemistableComparison.Out.cstSmooth_of_side}
\end{theorem}

\begin{proof}
(i) By the crystalline comparison theorem $C_{\mathrm{cris}}$ \cite{FontaineMessing1987,Faltings1989}, \cite{Tsuji1999}
(see also \cite{BhattMorrowScholze2018}), for a proper smooth $\OO_K$-scheme the $G_K$-representation
$\stepsIHQ^n((M_K)_{\Qlbar})(0)$ is crystalline: $\stepsIDcris(V_n(M))=(V_n(M)\otimes\stepsIBcris)^{G_K}$ has $K_0$-dimension
$\dim V_n(M)$, and a $K_0$-basis of it is a $G_K$-invariant $\stepsIBcris$-basis of $V_n(M)\otimes\stepsIBcris$
\cite[Exp.~III]{Fontaine1994rss}. Its coordinate matrix in the basis $b$ is the required $P$.

Classically, (ii) expresses the de Rham comparison theorem \cite{Faltings1988,Scholze2013phodge} with its strict
compatibility with filtrations, and (iii) the compatibility of $C_{\mathrm{cris}}$ with Frobenius, $D_n(M)$ playing the role
of the rational crystalline cohomology of the special fibre. In the development, (i)--(iii) are obtained together: the
three statements are derived by \lref{CrystallineSemistableComparison.cstSmooth_of_side} from the \'etale-side
properties of smooth good models, and these properties hold by \lref{UpperTheory.Gal.Seam.HT.etSideS} (see also
\lref{UpperTheory.Gal.KunnethLefschetz.Top.etSide_sm_holds}). The formal statements attached to the theorem are (i),
the conjunction of (i)--(iii), and (twice) the implication from the \'etale-side properties to (i)--(iii). Only (i) is used
in this section.
\end{proof}

\subsubsection{The comparison pairings of the formal proof}\label{sec:pf:steps1:pseam}

The formal proof of \cref{pf:steps1:cst} constructs comparison maps between $\ell$-adic cohomology and cohomology
theories with coefficients in $\stepsIAcris$ attached to models carrying a line bundle; one of its packages is the
following existence statement.

\begin{definition}[Line data and comparison data]\label{pf:steps1:def:linedata}
A \emph{line datum} on a scheme $X$ is a smooth surjective $f\colon Y\to X$ with a rule $u$ assigning to any two
$V$-points $a,b$ of $Y$ with $fa=fb$ a unit of $\Gamma(V,\OO_V)$, naturally in $V$ (transition functions of a line bundle
along $f$); it is \emph{tautological} if it is the line datum of $\OO(1)$ for some closed immersion of $X$ into a
projective space over the base, along the standard affine cover. For a model $M$ and a predicate $\mathrm{Amp}$ on line
data, the development attaches to $M$ and a line datum $D$ with $\mathrm{Amp}(M,D)$ graded $\stepsIAcris$-modules with
$G_K$-action $\mathcal A(M)$, $\mathcal H(M,D)$, $\mathcal H'(M,D)$ (the $\stepsIAcris$-valued cohomologies of the formal
proof), a Frobenius on $\mathcal A$ and $\mathcal H'$, pull-backs, cup products, and a predicate ``is a first Chern class
of $D$'' in degree $2$; we call this the \emph{comparison data} of $\mathrm{Amp}$.
\formal{CrystallineSemistableComparison.E.LineData}
\formal{CrystallineSemistableComparison.Pb.IsTautologicalLine}
\formal{CrystallineSemistableComparison.dat} \formal{CrystallineSemistableComparison.Model.HA}
\end{definition}

\begin{definition}[de Rham and crystalline pairings]\label{pf:steps1:def:pmap}
A \emph{de Rham pairing} for the comparison data of $\mathrm{Amp}$ on smooth good models consists of additive pairings
$\kappa_j\colon\mathcal H^j(M,D)\times\stepsIH^j_{\mathrm{dR}}(M_K/K)^\vee\to\stepsIBdR$ and $K$-bilinear cup products on de
Rham cohomology, such that for smooth good $M$: $\stepsIH^j_{\mathrm{dR}}(M_K/K)$ is finite-dimensional with
$\mathrm{Fil}^{j+1}=0$, and $\mathrm{Fil}^{d+1}=0$ if $M$ is connected of dimension $d$; cup products respect the Hodge
filtration; $\kappa$ is $\stepsIAcris$-semilinear in the first and $K$-linear in the second variable, $G_K$-equivariant,
natural for morphisms of models, compatible with the filtrations (if $x$ lies in the $n$-th filtration step and $\phi$
kills $\mathrm{Fil}^{r+1}$, then $t^r\kappa(x,\phi)\in t^n\mathcal P$, where $\mathcal P\subset\stepsIBdR$ is the subring of
the period data) and multiplicative ($\kappa(x\cup y,\phi)=\sum_{s,s'}\kappa(x,b_s^\vee)\kappa(y,b_{s'}'^\vee)
\phi(b_s\cup b'_{s'})$ for bases $b,b'$); and if $M$ is connected of dimension $d$, there are a first Chern class $x$
of $D$ in $\mathcal H^2$, its $d$-th cup power $z$ and $\phi$ with $\kappa(z,\phi)$ a unit of $\mathcal P$. For a
residue datum $r$ of $K$ (\cref{pf:steps1:cst}(iii)), a \emph{crystalline pairing} consists of an integer $h$ and pairings
$\psi_j\colon\mathcal H'^j(M,D)\times D_j(M)^\vee\to C$ with values in the crystalline period ring $C$ of $r$, where
$D_j(M)$ is the $K_0$-space of \cref{pf:steps1:cst}(iii), with finiteness, semilinearity
(through $\varphi^h$ on $\stepsIAcris$), $G_K$-equivariance, compatibility with Frobenius and with pull-backs, and a
factorisation property expressing $\varphi^h$-iterates through $\psi$.
\formal{CrystallineSemistableComparison.PMap} \formal{CrystallineSemistableComparison.Res}
\formal{CrystallineSemistableComparison.FMap}
\end{definition}

\begin{theorem}[Existence of comparison pairings]\label{pf:steps1:pseam}
For every prime $\ell$ and finite extension $K\subset\Qlbar$ of $\Ql$ there is a predicate $\mathrm{Amp}$ on pairs (model
$M$ over $\OO_K$, line datum $D$ on $M.X$) such that: every $D$ with $\mathrm{Amp}(M,D)$ has \'etale structure map; every
smooth good model admits some $D$ with $\mathrm{Amp}(M,D)$; there is a de Rham pairing for the comparison data of
$\mathrm{Amp}$ on smooth good models; and for every residue datum $r$ of $K$ there is a crystalline pairing for these
data (with the pull-backs on $\mathcal H'$ of the development).
\formallabel{CrystallineSemistableComparison.Pb.pSeam}
\formallabel{CrystallineSemistableComparison.Pb.pmap}
\formallabel{CrystallineSemistableComparison.Pb.fmap}
\formallabel{CrystallineSemistableComparison.Pb.kap_pow}
\end{theorem}

The decisive property is the last condition on de Rham pairings: the top power of the first Chern class of an ample line
bundle pairs to a unit. It comes from the following classical fact.

\begin{proposition}[Top powers of hyperplane classes]\label{pf:steps1:powK}
Let $K$ be a field of characteristic $0$, $X\to\Spec K$ connected and smooth of relative dimension $d$, $D$ a tautological
line datum on $X$ and $c\in\stepsIH^2_{\mathrm{dR}}(X/K)$ a de Rham first Chern class of $D$. Then $c^d\ne0$ in
$\stepsIH^{2d}_{\mathrm{dR}}(X/K)$.
\formallabel{CrystallineSemistableComparison.Pb.powK_ne_zero}
\formallabel{CrystallineSemistableComparison.Pb.powK_ne_zero_main}
\formallabel{CrystallineSemistableComparison.Pb.powK_ne_zero_of}
\end{proposition}

Here $\stepsIH^*_{\mathrm{dR}}(X/K)$ is the de Rham cohomology of the development (hypercohomology of the de Rham complex in its
\v{C}ech formalism, with its Hodge filtration), and a de Rham first Chern class of $D$ is a class represented, on covers
trivializing $D$, by the logarithmic derivatives of the transition functions. The proof is algebraic and reduces the statement
to a residue computation on a torus. For units $t_0,\dots,t_n$ of a ring $R$ over $K$ and derivations $\partial_0,\dots,\partial_n$
of $R$ over $K$ with $\partial_kt_j=\delta_{kj}$, and for nonempty $S\subseteq\{0,\dots,n+1\}$, let $\mathcal T(S)\subseteq R$ be
the $K$-subalgebra generated by the monomials $t^m=\prod_kt_k^{m_k}$ with $m_k\ge0$ whenever $k+1\notin S$ and $\sum_km_k\le0$ if
$0\notin S$. (With $t_k=y_{k+1}/y_0$ for homogeneous coordinates $y_0,\dots,y_{n+1}$, these are the monomials regular where
$y_i\ne0$ for $i\in S$.) For a subalgebra $B$ and an injective $\kappa\colon\{1,\dots,m\}\to\{0,\dots,n\}$ let
$\mathrm{Jac}_\kappa(B)$ be the additive subgroup generated by the elements $b\det(\partial_{\kappa(k)}b_l)_{k,l}$ with $b,b_l\in B$:
the coefficients, with respect to $dt_{\kappa(1)}\wedge\dots\wedge dt_{\kappa(m)}$, of the forms $b\,db_1\wedge\dots\wedge db_m$.

\begin{lemma}[Logarithmic derivatives of ratios]\label{pf:steps1:dlogratios}
In this situation let $g_0=t_0^{-1}$ and $g_l=t_{l-1}t_l^{-1}$ for $1\le l\le n$. Then
$(\prod_lg_l)^{-1}\det(\partial_kg_l)_{k,l}=\pm(\prod_kt_k)^{-1}$.
\formallabel{CrystallineSemistableComparison.Pb.dlog_ratios}
\end{lemma}

\begin{proof}
Dividing the $l$-th column by $g_l$, the left side is $\det(\partial_kg_l/g_l)$. Now $\partial_kg_0/g_0=-\delta_{k0}t_0^{-1}$ and
$\partial_kg_l/g_l=\delta_{k,l-1}t_{l-1}^{-1}-\delta_{kl}t_l^{-1}$ for $l\ge1$, so the matrix is $\mathrm{diag}(t_k^{-1})$ times the
matrix with $-1$ on the diagonal, $1$ just above it and $0$ elsewhere, whose determinant is $(-1)^{n+1}$.
\end{proof}

\begin{lemma}[Residues]\label{pf:steps1:resid}
Let $K$ be a field of characteristic $0$, $R$ a domain over $K$, $t_0,\dots,t_n\in R^\times$ algebraically independent over $K$,
$\partial_k$ derivations with $\partial_kt_j=\delta_{kj}$, and $B(S)\subseteq R$ ($\emptyset\ne S\subseteq\{0,\dots,n+1\}$)
finitely generated $K$-subalgebras with $\mathcal T(S)\subseteq B(S)$, $B(S)$ integral over $\mathcal T(S)$,
$B(S\smallsetminus\{j\})\subseteq B(S)$ for $S$ the full set, and $R\subseteq\mathrm{Frac}(B(\{0,\dots,n+1\}))$. Then there are no
$\omega_i\in\mathrm{Jac}_{\mathrm{id}}(B(\{0,\dots,n+1\}\smallsetminus\{i\}))$ and
$\eta_k\in\mathrm{Jac}_{\kappa_k}(B(\{0,\dots,n+1\}))$, $\kappa_k$ the increasing enumeration of $\{0,\dots,n\}\smallsetminus\{k\}$,
with
\[
 \Bigl(\prod_kt_k\Bigr)^{-1}=\sum_{i=0}^{n+1}\omega_i+\sum_{k=0}^n(-1)^k\partial_k\eta_k .
\]
\formallabel{CrystallineSemistableComparison.Pb.resid_alg}
\end{lemma}

\begin{proof}
Write $\mathcal T=\mathcal T(\{0,\dots,n+1\})=K[t^{\pm1}]$, $F=K(t)\subseteq L=\mathrm{Frac}(R)$. As $R\subseteq\mathrm{Frac}(B(\{0,
\dots,n+1\}))$ and the latter ring is generated by finitely many elements integral over $\mathcal T$, $L/F$ is finite, and separable
since $\mathrm{char}K=0$; let $\Tr\colon L\to F$ be the trace. Let $\mathrm{res}\colon\mathcal T\to K$ be the coefficient of
$(t_0\cdots t_n)^{-1}$; it vanishes on $\partial_k(\mathcal T)$, since $\partial_kt^m=m_kt^{m-e_k}$. We use three facts. (a) If
$\Tr(y)\in\mathcal T$, then $\Tr(\partial_ky)=\partial_k\Tr(y)$: $\partial_k$ extends uniquely to the separable extension $L$ of $F$,
and the trace commutes with it (\lref{CrystallineSemistableComparison.Pb.trace_deriv}, a routine verification). (b) For every
chart $S$ for which $\mathcal T(S)$ is a localization of a polynomial ring $K[u_0,\dots,u_n]$ with $\mathrm{Frac}(\mathcal T(S))=F$,
every finitely generated $B'\supseteq\mathcal T(S)$ integral over $\mathcal T(S)$ with $R\subseteq\mathrm{Frac}(B')$, and every
injective $\kappa$, the traces of the elements of $B'$-Jacobians with respect to the derivations dual to the $u_j$ lie in
$\mathcal T(S)$ (\lref{CrystallineSemistableComparison.Pb.trace_chart}, deferred; this is where integrality is used).
(c) Each $B(\{0,\dots,n+1\}\smallsetminus\{i\})$ is contained in such a $B'$ for $S=\{0,\dots,n+1\}\smallsetminus\{i\}$
(\lref{CrystallineSemistableComparison.Pb.chart_ext}, deferred).

Apply $\Tr$ and then $\mathrm{res}$ to the relation. The left side gives $[L:F]\,\mathrm{res}((\prod t_k)^{-1})=[L:F]\ne0$. For
$k\le n$, $\Tr(\eta_k)\in\mathcal T$ by (b) with $S$ the full set and $u=t$, so $\mathrm{res}(\Tr(\partial_k\eta_k))=
\mathrm{res}(\partial_k\Tr\eta_k)=0$ by (a). For $i\ge1$, take $u=t$ with $t_{i-1}$ as a polynomial variable: by (b), (c),
$\Tr(\omega_i)\in\mathcal T(\{0,\dots,n+1\}\smallsetminus\{i\})$, whose monomials have nonnegative exponent of $t_{i-1}$, so
$\mathrm{res}(\Tr\omega_i)=0$. For $i=0$ use the coordinates at infinity $u_0=t_0^{-1}$, $u_k=t_kt_0^{-1}$ ($k\ge1$), for which
$\mathcal T(\{1,\dots,n+1\})=K[u_0,u_1^{\pm1},\dots,u_n^{\pm1}]$; the Jacobian of the change of coordinates is
$\det(\partial u_j/\partial t_k)=-t_0^{-(n+2)}$, so $\Tr(\omega_0)=-t_0^{-(n+2)}a$ with $a\in K[u_0,u_k^{\pm1}]$ by (b), (c). A monomial
$u_0^\alpha\prod_{k\ge1}u_k^{\beta_k}$ ($\alpha\ge0$) times $t_0^{-(n+2)}$ equals $t_0^{-(n+2)-\alpha-\sum\beta_k}\prod_{k\ge1}t_k^{\beta_k}$, which is
$(t_0\cdots t_n)^{-1}$ only if $\beta_k=-1$ for all $k$ and $\alpha=-1$; so $\mathrm{res}(\Tr\omega_0)=0$. Hence the right side gives
$0$, a contradiction.
\end{proof}

\begin{proof}[Proof of \cref{pf:steps1:powK}]
If $d=0$, $c^0=1\ne0$ because $X\ne\emptyset$. Let $d=n+1\ge1$ and suppose $c^{n+1}=0$; let $\iota\colon X\to\mathbb P^N_K$ be the
closed immersion underlying $D$.

\emph{Step 1: coordinates} (\lref{CrystallineSemistableComparison.Pb.geo_setup}, deferred). There are $e\ge1$ and homogeneous
polynomials $y_0,\dots,y_{n+1}$ of degree $e$ such that the opens $U_k=X\cap\{y_k\ne0\}$ cover $X$ and all intersections
$U_S=\bigcap_{k\in S}U_k$ are affine, and a nonempty open $V\subseteq U_{\{0,\dots,n+1\}}$ with $R=\Gamma(V,\OO)$ a domain, on which
$t_k=y_{k+1}/y_0$ ($0\le k\le n$) are units, algebraically independent over $K$, with derivations $\partial_k$ of $R$ dual to them, such
that the subalgebras $B(S)\subseteq R$, images of the injective restriction maps $\Gamma(U_S,\OO)\to R$, satisfy the hypotheses of
\cref{pf:steps1:resid}. (The $y_k$ are chosen by Noether normalization of the homogeneous coordinate ring, so that
$(y_0:\dots:y_{n+1})$ is finite over $\mathbb P^{n+1}_K$ and $\Gamma(U_S,\OO)$ is integral over $\mathcal T(S)$.)

\emph{Step 2: the relation} (\lref{CrystallineSemistableComparison.Pb.firstChern_local_exact},
\lref{CrystallineSemistableComparison.Pb.cech_top_rel}, \lref{CrystallineSemistableComparison.Pb.cone_transfer}, deferred).
On the cover $(U_k)$, $c$ is represented by the \v{C}ech cocycle of logarithmic derivatives $d\log(y_k/y_{k'})$, which is locally
exact; computing $c^{n+1}$ in the filtered \v{C}ech complex of the cover, and using that the graded pieces of the Hodge filtration
have no higher cohomology on the affine pieces in the relevant degrees, the vanishing of $c^{n+1}$ means that the \v{C}ech
$(n+1)$-cocycle $d\log(y_0/y_1)\wedge d\log(y_1/y_2)\wedge\dots\wedge d\log(y_n/y_{n+1})$ on $U_{\{0,\dots,n+1\}}$ is the sum of a
\v{C}ech coboundary of $(n+1)$-forms on the $U_{\{0,\dots,n+1\}\smallsetminus\{j\}}$ and the differential of an $n$-form on
$U_{\{0,\dots,n+1\}}$. Restricted to $V$, with $g_l$ as in \cref{pf:steps1:dlogratios}, this gives $x_j\in
\mathrm{Jac}_{\mathrm{id}}(B(\{0,\dots,n+1\}\smallsetminus\{j\}))$ and $\eta_k\in\mathrm{Jac}_{\kappa_k}(B(\{0,\dots,n+1\}))$ with
$(\prod_lg_l)^{-1}\det(\partial_kg_l)=\sum_jx_j+\sum_k(-1)^k\partial_k\eta_k$.

\emph{Step 3.} By \cref{pf:steps1:dlogratios} the left side is $\pm(\prod t_k)^{-1}$; replacing $(x,\eta)$ by $(-x,-\eta)$ if
necessary, we get a relation excluded by \cref{pf:steps1:resid}. Hence $c^{n+1}\ne0$.
\end{proof}

\begin{remark}[A classical argument]\label{pf:steps1:powKrem}
For the classical algebraic de Rham cohomology the statement also follows from Hodge theory: $X$ is projective, $X$ and $c$ descend
to a finitely generated subfield of $K$, which embeds into $\C$; by Grothendieck's comparison theorem \cite{Grothendieck1966},
$\stepsIH^*_{\mathrm{dR}}(X/\C)\cong\stepsIH^*(X(\C),\C)$ compatibly with cup products, $c$ corresponds to a nonzero multiple of
the class of the Fubini--Study form $\omega$, and $\int\omega^d>0$ on every connected component of $X(\C)$. Applying it to the
development's de Rham cohomology would require identifying the latter, with its Chern classes, with the classical theory; the
algebraic proof above avoids this.
\end{remark}

\begin{proof}[Proof of \cref{pf:steps1:pseam}]
Let $\mathrm{Amp}(M,D)$ mean that $D$ is tautological for $M.X\to\Spec\OO_K$. Tautological line data are defined along the
standard affine cover, so their structure maps are \'etale (\lref{CrystallineSemistableComparison.Pb.isO1_etale},
stating this), and every projective model, in particular every smooth good model, carries one
(\lref{CrystallineSemistableComparison.Pb.isO1_exists}, stating this).

\emph{De Rham pairing.} The development constructs a family of pairings $\kappa$ with the linearity, equivariance,
naturality, filtration and multiplicativity properties, together with the compatibility $\kappa(x,\phi)=\phi(c)$ whenever
$x$ is a first Chern class of $D$ in $\mathcal H^2$ and $c$ a de Rham first Chern class of $D$
(\lref{CrystallineSemistableComparison.Pb.crystalline_deRham_firstChern_exists}, deferred). Finiteness of de Rham
cohomology, $\mathrm{Fil}^{j+1}\stepsIH^j_{\mathrm{dR}}=0$ and $\mathrm{Fil}^{d+1}=0$ in dimension $d$ are
\lref{CrystallineSemistableComparison.Pb.m_fin} and \lref{CrystallineSemistableComparison.Pb.m_top}
(deferred), and the cup product on de Rham cohomology respects the filtration. For the last property let $M$ be connected
of dimension $d$ and $D$ tautological. There are a first Chern class $x\in\mathcal H^2(M,D)$ of $D$
(\lref{CrystallineSemistableComparison.Pb.firstChern_acris_family_exists}, deferred) and a de Rham first Chern class
$c$ of the induced tautological datum on $M_K$ (\lref{CrystallineSemistableComparison.Pb.isFirstChernDeRham_exists},
deferred). Let $z$ be the $d$-th cup power of $x$. By multiplicativity and the compatibility above,
$\kappa(z,\phi)=\phi(c^d)$ for all $\phi$. By \cref{pf:steps1:powK}, $c^d\ne0$, so there is a linear form $\phi$ with
$\phi(c^d)=1$, and $\kappa(z,\phi)=1$ is a unit.

\emph{Crystalline pairings.} Given $r$, a Teichm\"uller congruence of some level $N$ holds
(\lref{CrystallineSemistableComparison.Pb.teich}, deferred), a Frobenius factorisation of level $N$ exists
(\lref{CrystallineSemistableComparison.Pb.fac_ex}, stating this), and from these the crystalline pairing is
constructed (\lref{CrystallineSemistableComparison.Pb.fmap_of}, deferred).
\end{proof}

\subsubsection{Proof of $\sfname{CrysStep}$}

We use the following results proved elsewhere in the development.
\begin{enumerate}[label=(\alph*),nosep]
\item \emph{Base change preserves niceness} (\lref{UpperTheory.Gal.Seam.niceBc}).
\item \emph{The local tower lemma} (\lref{UpperTheory.Gal.Seam.locTower_holds}): let $k$ be a field of
  characteristic $0$ with algebraic closure $\Omega$, $\tau\colon k\to K$, $\Omega'$ algebraically closed over $K$ and
  $j\colon\Omega\to\Omega'$ extending $\tau$; if $\rho_0$ is a matrix of the $\theta$-piece of a nice tower $H$ for
  $\Aut_k(\Omega)$, there is a matrix $\rho_1$ of the $\theta$-piece of $\mathrm{bc}_\tau H$ for $\Aut_K(\Omega')$ with
  $\rho_1(\sigma')=\rho_0(\sigma)$ whenever $\sigma'\circ j=j\circ\sigma$.
\item \emph{Base change of covers} (\lref{UpperTheory.Gal.Seam.bcCover}): a cover over $O$ base-changes along
  compatible maps $O\to O'$, $k\to K$ to a cover over $O'$ of the base-changed tower, preserving ``onto'', ``good'' and the
  relations ``lies under''.
\item \emph{The integral embedding} (\lref{UpperTheory.Gal.Seam.intMap}): an inducing bridge gives a ring map
  $\stepsIOw\to\OO_{K_w}$ compatible with $\tau_w$ and mapping irreducible elements to irreducible elements.
\item \emph{Smooth projective models are good} (\lref{UpperTheory.Gal.Seam.goodOf}).
\item \emph{The piece embedding} (\lref{UpperTheory.Gal.Seam.pieceIn_holds}): for an onto cover $C$ over $\OO_K$ of a
  nice tower over $K$, with model $M_C=(Y\to\Spec\OO_K)$, the pull-back $\stepsIHtot((M_C)_{K,\Qlbar})\to\stepsIHtot(Z_{\Qlbar})$
  along $Z\to Y_K$ is injective (after extension of scalars to $\Qlbar$) and $G_K$-equivariant, and the $\theta$-piece lies
  in its image.
\item \emph{Vanishing in negative degrees} (\lref{UpperTheory.Gal.Seam.negVan}).
\item \emph{Admissibility transfer} (\lref{UpperTheory.Gal.Seam.isAdmP_of_isAdm}, with the tie data
  \lref{UpperTheory.Lit.Cmp.locTied6}): for a tied package $P$, $B_{\mathrm{cr}}$-admissibility implies
  $P$-admissibility; and \lref{UpperTheory.Gal.Seam.isCrys_of_isAdmP}: if $r\colon G_K\to\GL_d(\Ql)$ is
  $P$-admissible, its extension of scalars to any finite $LL\supseteq\Ql$ is crystalline for $P$.
\item \emph{Subquotients} (\lref{GlobalTheory.Per.CrysSubQuot.crys_sub_quot} and
  \lref{GlobalTheory.TateDualCrys.isCrys_of_isSubOfSums}): for a tied package, subrepresentations and quotients of
  crystalline representations are crystalline, and a representation embedding into a direct sum of crystalline ones is
  crystalline.
\end{enumerate}

\begin{lemma}[Embedding of the piece into the cohomology of a model]\label{pf:steps1:matin}
Let $K$ be finite over $\Ql$, $H$ a nice tower over $K$, $\theta\colon\Gamma\to\Qlbar$, $C$ an onto cover of $(H,\theta)$ over
$\OO_K$ with model $M_C$, and $\rho_1\colon G_K\to\GL_d(\Qlbar)$ a matrix of the $\theta$-piece. Then there are $N$, bases
$b_n$ of $V_n(M_C)$ $(0\le n<N)$ and a matrix $\mathcal M$ defining an injective map
$\Qlbar^d\to\bigoplus_n\Qlbar^{\dim V_n}$ such that $\bigl(\bigoplus_{n<N}\mathrm{mat}_n(b_n)(\sigma)\bigr)\mathcal M=\mathcal
M\rho_1(\sigma)$ for all $\sigma\in G_K$.
\formallabel{UpperTheory.Gal.Seam.matIn_holds} \formallabel{UpperTheory.Gal.Seam.matInS}
\end{lemma}

\begin{proof}
$M_C$ is proper over $\OO_K$, so its geometric generic fibre is proper over $\Qlbar$, and each $V_n(M_C)$ is
finite-dimensional by \cref{pf:steps1:vfin}; $V_n=0$ for $n<0$ by (g). By (f) the pull-back $\Pi$ from
$\Qlbar\otimes\bigoplus_nV_n(M_C)$ to $\Qlbar\otimes\stepsIHtot(Z_{\Qlbar})$ is injective and equivariant, and the basis vectors
$e_1,\dots,e_d$ of the piece defining $\rho_1$ lie in its image. Their preimages involve only finitely many degrees, say
$n<N$. Writing the preimages in the bases $b_n$ gives the columns of $\mathcal M$; injectivity of $\Pi$ and linear
independence of the $e_j$ give injectivity of $\mathcal M$, and equivariance of $\Pi$ together with
$\sigma^*e_j=\sum_i\rho_1(\sigma)_{ij}e_i$ gives the intertwining relation.
\end{proof}

\begin{lemma}[Subrepresentations of admissible representations]\label{pf:steps1:subadm}
Let $P$ be a tied period package, $B_w$ a bridge for $G_w\subseteq G_E$, $R_n\colon G_{K_w}\to\GL_{a_n}(\Ql)$ $(n<N)$
$P$-admissible homomorphisms, $r\colon G_w\to\GL_d(\Qlbar)$ a homomorphism and $\mathcal M$ an injective matrix with
$\bigl(\bigoplus_nR_n(e(g))\bigr)\mathcal M=\mathcal Mr(g)$ for $g\in G_w$. Then $r$ is crystalline for $(P,B_w)$.
\formallabel{UpperTheory.Gal.Seam.subAdm} \formallabel{UpperTheory.Gal.Seam.isAdmP_block}
\formallabel{GlobalTheory.CrysSubRep.isCrysAt_of_sub}
\end{lemma}

\begin{proof}
The block sum $R=\bigoplus_nR_n$ is $P$-admissible: take block diagonal matrices. Hence the scalar extension of $R$ to the
field $LL$ of the bridge is crystalline by (h). Let $L\supseteq LL$ be the finite extension generated by the entries of
$\mathcal M$, and $\mathcal N$ a left inverse of $\mathcal M$ over $L$. Then $r(g)=\mathcal NR(e(g))\mathcal M$ has entries in
$L$, so $r$ has a model $r'$ over $L$; via $\mathcal M$, $r'$ is a subrepresentation of the $L$-linear crystalline
representation $R_L$, hence crystalline by (i).
\end{proof}

\begin{lemma}[Graded representations]\label{pf:steps1:grcrys}
Let $P$ be tied. If $r$ is graded of $r_0$ and $r_0$ is crystalline for $(P,B_w)$, then so is $r$.
\formallabel{UpperTheory.Gal.Seam.grCrys} \formallabel{UpperTheory.Gal.Seam.isCrys_gr}
\end{lemma}

\begin{proof}
Let $c$ be the grading function and $r_0'$ a crystalline model of $r_0$ over $L$. The entries of $r$ are entries of $r_0$
or $0$, so $r$ has a model $r'$ over $L$, and $r'$ is graded of $r_0'$. For $v\in\N$ let $U_v=\{x:x_j=0\text{ if
}c(j)>v\}$ and $W_v=\{x:x_j=0\text{ if }c(j)\ne v\}$. Block triangularity shows that $U_v$ is $r_0'$-stable and that the
coordinate projection $U_v\to W_v$ intertwines $r_0'|_{U_v}$ with $r'|_{W_v}$ and is surjective. Hence $r'|_{W_v}$ is a
quotient of a subrepresentation of a crystalline representation, so it is crystalline by (i), and
$r'=\bigoplus_vr'|_{W_v}$ is crystalline.
\end{proof}

\begin{theorem}[$\sfname{CrysStep}$]\label{pf:steps1:crys}
Let $E$ be a number field with algebraic closure $\Omega$, $\ell$ a prime, $H$ a nice tower over $E$,
$\theta\colon\Gamma\to\Qlbar$, $\rho\colon G_E\to\GL_d(\Qlbar)$ continuous and attached to the $\theta$-piece, and $w\mid\ell$ a
finite place such that $(H,\theta)$ has an onto cover over $\stepsIOw$. Then $\rho$ is crystalline at $w$.
\formallabel{UpperTheory.Gal.Seam.crysStep_holds} \formallabel{UpperTheory.Gal.Seam.crysStep}
\formallabel{UpperTheory.Gal.Seam.crysStep_of}
\end{theorem}

\begin{proof}
Let $\rho_0$ be a matrix of the $\theta$-piece of which $\rho$ is graded, $C$ the onto cover over $\stepsIOw$, $P$ a tied
package, $A$ a place of $\Omega$ over $w$ and $B_w$ an inducing bridge; put $K=K_w$.
\begin{enumerate}[label=(\arabic*),nosep]
\item By (d) there is $\phi\colon\stepsIOw\to\OO_K$ compatible with $\tau_w$; by (a), $\mathrm{bc}_{\tau_w}H$ is nice; by (c),
  $C$ base-changes to an onto cover $C'$ of $(\mathrm{bc}_{\tau_w}H,\theta)$ over $\OO_K$, with model $M_{C'}$.
\item By (b), with $j\colon\Omega\to\Qlbar$, there is a matrix $\rho_1$ of the $\theta$-piece of $\mathrm{bc}_{\tau_w}H$ for $G_K$
  with $\rho_1(e(g))=\rho_0(g)$ for $g\in D_A$, since $e(g)\circ j=j\circ g$.
\item By \cref{pf:steps1:matin} there are bases $b_n$ of $V_n(M_{C'})$ ($n<N$) and an injective $\mathcal M$ with
  $(\bigoplus_n\mathrm{mat}_n(b_n)(\sigma))\mathcal M=\mathcal M\rho_1(\sigma)$.
\item By (e), $M_{C'}$ is a smooth good model; by \cref{pf:steps1:cst} each $\mathrm{mat}_n(b_n)$ is
  $B_{\mathrm{cr}}$-admissible, hence $P$-admissible by (h), since $P$ is tied.
\item By (2)--(4) and \cref{pf:steps1:subadm} (with $r=\rho_0|_{D_A}$), $\rho_0|_{D_A}$ is crystalline for $(P,B_w)$.
\item Restriction to $D_A$ preserves gradings, so \cref{pf:steps1:grcrys} shows that $\rho|_{D_A}$ is crystalline for $(P,B_w)$.
\end{enumerate}
As $P,A,B_w$ were arbitrary, $\rho$ is crystalline at $w$.
\end{proof}

\begin{remark}
$\sfname{CrysStep}$ needs only an onto cover, not a good one: crystallinity does not involve the Hecke correspondences on the
model, only a smooth projective model through whose cohomology the piece factors. The hypotheses that $\rho$ is continuous and
that $w\mid\ell$ are not used in the proof; at places $w\nmid\ell$ no bridge induces a place over $w$, so the conclusion is
vacuous there.
\end{remark}

% ===========================================================================
\subsection{The Hodge--Tate step}\label{sec:pf:steps1:jump}

In this subsection $\sfname{JumpStep}$ is deduced from the Hodge--Tate comparison theorem. The idea is the
following. Over a finite extension $K$ of $\Ql$, a smooth projective model of the tower gives a
Hodge--Tate decomposition $V_n\otimes\stepsICl\cong\bigoplus_a\stepsIH^{n-a}(\Omega^a)\otimes_K\stepsICl(-a)$ which is
natural in morphisms of models. Applied to finite \'etale Galois covers splitting the Hecke correspondences,
naturality shows that the decomposition intertwines the Hecke operators on \'etale cohomology with the Hecke
operators (trace after pull-back) on Hodge cohomology (\cref{pf:steps1:loce}). Hence a Hecke eigenclass in
$\stepsIH^b(\Omega^a)$ produces an eigenvector of the Galois representation on which $G_K$ acts through
$\stepsIchar^{-a}$, so that $a$ is a Hodge--Tate weight. Since the tower over $E$ only has an integral model
after passing to the quotient $Y=Z/\Delta$ of the cover, the argument is run on the ``low tower'' formed by
$Y$ and the correspondences lying under the Hecke correspondences, and transported back.

Throughout, $K\subset\Qlbar$ is a finite extension of $\Ql$, $\stepsICl$ the completion of $\Qlbar$, $G_K$ acts
on $\stepsICl$, and $\stepsIchar$ is the $\ell$-adic cyclotomic character.

\begin{definition}[Hodge--Tate occurrence]\label{pf:steps1:def:htocc}
Let $\rho\colon G_K\to M_d(\Qlbar)$ be a map and $a\in\Z$. We say that $a$ \emph{occurs as a Hodge--Tate
weight of} $\rho$, and write $a\in\mathrm{HT}(\rho)$, if there are a finite set $S\subset\Qlbar$ and a nonzero
$z\in\stepsICl^d$ such that $\rho(\sigma)\,\sigma(z)=\stepsIchar(\sigma)^{-a}z$ for every $\sigma\in G_K$ fixing
$S$ pointwise. \formal{PositivityTheory.AutomorphicGaloisHodgeTate.G8.HTOcc}
\end{definition}

Thus $a\in\mathrm{HT}(\rho)$ if $\stepsICl(-a)$ occurs in the semilinear representation $\stepsICl^d$ of an
open subgroup of $G_K$; with this convention the cohomology $V_n$ of a smooth projective variety has its
Hodge--Tate weights in $[0,n]$.

\begin{definition}[Frames and Hodge--Tate matrices]\label{pf:steps1:def:frame}
Let $q\colon X\to\Spec K$ be a $K$-scheme and $n\ge0$. A \emph{frame} $F$ of $(X,n)$ consists of a
$\Ql$-basis $b_1,\dots,b_d$ of $V_n(X)=\stepsIHQ^n(X_{\Qlbar})(0)$ and, for $0\le a\le n$, $K$-bases of
$\stepsIH^{n-a}(X,\Omega^a)$ of sizes $e_a$; let $\mathrm{Col}(F)$ be the set of pairs $(a,i)$, $i\le e_a$. A
matrix $\Phi\in M_{d\times\mathrm{Col}(F)}(\stepsICl)$ is \emph{Hodge--Tate} for $F$ if it is invertible and
\[
 \rho_F(\sigma)\,\sigma(\Phi)=\Phi\cdot\mathrm{diag}\bigl(\stepsIchar(\sigma)^{-a}\bigr)_{(a,i)\in\mathrm{Col}(F)}
 \qquad(\sigma\in G_K),
\]
where $\rho_F(\sigma)$ is the matrix of the Galois action on $V_n(X)$ in the basis $b$. For a $K$-morphism
$f\colon X'\to X$ and frames $F,F'$, $\mathrm{et}(f)$ is the matrix of $f^*\colon V_n(X)\to V_n(X')$ and
$\mathrm{hdg}(f)$ the block diagonal matrix of $f^*$ on $\bigoplus_a\stepsIH^{n-a}(\Omega^a)$, both viewed in
$\stepsICl$. \formal{GlobalTheory.HTLink.Frame} \formal{GlobalTheory.HTLink.Frame.IsHT}
\end{definition}

If $\Phi$ is Hodge--Tate, the $\stepsICl$-linear map $\Phi\colon\bigoplus_a\stepsIH^{n-a}(X,\Omega^a)\otimes_K\stepsICl
\to V_n(X)\otimes\stepsICl$ given by the matrix is an isomorphism, and for every $y$ in the $a$-th summand with
coordinates fixed by $\sigma$, $\sigma(\Phi y)=\stepsIchar(\sigma)^{-a}\Phi y$.

\begin{theorem}[Hodge--Tate comparison in the form used]\label{pf:steps1:htspec}
Let $\ell$ be a prime, $K\subset\Qlbar$ finite over $\Ql$ and $n\ge0$.
\begin{enumerate}[label=\textup{(\roman*)},nosep]
\item For every smooth good model $M$ over $\OO_K$ of pure relative dimension there is a frame $F$ of
  $(M_K,n)$ with $\sum_ae_a=d$, i.e.\ $\dim_{\Ql}V_n(M)=\sum_{a=0}^n\dim_K\stepsIH^{n-a}(M_K,\Omega^a)$.
\item There is a choice $\Phi(M,F)$, for all such $M$ and all frames $F$ of $(M_K,n)$, of Hodge--Tate matrices
  for $F$ which is natural: for every morphism of models $u\colon M'\to M$ over $\OO_K$ and frames $F,F'$,
  $\mathrm{et}(u)\,\Phi(M,F)=\Phi(M',F')\,\mathrm{hdg}(u)$.
\end{enumerate}
\formallabel{UpperTheory.Gal.Seam.HT.htSpecGood}
\end{theorem}

\begin{proof}
This is the Hodge--Tate decomposition for smooth proper varieties with good reduction \cite{Faltings1988}
(see also \cite{Scholze2013phodge}): there is a functorial $G_K$-equivariant isomorphism
$V_n(M)\otimes_{\Ql}\stepsICl\cong\bigoplus_a\stepsIH^{n-a}(M_K,\Omega^a)\otimes_K\stepsICl(-a)$. Taking
$\Phi(M,F)$ to be its matrix in the bases of $F$ gives (ii), and comparing dimensions gives (i).
\end{proof}

\begin{remark}[The formal proof]\label{pf:steps1:htspecrem}
The development proves \cref{pf:steps1:htspec} from scratch, in about $26{,}000$ lines, by the route of
\cref{sec:pf:steps1:htdev} below: it constructs a Hodge--Tate pairing and shows that it is injective
(\cref{pf:steps1:htinj}), upgrades it to a bijection by a K\"unneth and diagonal-class argument on self-products together
with a dimension count (\cref{pf:steps1:htbij}), and turns it into natural period matrices using the natural
identification of its Hodge cohomology with sheaf cohomology (\cref{pf:steps1:cech};
\lref{UpperTheory.Gal.Seam.Spec.htSpecGood_of}, deferred, chooses the frames and verifies the naturality of the
resulting matrices).
\end{remark}

\subsubsection{The Hodge--Tate comparison of the development}\label{sec:pf:steps1:htdev}

We prove the formal Hodge--Tate statements on which the formal proof of \cref{pf:steps1:htspec} rests. The construction of
the pairing is the comparison theorem itself, and we cite the literature for it; the injectivity and the bijectivity are
deduced in the development by arguments of its own, which we give in full. Throughout, $K\subset\Qlbar$ is finite over
$\Ql$. For a scheme $X$ over a field $k$ and $a,n\ge0$ let $\stepsIHdg^{a,n}(X)$ be the Hodge cohomology used by the
development: the degree-$n$ cohomology, in the \v{C}ech formalism of the development, of the $a$-th graded piece
$\Omega^a_{X/k}[-a]$ of the de Rham complex; by \cref{pf:steps1:cech} below, $\stepsIHdg^{a,a+j}(X)\cong
\stepsIH^j(X,\Omega^a_{X/k})$ for smooth proper $X$. It carries pull-backs, $k$-bilinear cup products
$\stepsIHdg^{a,i}\times\stepsIHdg^{a',j}\to\stepsIHdg^{a+a',i+j}$ and a unit $1\in\stepsIHdg^{0,0}$. For a model $M$ we
write $\stepsIHdg^{a,n}(M)=\stepsIHdg^{a,n}(M_K)$. Classes $x\in T_i(M)$, $y\in T_j(M)$ have cup products $z\in
T_{i+j}(M)$ (levelwise cup products of compatible families).

\begin{definition}[Hodge--Tate pairings]\label{pf:steps1:def:htpre}
A \emph{Hodge--Tate pairing} $P$ for $K$ consists of biadditive maps
$h=h_M^{a,n}\colon V_n(M)\times\stepsIHdg^{a,n}(M)^\vee\to\stepsICl$, for all models $M$ and $a,n\ge0$, which are
$\Ql$-linear in the first and $K$-linear in the second variable, such that for all smooth good models $M$, $N$:
\begin{enumerate}[label=\textup{(\roman*)},nosep]
\item $\stepsIHdg^{a,n}(M)$ is finite-dimensional, and zero if $n<a$;
\item $\stepsIchar(\sigma)^a\,h(\sigma v,\phi)=\sigma\bigl(h(v,\phi)\bigr)$ for $\sigma\in G_K$;
\item $h_M(u^*v,\phi)=h_N(v,\phi\circ u^*)$ for every morphism $u\colon M.X\to N.X$ over $\OO_K$;
\item if $z\in T_{i+j}(M)$ is a cup product of $x\in T_i(M)$ and $y\in T_j(M)$, then for all bases $(e_{a,s})_s$ of
  $\stepsIHdg^{a,i}(M)$ and $(e'_{a,s'})_{s'}$ of $\stepsIHdg^{a,j}(M)$, with dual bases $(e^*_{a,s})$, $(e'^*_{a,s'})$,
  \begin{equation}\label{pf:steps1:eq:htcup}
   h^{c,i+j}(z,\phi)=\sum_{a+a'=c}\ \sum_{s,s'}h^{a,i}(x,e^*_{a,s})\,h^{a',j}(y,e'^*_{a',s'})\,\phi(e_{a,s}\cup e'_{a',s'});
  \end{equation}
\item $h^{0,0}(1,\phi)=\phi(1)$ for the unit class $1\in T_0(M)$;
\item for every line datum $D$ on $M.X$ with \'etale structure map whose trivializing opens cover $M.X$, $D$ has \'etale first
  Chern classes $v\in T_2(M)$ (for every choice of roots of unity), its generic fibre has a unique logarithmic first Chern
  class $c\in\stepsIHdg^{1,2}(M)$, and $h^{1,2}(v,\phi)=\phi(c)$, $h^{a,2}(v,\cdot)=0$ for $a\ne1$.
\end{enumerate}
Let $M$ be smooth good. $P$ is \emph{injective in degree $n$} for $M$ if for every basis $(b_k)$ of $V_n(M)$ and every
$z\in\stepsICl^{\dim V_n}$ with $\sum_kz_kh^{a,n}(b_k,\phi)=0$ for all $a$ and $\phi$ one has $z=0$; for $M$ connected of
dimension $d$, $P$ is \emph{injective in top degree} if the same holds with $n=2d$ and only $a=d$; and $P$ is
\emph{surjective in degree $n$} if for all bases $(e_{a,s})_s$ of $\stepsIHdg^{a,n}(M)$ and all $z_{a,s}\in\stepsICl$ with
$\sum_{a\le n}\sum_sz_{a,s}h^{a,n}(v,e^*_{a,s})=0$ for every $v\in V_n(M)$, all $z_{a,s}$ vanish.
\formal{CrystallineSemistableComparison.HT.HTPre} \formal{CrystallineSemistableComparison.HT.HTPre.Inj}
\formal{CrystallineSemistableComparison.HT.HTPre.Top} \formal{CrystallineSemistableComparison.Onto.Surj}
\end{definition}

Thus injectivity in degree $n$ means that $v\mapsto\sum_{a,s}h(v,e^*_{a,s})e_{a,s}$ is an injective map
$V_n(M)\otimes\stepsICl\to\bigoplus_a\stepsIHdg^{a,n}(M)\otimes_K\stepsICl$, and surjectivity that this map is
surjective. For the classical comparison isomorphism of \cite{Faltings1988} both hold.

\begin{lemma}[\v{C}ech comparison]\label{pf:steps1:cech}
Let $k$ be a field of characteristic $0$. There are $k$-linear isomorphisms
$e_X\colon\stepsIHdg^{a,a+j}(X)\to\stepsIH^j(X,\Omega^a_{X/k})$, for all smooth proper $k$-schemes $X$ and all $a,j\ge0$, such
that $e_X(u^*x)=u^*e_Y(x)$ for every $k$-morphism $u\colon X\to Y$.
\formallabel{CrystallineSemistableComparison.Om.Nt2F.hodge_iso_natural_cech}
\end{lemma}

\begin{proof}
For every finite cover $\mathcal U$ of $X$ by affine opens whose finite intersections are affine, the development
identifies $\stepsIHdg^{a,a+j}(X)$ with the $j$-th cohomology of the \v{C}ech complex $\check C^\bullet(\mathcal U,\Omega^a)$,
compatibly with pull-backs along morphisms $u$ and refinements $\mathcal U\to u^{-1}\mathcal W$
(\lref{CrystallineSemistableComparison.Om.GlobPf2.glob_iso}, deferred; this is internal to the \v{C}ech formalism of the
development, whose graded pieces are built from differential forms on affine pieces). Since $X$ is noetherian and
separated and $\Omega^a$ is quasi-coherent, the natural map $\check\stepsIH^j(\mathcal U,\Omega^a)\to\stepsIH^j(X,\Omega^a)$ is
an isomorphism \cite[III.4.5]{Hartshorne1977}; these maps are compatible with refinements and with pull-backs along
morphisms $u$ for which $\mathcal U$ refines $u^{-1}\mathcal W$ \cite[III.4.4]{Hartshorne1977}. Composing gives $e_X$, and
it is independent of the cover because any two covers have a common refinement.
\end{proof}

\begin{lemma}[Top powers in Hodge cohomology]\label{pf:steps1:ktop}
Let $X$ be a connected smooth scheme of relative dimension $d$ over a field $k$ of characteristic $0$, $D$ a tautological
line datum on $X$ and $c\in\stepsIHdg^{1,2}(X)$ a logarithmic first Chern class of $D$. Then $c^d\ne0$ in
$\stepsIHdg^{d,2d}(X)$.
\formallabel{CrystallineSemistableComparison.HT.kTop}
\end{lemma}

\begin{proof}
Let $c'$ be a de Rham first Chern class of $D$ (\lref{CrystallineSemistableComparison.Pb.isFirstChernDeRham_exists},
deferred). Choose a cover of $X$ by $d+1$ affine opens trivializing $D$; on it both $c$ and $c'$ are represented by the
same \v{C}ech cocycle, formed by the logarithmic derivatives of the transition functions of $D$, so $c^d$ and $c'^d$ are
represented by the same cocycle of $d$-forms on the $(d+1)$-fold intersection. In total degree $2d$ the de Rham cohomology
of $X$ computed on this cover is a quotient of the \v{C}ech cohomology $\check\stepsIH^d(\Omega^d)$, the other terms of the
double complex vanishing for degree reasons, and $c'^d$ is the image of $c^d$. As $c'^d\ne0$ by \cref{pf:steps1:powK}, $c^d\ne0$.
\end{proof}

\begin{theorem}[Hodge--Tate pairings are injective]\label{pf:steps1:htinj}
There is a Hodge--Tate pairing $P$ for $K$, and every Hodge--Tate pairing is injective in every degree for every smooth good
model and injective in top degree for every smooth good model connected of some dimension.
\formallabel{UpperTheory.Gal.Seam.HT.hodgeTate_comparison_smooth_projective}
\formallabel{CrystallineSemistableComparison.HT.hodgeTate_of_etale_side}
\formallabel{CrystallineSemistableComparison.HT.ht_of}
\formallabel{CrystallineSemistableComparison.HT.inj_of}
\end{theorem}

\begin{proof}
\emph{Existence.} The Hodge--Tate decomposition of the \'etale cohomology of smooth proper varieties over $K$, with its
functoriality and its compatibility with cup products and Chern classes \cite{Faltings1988} (see also
\cite{Scholze2013phodge}), provides a pairing with properties (ii)--(vi) of \cref{pf:steps1:def:htpre} (the logarithmic
first Chern class of (vi) is the Hodge component of $c_1$; the sign convention may be adjusted by multiplying $h^{a,n}$ by
$(-1)^a$, which preserves (ii)--(v)); (i) is the finiteness of coherent cohomology of proper schemes and the vanishing of
$\stepsIH^{n-a}$ for $n<a$. The development constructs such a pairing from its own comparison maps
(\lref{CrystallineSemistableComparison.HT.pre_of}, deferred, which assembles it from lifts of \'etale classes through
the Hodge--Tate filtration of the period ring, \lref{CrystallineSemistableComparison.HT.half1}, and from maps of graded
pieces to Hodge cohomology, \lref{CrystallineSemistableComparison.HT.half2}, both deferred).

Now let $P$ be any Hodge--Tate pairing. We use the \'etale side for smooth good models
(\lref{UpperTheory.Gal.Seam.HT.etSideS}): cup products exist, $V_n(M)$ is finite-dimensional, every smooth good model has
finitely many smooth good models $u_\alpha\colon M_\alpha\to M$, each connected of some dimension, such that
$v\in V_n(M)$ vanishes if all $u_\alpha^*v$ vanish, and \emph{Poincar\'e duality} holds for $M$ connected of dimension $d$:
$V_i(M)=0$ for $i>2d$, and for $i+j=2d$ there are $g_1,\dots,g_r$ spanning $V_{2d}(M)$, permuted by $G_K$ up to nonzero
scalars (for all $s,s'$ there are $\sigma\in G_K$ and $\zeta\in\Ql^\times$ with $\sigma g_s=\zeta g_{s'}$), such that the
linear forms $f$ on $V_i(M)$ for which there are $s$ and $y\in T_j(M)$ with $x\cup y=f(x)\,g_s$ for all $x\in T_i(M)$ span
the dual of $V_i(M)$.

\emph{Step 1: a nonzero top value.} Let $M$ be connected of dimension $d$. There are $v\in V_{2d}(M)$ and $\phi$ with
$h^{d,2d}(v,\phi)\ne0$. Indeed, $M.X$ carries a tautological line datum $D$ with \'etale structure map whose trivializing opens
cover $M.X$ (\lref{CrystallineSemistableComparison.Pb.isO1_exists},
\lref{CrystallineSemistableComparison.Pb.isO1_etale}); let $v_1\in T_2(M)$ and $c\in\stepsIHdg^{1,2}(M)$ be as in (vi).
By induction on $k$, using \eqref{pf:steps1:eq:htcup} and (v), a $k$-fold cup power $w_k\in T_{2k}(M)$ of $v_1$ satisfies
$h^{k,2k}(w_k,\phi)=\phi(c^k)$ and $h^{a,2k}(w_k,\cdot)=0$ for $a\ne k$: in \eqref{pf:steps1:eq:htcup} for $w_{k+1}=w_k\cup v_1$ only
the terms with $a=k$, $a'=1$ survive, and they sum to $\phi(c^k\cup c)$. By \cref{pf:steps1:ktop}, $c^d\ne0$, so
$\phi(c^d)=1$ for some $\phi$, and $v=w_d$ works.

\emph{Step 2: duality.} Let $M$ be connected of dimension $d$, $n+j=2d$, $(b_k)$ a basis of $V_n(M)$ and $z\in\stepsICl^{\dim V_n}$
such that $\sum_kz_kh^{p,n}(b_k,\psi)=0$ for all $\psi$ and all $p$ with $d-p\le j$. Then $z=0$. Take $v,\phi$ as in Step 1 and
write $v$ as a combination of the $g_s$; then $h^{d,2d}(g_{s_0},\phi)\ne0$ for some $s_0$, and by (ii) and the transitivity of the
Galois action on the $g_s$ up to scalars, $h^{d,2d}(g_s,\phi)\ne0$ for all $s$. Let $f$, $s$, $y$ be as in Poincar\'e duality.
By \eqref{pf:steps1:eq:htcup} applied to $b_k\cup y=f(b_k)g_s$,
\[
 \Bigl(\sum_kz_kf(b_k)\Bigr)h^{d,2d}(g_s,\phi)=\sum_{p+q=d}\ \sum_{t,t'}\Bigl(\sum_kz_kh^{p,n}(b_k,e^*_{p,t})\Bigr)h^{q,j}(y,e'^*_{q,t'})\,
 \phi(e_{p,t}\cup e'_{q,t'})=0,
\]
because the inner sums vanish for $q\le j$, while $\stepsIHdg^{q,j}(M)=0$ for $q>j$ by (i). Hence $\sum_kz_kf(b_k)=0$ for a
spanning set of linear forms $f$ on $V_n(M)$, and $z=0$ since $(b_k)$ is a basis.

\emph{Step 3: conclusion.} For $M$ connected of dimension $d$, injectivity in degree $n$ is trivial for $n>2d$ ($V_n(M)=0$) and
follows from Step 2 with $j=2d-n$ otherwise; injectivity in top degree is Step 2 with $n=2d$, $j=0$, where the only
condition needed is the one for $p=d$. For a general smooth good $M$, let $z$ satisfy the hypothesis of injectivity in degree
$n$ and put $x=\sum_kz_kb_k\in V_n(M)\otimes\stepsICl$. By (iii), $\sum_kz_kh_{M_\alpha}(u_\alpha^*b_k,\phi')=
\sum_kz_kh_M(b_k,\phi'\circ u_\alpha^*)=0$ for all $\alpha$ and $\phi'$, so expanding $u_\alpha^*b_k$ in a basis of
$V_n(M_\alpha)$ and using injectivity for the connected model $M_\alpha$ gives $u_\alpha^*x=0$ for all $\alpha$; as the
$u_\alpha^*$ jointly detect $V_n(M)$, also after the flat extension of scalars to $\stepsICl$, $x=0$, i.e.\ $z=0$.
\end{proof}

\begin{lemma}[Dimension count]\label{pf:steps1:rank}
Let $P$ be a Hodge--Tate pairing, $M$ a smooth good model and $n\ge0$ such that $V_n(M)$ is finite-dimensional and $P$ is
injective and surjective in degree $n$ for $M$. Then $\dim_{\Ql}V_n(M)=\sum_{a=0}^n\dim_K\stepsIHdg^{a,n}(M)$.
\formallabel{CrystallineSemistableComparison.Onto.rank_equality}
\end{lemma}

\begin{proof}
Choose bases $(b_k)$ of $V_n(M)$ and $(e_{a,s})_s$ of $\stepsIHdg^{a,n}(M)$, $a\le n$, and consider the matrix
$H=(h^{a,n}(b_k,e^*_{a,s}))_{k,(a,s)}$ over $\stepsICl$. Since every $\phi$ is a $K$-combination of the $e^*_{a,s}$ (and
$\stepsIHdg^{a,n}=0$ for $a>n$), a relation $\sum_kz_kH_{k,(a,s)}=0$ for all $(a,s)$ gives $\sum_kz_kh^{a,n}(b_k,\phi)=0$ for all
$a,\phi$, hence $z=0$ by injectivity: the rows of $H$ are linearly independent. Since $h$ is $\Ql$-linear in $v$, a relation
$\sum_{a,s}z_{a,s}H_{k,(a,s)}=0$ for all $k$ gives the hypothesis of surjectivity, hence $z=0$: the columns are linearly
independent. So $H$ is square.
\end{proof}

\begin{lemma}[Retracts]\label{pf:steps1:retract}
\begin{enumerate}[label=\textup{(\roman*)},nosep]
\item Every smooth good model $M$ of relative dimension $d$ is a retract of a smooth good model $N$ of relative dimension
  $d+2$: there are morphisms $p\colon N.X\to M.X$ and $s\colon M.X\to N.X$ over $\OO_K$ with $ps=\mathrm{id}$.
\item If $M$ is a retract of $N$ as in \textup{(i)} and a Hodge--Tate pairing $P$ is surjective in degree $n$ for $N$, then it is
  surjective in degree $n$ for $M$.
\end{enumerate}
\formallabel{CrystallineSemistableComparison.Onto.retract_pf}
\formallabel{CrystallineSemistableComparison.Onto.ontoLow_pf}
\end{lemma}

\begin{proof}
(i) Take $N.X=M.X\times_{\OO_K}\mathbb P^2_{\OO_K}$ with $p$ the projection and $s=(\mathrm{id},x_0)$ for an $\OO_K$-point $x_0$ of
$\mathbb P^2$. It is smooth of relative dimension $d+2$, projective (products of projective morphisms are projective), and
strictly semistable for the same uniformizer, as every smooth $\OO_K$-scheme is.

(ii) Let $(e_{a,t})_t$ be bases of $\stepsIHdg^{a,n}(M)$ and $z_{a,t}$ with $\sum_{a,t}z_{a,t}h_M(v,e^*_{a,t})=0$ for all $v\in V_n(M)$.
Let $\Lambda=\sum_{a,t}z_{a,t}\,e^*_{a,t}\circ s^*$, a $\stepsICl$-valued functional on $\bigoplus_a\stepsIHdg^{a,n}(N)$ (via
$K\subset\stepsICl$). By (iii) of \cref{pf:steps1:def:htpre}, $\sum_{a,t}z_{a,t}h_N(w,e^*_{a,t}\circ s^*)=
\sum_{a,t}z_{a,t}h_M(s^*w,e^*_{a,t})=0$ for all $w\in V_n(N)$. Expanding $e^*_{a,t}\circ s^*$ in the dual basis of a basis
of $\stepsIHdg^{a,n}(N)$ and using $K$-linearity of $h$ in the second variable, surjectivity for $N$ shows $\Lambda=0$.
Composing with $p^*$ and using $s^*p^*=(ps)^*=\mathrm{id}$ gives $\sum_{a,t}z_{a,t}e^*_{a,t}=0$, so all $z_{a,t}$ vanish.
\end{proof}

The surjectivity is proved with the class of the diagonal. We collect the inputs.

\begin{lemma}[Inputs for the diagonal argument]\label{pf:steps1:diaginputs}
Let $M$ be a smooth good model of relative dimension $d+1$ with $d\ge1$.
\begin{enumerate}[label=\textup{(\roman*)},nosep]
\item \textup{(The blown-up square)} There are a smooth good model $Q$ with $Q.X=M.X\times_{\OO_K}M.X$, projections $p_1,p_2$ and
  diagonal $\Delta$, a smooth good model $Y$ with a morphism $\pi\colon Y.X\to Q.X$ over $\OO_K$, and a line datum $E$ on $Y.X$
  with \'etale structure map whose trivializing opens cover $Y.X$, such that on generic fibres $\pi$ is the blow-up of $Q_K$
  along $\Delta(M_K)$ with exceptional line datum $E_K$: $\pi$ is an isomorphism over the complement $U$ of the diagonal, $E$
  has a section which is invertible over $\pi^{-1}U$, and over finitely many affine opens covering the diagonal $\pi$ is given by
  the standard blow-up charts of codimension $d+1$, the section being the standard local equation.
\item \textup{(The diagonal class)} For every \'etale first Chern class $v\in T_2(Y)$ of $E$ there are $c\in T_{2d+2}(Q)$ and
  $w\in T_{2d}(Y)$ such that $\pi^*c$ and the cup product $v\cup w$ have the same image in $V_{2d+2}(Y)$, and finitely many
  $a_0(s)\in T_{2d+2}(M)$ with linearly independent images in $V_{2d+2}(M)$, $b_0(s)\in T_0(M)$ with $\sum_sb_0(s)=N\cdot1$
  in $V_0(M)$ for some $N\in\Ql^\times$, and $a(s)\in T_{i(s)}(M)$, $b(s)\in T_{j(s)}(M)$ with $i(s)+j(s)=2d+2$ and
  $j(s)\ge1$, such that in $V_{2d+2}(Q)$
  \begin{equation}\label{pf:steps1:eq:diagdec}
   c=\sum_sp_1^*a_0(s)\cup p_2^*b_0(s)+\sum_sp_1^*a(s)\cup p_2^*b(s).
  \end{equation}
\item \textup{(Finiteness)} $V_n(N)$ is finite-dimensional for every smooth good model $N$ and every $n$.
\item \textup{(Hodge cohomology)} For schemes over a field $k$ of characteristic $0$: cup products on $\stepsIHdg$ are bilinear,
  associative, unital, graded commutative ($x\cup y=(-1)^{ij}y\cup x$ for $x\in\stepsIHdg^{a,i}$, $y\in\stepsIHdg^{a',j}$) and
  compatible with pull-backs. \emph{Slant products:} if $X$, $X'$ are smooth and proper over $k$, $R=X\times_kX'$ with projections
  $p_1,p_2$, and $\phi\in\stepsIHdg^{a,i}(X)^\vee$, then for all $a',j$ there is a $k$-linear map
  $\mathrm{sl}\colon\stepsIHdg^{a+a',i+j}(R)\to\stepsIHdg^{a',j}(X')$ with $\mathrm{sl}(p_1^*u\cup p_2^*u')=\phi(u)u'$ for
  $u\in\stepsIHdg^{a,i}(X)$, $u'\in\stepsIHdg^{a',j}(X')$, and $\mathrm{sl}(p_1^*u\cup p_2^*u')=0$ for $u\in\stepsIHdg^{b,i'}(X)$,
  $u'\in\stepsIHdg^{b',j'}(X')$ with $b+b'=a+a'$, $i'+j'=i+j$ and $(b,i')\ne(a,i)$. \emph{Vanishing on blow-ups:} if $Z$ is smooth
  and proper over $k$ and $\pi\colon Y'\to Z$ is a blow-up along a closed immersion $\iota\colon W\to Z$ with exceptional line
  datum $E'$ as in \textup{(i)}, of positive codimension, $e$ is a logarithmic first Chern class of $E'$, $c'\in
  \stepsIHdg^{a,n}(Z)$ satisfies $\pi^*c'=0$ or $\pi^*c'=e\cup w'$ for some $w'$, and $y\in\stepsIHdg^{a',n'}(Z)$ satisfies
  $\iota^*y=0$, then $c'\cup y=0$.
\end{enumerate}
\formallabel{UpperTheory.Gal.KunnethLefschetz.Top.diagEt3_holds}
\formallabel{UpperTheory.Gal.KunnethLefschetz.Top.etFin_holds}
\end{lemma}

\begin{proof}
(i) is \lref{UpperTheory.Gal.KunnethLefschetz.G.sqEx3} (deferred; it constructs the blow-up of $M.X\times_{\OO_K}M.X$ along
the diagonal, which is smooth and projective over $\OO_K$ since the diagonal is a smooth closed subscheme, by explicit
charts). (ii) and (iii) are derived in the development from the \'etale propositions C2--C5 and D1, which hold
(\lref{UpperTheory.Gal.KunnethLefschetz.core_holds}, stating this), by
\lref{UpperTheory.Gal.KunnethLefschetz.Core.diagEt3} and \lref{UpperTheory.Gal.KunnethLefschetz.Core.etFin};
classically $c$ is the cycle class of the diagonal, which pulls back along $\pi$ to a multiple of the class of the exceptional
divisor, and \eqref{pf:steps1:eq:diagdec} is its K\"unneth decomposition. (iv) is derived from
\lref{CrystallineSemistableComparison.Onto.parts} (deferred), which establishes thirty-one properties of Hodge cohomology
of smooth schemes over fields of characteristic $0$ (Mayer--Vietoris and descent along covers, the structure of differential
forms on affine pieces, logarithmic derivatives, K\"unneth formulas, cup products, first Chern classes and blow-up charts);
classically, the slant product is $\phi\otimes\mathrm{id}$ on the K\"unneth decomposition of the coherent cohomology of
$\Omega^\bullet_R$, and the vanishing on blow-ups reflects that $e$ is supported on the exceptional divisor, which maps to $W$.
\end{proof}

From now on let $M$, $Q$, $Y$, $\pi$, $E$ be as in \cref{pf:steps1:diaginputs}(i), let $v\in T_2(Y)$ be an \'etale first Chern
class of $E$ (it exists by (vi) of \cref{pf:steps1:def:htpre}), $e\in\stepsIHdg^{1,2}(Y)$ the logarithmic first Chern class of
$E_K$, and $c$, $w$, $a_0$, $b_0$, $a$, $b$, $N$ as in \cref{pf:steps1:diaginputs}(ii). We fix a Hodge--Tate pairing
$P$; its maps are written $h_M$, $h_Q$, $h_Y$.

\begin{lemma}[Classes vanishing on the diagonal]\label{pf:steps1:kill}
For all $A,a',n'\ge0$, every $y\in\stepsIHdg^{a',n'}(Q)$ with $\Delta^*y=0$ and every linear form $\Theta$ on
$\stepsIHdg^{A+a',2d+2+n'}(Q)$, $h_Q^{A,2d+2}\bigl(c,\Theta(\,\cdot\cup y)\bigr)=0$.
\formallabel{CrystallineSemistableComparison.Onto.stepA}
\formallabel{CrystallineSemistableComparison.Onto.mkDat}
\end{lemma}

\begin{proof}
Let $\Theta'=\Theta(\,\cdot\cup y)$, a linear form on $\stepsIHdg^{A,2d+2}(Q)$. By the vanishing on blow-ups, $\Theta'(C)=0$ whenever
$\pi^*C=0$ or $\pi^*C=e\cup w'$ for some $w'$. Hence there is a linear form $\phi$ on $\stepsIHdg^{A,2d+2}(Y)$ with
$\Theta'=\phi\circ\pi^*$ which vanishes on $e\cup\stepsIHdg^{A-1,2d}(Y)$: define $\phi$ on $\pi^*\stepsIHdg^{A,2d+2}(Q)+
e\cup\stepsIHdg^{A-1,2d}(Y)$ by $\phi(\pi^*C+e\cup w')=\Theta'(C)$, which is well defined by what was just said, and extend it
linearly. By (iii) of \cref{pf:steps1:def:htpre} and \cref{pf:steps1:diaginputs}(ii), $h_Q(c,\Theta')=h_Y(\pi^*c,\phi)=
h_Y(v\cup w,\phi)$, and by \eqref{pf:steps1:eq:htcup} and (vi) only the terms with $a=1$ survive:
$h_Y(v\cup w,\phi)=\sum_{s'}h_Y^{A-1,2d}(w,e'^*_{s'})\,\phi(e\cup e'_{s'})=0$.
\end{proof}

For $x\in\stepsIHdg^{a',n}(M)$, $\alpha\ge0$ and $\phi'\in\stepsIHdg^{\alpha,2d+2}(M)^\vee$ put, for $X\in T_i(M)$,
\[
 \kappa(X)=(-1)^{n}\,h_M^{\alpha-a',i}\bigl(X,\phi'(\,\cdot\cup x)\bigr)\quad\text{if }a'\le\alpha\text{ and }i+n=2d+2,
 \qquad\kappa(X)=0\text{ otherwise},
\]
and for $\psi\in\stepsIHdg^{\beta,n}(M)^\vee$ let $G_x(\psi)\in\stepsIHdg^{0,0}(M)^\vee$ be $u\mapsto\psi(u\cup x)$ if $\beta=a'$ and $0$
otherwise.

\begin{lemma}[The diagonal identity]\label{pf:steps1:eid}
For all $x$, $\alpha$, $\phi'$ as above, all $\beta$ and all $\psi\in\stepsIHdg^{\beta,n}(M)^\vee$,
\begin{equation}\label{pf:steps1:eq:eid}
\begin{aligned}
 &\sum_s[n=0]\,\kappa(a_0(s))\,h_M^{\beta,n}(b_0(s),\psi)+\sum_s[j(s)=n]\,\kappa(a(s))\,h_M^{\beta,n}(b(s),\psi)\\
 &\qquad=\sum_sh_M^{\alpha,2d+2}(a_0(s),\phi')\,h_M^{0,0}\bigl(b_0(s),G_x(\psi)\bigr),
\end{aligned}
\end{equation}
where $[\,\cdot\,]$ is $1$ if the condition holds and $0$ otherwise.
\formallabel{CrystallineSemistableComparison.Onto.eid}
\end{lemma}

\begin{proof}
If $a'>\alpha+\beta$, every $\kappa$ and $G_x(\psi)$ vanish and both sides are $0$. Otherwise let $A=\alpha+\beta-a'$, let
$\mathrm{sl}\colon\stepsIHdg^{A+a',2d+2+n}(Q)\to\stepsIHdg^{\beta,n}(M)$ be the slant product with $\phi'$
(\cref{pf:steps1:diaginputs}(iv)) and $\Theta=\psi\circ\mathrm{sl}$. The class $y=p_1^*x-p_2^*x$ satisfies
$\Delta^*y=x-x=0$, so by \cref{pf:steps1:kill}
\[
 h_Q\bigl(c,\Theta(\,\cdot\cup p_1^*x)\bigr)=h_Q\bigl(c,\Theta(\,\cdot\cup p_2^*x)\bigr).
\]
We expand $c$ by \eqref{pf:steps1:eq:diagdec} and compute the terms. Let $X\in T_i(M)$, $Y'\in T_j(M)$ with $i+j=2d+2$. By
\eqref{pf:steps1:eq:htcup} on $Q$ and (iii) for $p_1,p_2$,
\[
 h_Q\bigl(p_1^*X\cup p_2^*Y',\Theta''\bigr)=\sum_{b+b'=A}\ \sum_{t,t'}h_M^{b,i}(X,e^*_{b,t})\,h_M^{b',j}(Y',e'^*_{b',t'})\,
 \Theta''\bigl(p_1^*e_{b,t}\cup p_2^*e'_{b',t'}\bigr)
\]
for every linear form $\Theta''$. For $\Theta''=\Theta(\,\cdot\cup p_1^*x)$, associativity and graded commutativity give
$p_1^*u\cup p_2^*u'\cup p_1^*x=(-1)^{jn}p_1^*(u\cup x)\cup p_2^*u'$, so by the properties of $\mathrm{sl}$ the term is nonzero only
if $u\cup x$ has bidegree $(\alpha,2d+2)$, i.e.\ $b=\alpha-a'$ and $j=n$; summing, with $K$-linearity of $h$ in the second
variable, gives $[j=n]\,\kappa(X)\,h_M^{\beta,n}(Y',\psi)$. For $\Theta''=\Theta(\,\cdot\cup p_2^*x)$, $p_1^*u\cup p_2^*u'\cup
p_2^*x=p_1^*u\cup p_2^*(u'\cup x)$, and the term is nonzero only if $u$ has bidegree $(\alpha,2d+2)$, i.e.\ $i=2d+2$ and $j=0$;
for $j=0$ the sum is $h_M^{\alpha,2d+2}(X,\phi')\,h_M^{0,0}(Y',G_x(\psi))$, and for $j\ge1$ it is $0$. Inserting these values for
the terms of \eqref{pf:steps1:eq:diagdec} gives \eqref{pf:steps1:eq:eid}.
\end{proof}

\begin{proposition}[Surjectivity from top-degree injectivity]\label{pf:steps1:onto}
Assume the conclusions of \cref{pf:steps1:diaginputs}. Let $P$ be a Hodge--Tate pairing and $M$ a smooth good model of relative
dimension $d+1$ with $d\ge1$ such that $P$ is injective in degree $2d+2$ for $M$. Then $P$ is surjective in every degree for $M$.
\formallabel{CrystallineSemistableComparison.Onto.ontoOf3_pf}
\formallabel{CrystallineSemistableComparison.Onto.core}
\formallabel{CrystallineSemistableComparison.Onto.surj_of_core}
\end{proposition}

\begin{proof}
Fix $n$. For $v\in V_n(M)$ let $g_v$ be the family of linear forms $(\beta,\psi)\mapsto h_M^{\beta,n}(v,\psi)$, and let $S$ be
the $\stepsICl$-span of the $g_v$ in the space of all families $(\beta,\psi)\mapsto\stepsICl$. For $x\in\stepsIHdg^{a',n}(M)$
let $t_x$ be the family $(\beta,\psi)\mapsto N\psi(x)$ for $\beta=a'$ and $0$ for $\beta\ne a'$.

\emph{Claim: $t_x\in S$.} For every $\alpha$ and $\phi'$ the left side of \eqref{pf:steps1:eq:eid}, as a function of
$(\beta,\psi)$, is a combination of the $g_{b_0(s)}$ and $g_{b(s)}$, so the right side $R_{\phi'}$ lies in $S$. Let $\Phi$ map
$(c_s)\in\stepsICl^{m_0}$ to the family $(\beta,\psi)\mapsto\sum_sc_s\,h_M^{0,0}(b_0(s),G_x(\psi))$; then
$R_{\phi'}=\Phi(\mathrm{vec}(\phi'))$ with $\mathrm{vec}(\phi')=(h_M^{\alpha,2d+2}(a_0(s),\phi'))_s$. The vector $(1,\dots,1)$ lies
in the span of the vectors $\mathrm{vec}(\phi')$: otherwise some $(z_s)$ satisfies $\sum_sz_sh_M^{\alpha,2d+2}(a_0(s),\phi')=0$ for
all $\alpha,\phi'$ and $\sum_sz_s\ne0$; extending the linearly independent $a_0(s)$ to a basis of the finite-dimensional
$V_{2d+2}(M)$ and the coefficients by $0$, injectivity in degree $2d+2$ gives $z=0$, a contradiction. Hence
$\Phi(1,\dots,1)\in S$, and by (v) of \cref{pf:steps1:def:htpre} and $\sum_sb_0(s)=N\cdot1$,
$\Phi(1,\dots,1)(\beta,\psi)=h_M^{0,0}(N\cdot1,G_x(\psi))=N\,G_x(\psi)(1)=t_x(\beta,\psi)$.

\emph{Surjectivity.} Let $(e_{a,t})_t$ be bases of $\stepsIHdg^{a,n}(M)$ and $z_{a,t}$ with $\sum_{a\le n}\sum_tz_{a,t}
h_M^{a,n}(v,e^*_{a,t})=0$ for all $v$. The linear map $L(g)=\sum_{a\le n}\sum_tz_{a,t}\,g(a,e^*_{a,t})$ vanishes on every $g_v$,
hence on $S$, hence on $t_x$ for $x=e_{a',t_0}$; but $L(t_x)=N\,z_{a',t_0}$. As $N\ne0$, $z_{a',t_0}=0$.
\end{proof}

\begin{theorem}[Hodge--Tate bijectivity]\label{pf:steps1:htbij}
There is a Hodge--Tate pairing $P$ for $K$ such that for every smooth good model $M$: $P$ is injective in every degree; if $M$ is
connected of dimension $d$, $P$ is injective in top degree; and if $M$ is smooth of pure relative dimension $d$, then for every
$n$, $P$ is surjective in degree $n$ and $\dim_{\Ql}V_n(M)=\sum_{a=0}^n\dim_K\stepsIHdg^{a,n}(M)$.
\formallabel{UpperTheory.Gal.Seam.HT.htBij}
\formallabel{CrystallineSemistableComparison.Onto.htBij_of}
\formallabel{CrystallineSemistableComparison.Onto.htBij_of_parts}
\formallabel{CrystallineSemistableComparison.Onto.ontoAll_of}
\end{theorem}

\begin{proof}
Let $P$ be as in \cref{pf:steps1:htinj}; it is injective in every degree for every smooth good model, and injective in top
degree for connected ones. Let $M$ be smooth of relative dimension $d$. By \cref{pf:steps1:retract}(i), $M$ is a retract of a
smooth good model $N$ of relative dimension $d+2=(d+1)+1$ with $d+1\ge1$; $P$ is injective in degree $2(d+2)$ for $N$, so it is
surjective in every degree for $N$ by \cref{pf:steps1:onto}, and hence for $M$ by \cref{pf:steps1:retract}(ii). The dimension
formula follows from \cref{pf:steps1:rank}, $V_n(M)$ being finite-dimensional by \cref{pf:steps1:diaginputs}(iii).
\end{proof}

\subsubsection{Hecke operators and the Hodge--Tate decomposition}

For a tower $H$ over $K$ and $g\in\Gamma$ with $s_g$ finite \'etale, the \emph{Hodge operator} of $g$ is
$\mathrm{tr}_{s_g}\circ t_g^*$ on $\stepsIH^b(Z,\Omega^a)$.

\begin{definition}[Hodge--Tate data for a tower]\label{pf:steps1:def:htcover}
Let $H$ be a tower over $K$ with $Z$ proper, $(g_j)_{j\in J}$ a family of indices and $n\ge0$. \emph{Hodge--Tate
data in degree $n$} consist of a frame $F$ of $(Z,n)$ with a Hodge--Tate matrix $\Phi$, and, for every $j\in J$,
a finite \'etale Galois cover $c\colon W\to Z$ (with finite group $G$) and sections $\sigma_1,\dots,\sigma_r$
splitting $s_{g_j}$ over $c$, a frame $F'$ of $(W,n)$ and a matrix $\Phi'$ such that
$\mathrm{et}(c)\,\Phi=\Phi'\,\mathrm{hdg}(c)$ and $\mathrm{et}(t_{g_j}\sigma_i)\,\Phi=\Phi'\,\mathrm{hdg}(t_{g_j}\sigma_i)$
for all $i$. \formal{PositivityTheory.AutomorphicGaloisHodgeTate.G8.HTCover}
\end{definition}

\begin{lemma}[Traces and splittings]\label{pf:steps1:trsplit}
Let $s\colon Y\to Z$ be a finite \'etale morphism of schemes, $c\colon W\to Z$ a morphism, and $(\sigma_j)_{j\in J}$ a finite
family of morphisms $\sigma_j\colon W\to Y$ with $s\sigma_j=c$. The family \emph{splits $s$ over $c$} if the induced morphism
$\coprod_{j\in J}W\to Y\times_ZW$ is an isomorphism.
\begin{enumerate}[label=\textup{(\roman*)},nosep]
\item \textup{(Algebra)} An \'etale ring map $D\to E$ such that every ring map from $D$ to an algebraically closed field has
  exactly one extension to $E$ is bijective. If $A\to A'$ is \'etale, $C$ is an $A$-algebra and $\tau_j\colon A'\to C$
  ($j\in J$) are $A$-algebra maps such that for all ring maps $x\colon C\to\Omega$, $y\colon A'\to\Omega$ to an
  algebraically closed field which agree on $A$ there is exactly one $j$ with $x\tau_j=y$, then $C\otimes_AA'\to C^J$,
  $c\otimes b\mapsto(c\,\tau_j(b))_j$, is bijective. If moreover $A'$ is finite over $A$ and such a bijection exists, the image
  in $C$ of the trace $\Tr_{A'/A}(b)$ is $\sum_j\tau_j(b)$ for every $b\in A'$.
\item \textup{(Functions)} If $(\sigma_j)$ splits $s$ over $c$, then $c^*(\Tr_s b)=\sum_j\sigma_j^*b$ in $\OO(c^{-1}V)$ for every
  affine open $V\subseteq Z$ and every $b\in\OO(s^{-1}V)$.
\item \textup{(Hodge cohomology)} Let $k$ be a field, $a,b\ge0$, and let $s$, $c$, $\sigma_j$ be morphisms of $k$-schemes. Let
  $\mathrm{tr}_s\colon\stepsIH^b(Y,\Omega^a)\to\stepsIH^b(Z,\Omega^a)$ be the map induced by the trace of differential forms
  along $s$. If $(\sigma_j)$ splits $s$ over $c$, then $c^*\circ\mathrm{tr}_s=\sum_j\sigma_j^*$.
\item \textup{(Splittings)} If $(\sigma_j)$ splits $s$ over $c$, then $(\sigma_jg)$ splits $s$ over $cg$ for every $g\colon W'\to
  W$. For a cartesian square with $s'\colon Y'=Y\times_ZZ'\to Z'$, $v\colon Y'\to Y$ and $g\colon Z'\to Z$, a family
  $(\sigma_j\colon W\to Y')$ over $c\colon W\to Z'$ splits $s'$ over $c$ if and only if $(v\sigma_j)$ splits $s$ over $gc$.
  Consequently, base change along any morphism of base schemes carries finite \'etale morphisms to finite \'etale morphisms
  and splittings to splittings.
\item \textup{(\'Etale cohomology)} Let $\ell$ be a prime and $B$ a scheme with $\ell\in\Gamma(B,\OO_B)^\times$. For schemes and
  morphisms over $B$, pull-back along a finite \'etale Galois cover is injective on each $\stepsIHQ^n(-)(i)$ and on
  $\stepsIHtot$, and every finite \'etale $s$ has a trace map $\mathrm{tr}_s$ on these groups with
  $c^*\circ\mathrm{tr}_s=\sum_j\sigma_j^*$ whenever $(\sigma_j)$ splits $s$ over $c$.
\end{enumerate}
\formallabel{UpperTheory.Gal.G3.algIsoPts_holds}
\formallabel{UpperTheory.Gal.G3.algSplit_holds}
\formallabel{UpperTheory.Gal.G3.algTr_holds}
\formallabel{UpperTheory.Gal.G3.hodgeFun_holds}
\formallabel{UpperTheory.Gal.G3.hodgeTr_holds}
\formallabel{UpperTheory.Gal.G3.traceSplit_holds}
\formallabel{UpperTheory.Gal.crystallineComparison.Br.splitsBasic}
\formallabel{UpperTheory.Gal.crystallineComparison.Br.bcSplits}
\formallabel{UpperTheory.Gal.crystallineComparison.Br.etnGood}
\formallabel{UpperTheory.Gal.crystallineComparison.Br.etGood}
\end{lemma}

\begin{proof}
We use the following results of the development: an \'etale morphism of schemes inducing a bijection on geometric points is
an isomorphism (\lref{UpperTheory.Gal.G3.pIsoPts}); a family splits $s$ over $c$ if and only if for every algebraically
closed field $\Omega$ and all geometric points $x$ of $W$, $y$ of $Y$ with $sy=cx$ there is exactly one $j$ with
$\sigma_jx=y$ (\lref{UpperTheory.Gal.G3.pPtsOfSplits} and \lref{UpperTheory.Gal.G3.pSplitsPts}; the morphism
$\coprod_JW\to Y\times_ZW$ is \'etale, so it is an isomorphism exactly when it is bijective on geometric points).

(i) The first assertion is the case of $\Spec E\to\Spec D$ of the fact just quoted
(\lref{UpperTheory.Gal.G3.pAlgIsoPts}). For the second, $C\otimes_AA'$ and $C^J$ are \'etale $C$-algebras, so the
$C$-algebra map between them is \'etale; a ring map $C\otimes_AA'\to\Omega$ is a pair $(x,y)$ as in the statement, and its
extensions to $C^J$ correspond to the indices $j$ with $x\tau_j=y$, of which there is exactly one; so the first assertion
applies (\lref{UpperTheory.Gal.G3.pAlgSplit}). For the third, $A'$ is finite projective over $A$, the trace commutes
with the base change $A\to C$, and the trace of the $C$-algebra $C^J$ at $(\tau_j(b))_j$ is $\sum_j\tau_j(b)$
(\lref{UpperTheory.Gal.G3.pAlgTr}).

(ii) Let $A=\OO(V)$ and $A'=\OO(s^{-1}V)$, which is finite \'etale over $A$, and let $V'\subseteq c^{-1}V$ be an affine open,
$C=\OO(V')$. The maps $\tau_j=\sigma_j^*\colon A'\to C$ satisfy the hypothesis of (i) by the characterization of splittings
by geometric points, so the restriction of $c^*(\Tr_sb)$ to $V'$ is $\sum_j\sigma_j^*b|_{V'}$. As the affine opens $V'$ cover
$c^{-1}V$, the identity holds in $\OO(c^{-1}V)$ (\lref{UpperTheory.Gal.G3.pHodgeFun}).

(iii) By \lref{UpperTheory.Gal.G3.pHodgeTr} (if the trace of functions along a finite \'etale $s$ is compatible with
splittings in the sense of (ii), then so is the trace on $\stepsIH^b(-,\Omega^a)$: on an affine open a form on $s^{-1}V$ is a
sum of forms $b\,s^*\eta$, whose trace is $\Tr_s(b)\eta$, and the identity passes to cohomology), (iii) follows from (ii).

(iv) By the characterization of splittings by geometric points, the first assertion holds because geometric points of $W'$
map to geometric points of $W$, and the second because geometric points of $Y'$ over a geometric point of $Z'$ are the
geometric points of $Y$ over its image in $Z$ (\lref{UpperTheory.Gal.G3.pSplitsBasic}). Finite \'etale morphisms are
stable under base change, and the base change of a splitting is a splitting by the second assertion
(\lref{UpperTheory.Gal.G3.pBcSplits}).

(v) Finite \'etale morphisms of constant degree $d$ are stable under base change, composites have degree the product, and a
finite \'etale Galois cover with group $G$ has constant degree $\#G$ (\lref{UpperTheory.Gal.G3.pFEtBasic}); the trace
on $\ell$-adic cohomology along finite \'etale morphisms satisfies the splitting identity
(\lref{UpperTheory.Gal.G3.pEtTr}). For a Galois cover $c$ of degree $\#G$, $\mathrm{tr}_c\circ c^*$ is multiplication by
$\#G$, which is invertible in $\Ql$, so $c^*$ is injective; this gives the assertion in each degree and twist
(\lref{UpperTheory.Gal.G3.pEtnGood}), and the assertion for $\stepsIHtot$ follows by taking direct sums of the traces
(\lref{UpperTheory.Gal.G3.pEtGood}).
\end{proof}

\begin{lemma}[Inputs on Hodge and \'etale cohomology]\label{pf:steps1:inputs}
\begin{enumerate}[label=\textup{(\roman*)},nosep]
\item \textup{(Base change of Hodge cohomology)} For a field extension $\phi\colon k\to k'$ and a cartesian square
  $X'=X\times_kk'$, there are $\phi$-semilinear maps $\stepsIH^b(X,\Omega^a)\to\stepsIH^b(X',\Omega^a)$ compatible
  with pull-backs and with traces along finite \'etale maps, carrying bases to bases when $X$ is proper.
\item \textup{(Operators on \'etale cohomology)} Let $H$ be a tower over $K$ with $Z$ proper, $g\in\Gamma$ with $s_g$
  finite \'etale, and $c\colon W\to Z$ a finite \'etale Galois cover with sections $\sigma_i$ splitting $s_g$; assume
  that $(s_g,t_g)_{\Qlbar}$ has some operator. Then for every $n$, $c^*$ is injective on $V_n$, and there is an
  endomorphism $T$ of $V_n(Z)$ which is the degree-$n$ component of $\stepsIhop(g)$ and satisfies
  $c^*\circ T=\sum_i(t_g\sigma_i)^*$.
\end{enumerate}
\formallabel{UpperTheory.Gal.G3.hodgeBC_holds}
\formallabel{UpperTheory.Gal.G3.coverOp_holds}
\end{lemma}

\begin{proof}
(i) is flat base change for the coherent cohomology of the sheaves $\Omega^a$, which commute with base
change \cite[III.9.3]{Hartshorne1977}; compatibility with pull-backs and traces is a routine verification
with \v{C}ech cocycles.
(ii): $c$ is finite \'etale and surjective, so $c^*$ is injective by \cref{pf:steps1:trsplit}(v). The sum
$\sum_i(t_g\sigma_i)^*x$ is invariant under the Galois group of $c$ (which permutes the
splitting sections), hence by Galois descent for the finite \'etale Galois cover $c$ it is of the form $c^*(T'x)$ for a
unique $T'x$; the map $T'$ so defined is an operator of the correspondence $(s_g,t_g)$. By definition $\stepsIhop(g)$ is
also an operator of this correspondence, and any two operators of the same correspondence coincide
(\lref{UpperTheory.Gal.G3.corrUniq_holds}, deferred). The base change of the splitting data from $K$ to $\Qlbar$ is
routine.
\end{proof}

\begin{proposition}[The Hodge--Tate spectral step]\label{pf:steps1:loce}
Let $\iota_{\C}\colon\Qlbar\xrightarrow\sim\C$, let $H$ be a nice tower over $K$, $H'$ its base change to $\C$
along $\iota_{\C}|_K$, $\theta\colon\Gamma\to\Qlbar$ and $a\in\N$. Assume that every $(s_g,t_g)_{\Qlbar}$ has an
operator on \'etale cohomology, that $H$ has Hodge--Tate data in every degree (for $J=\Gamma$), and that
$\iota_{\C}\theta$ occurs in Hodge degree $a$ for $H'$. Then $a\in\mathrm{HT}(\rho)$ for every matrix $\rho$ of the
$\theta$-piece of $H$ for $G_K$.
\formallabel{PositivityTheory.AutomorphicGaloisHodgeTate.G8.locE_of}
\formallabel{PositivityTheory.AutomorphicGaloisHodgeTate.G8.occE}
\formallabel{CrystallineSemistableComparison.Jump.corr_ht}
\end{proposition}

\begin{proof}
Let $x'\in\stepsIH^b(Z_\C,\Omega^a)$ be a nonzero class with $\mathrm{tr}_{s_g}t_g^*x'=\iota_{\C}(\theta(g))x'$ for all $g$, and
put $n=a+b$. Let $F,\Phi$ and, for each $g$, $(W,c,\sigma_i,F',\Phi')$ be Hodge--Tate data in degree $n$, and view
$\Phi$ as the isomorphism $\bigoplus_r\stepsIH^{n-r}(Z,\Omega^r)\otimes_K\stepsICl\to V_n(Z)\otimes\stepsICl$.

\emph{Step 1: intertwining.} Fix $g$, let $T_g$ be the endomorphism of $V_n(Z)$ given by
\cref{pf:steps1:inputs}(ii) and $M_g=\bigoplus_r\mathrm{tr}_{s_g}t_g^*$ on Hodge cohomology. Then
$T_g\circ\Phi=\Phi\circ M_g$. Indeed, by \cref{pf:steps1:inputs}(ii), naturality of the Hodge--Tate data,
\cref{pf:steps1:trsplit}(iii) and again naturality,
\[
 c^*T_g\Phi=\sum_i(t_g\sigma_i)^*\Phi=\sum_i\Phi'(t_g\sigma_i)^*=\Phi'c^*\mathrm{tr}_{s_g}t_g^*=c^*\Phi M_g ,
\]
and $c^*$ is injective on $V_n\otimes\stepsICl$.

\emph{Step 2: the Hodge eigenvector over $\Qlbar$.} By \cref{pf:steps1:inputs}(i) the basis of
$\stepsIH^b(Z,\Omega^a)$ in $F$ is carried to a basis of $\stepsIH^b(Z_\C,\Omega^a)$, compatibly with the Hodge
operators. Let $y'\in\C^{e_a}$ be the coordinates of $x'$; then $(M_g)_ay'=\iota_{\C}(\theta(g))y'$, where $(M_g)_a$
has entries in $K$. Applying $\iota_{\C}^{-1}$ gives $0\ne y\in\Qlbar^{e_a}$ with $(M_g)_ay=\theta(g)y$ for all $g$.
Let $Y\in\bigoplus_r\Qlbar^{e_r}$ be $y$ placed in the summand $r=a$.

\emph{Step 3: the Galois eigenvector.} Put $v=\Phi Y\in V_n(Z)\otimes\stepsICl$. Then $v\ne0$ ($\Phi$ is invertible),
$\stepsIhop(g)v=T_gv=\Phi M_gY=\theta(g)v$ for all $g$ by Step 1, so $v$ lies in the $\theta$-piece tensored with
$\stepsICl$ (joint eigenspaces commute with the flat extension $\Qlbar\to\stepsICl$); and for $\sigma\in G_K$ fixing the
finite set $S$ of coordinates of $y$, the Hodge--Tate property gives
$\sigma(v)=\sigma(\Phi)\,Y=\rho_F(\sigma)^{-1}\Phi\,\mathrm{diag}(\stepsIchar(\sigma)^{-r})Y=\stepsIchar(\sigma)^{-a}\rho_F(\sigma)^{-1}v$,
i.e.\ the semilinear action $v\mapsto\rho_F(\sigma)\sigma(v)$ of $\sigma$ on $V_n(Z)\otimes\stepsICl$ multiplies $v$ by
$\stepsIchar(\sigma)^{-a}$. Let $z\in\stepsICl^{d}$ be the coordinates of $v$ (viewed in $\stepsIHtot\otimes\stepsICl$) in the basis of
the piece defining $\rho$. The action of $\sigma$ on that basis is given by $\rho(\sigma)$; if $\sigma$ also fixes the finitely many
coordinates of the basis vectors with respect to a $\Ql$-basis of the cohomology groups involved, the semilinear action reads
$z\mapsto\rho(\sigma)\sigma(z)$ in these coordinates. Enlarging $S$ by these coordinates we get $\rho(\sigma)\sigma(z)
=\stepsIchar(\sigma)^{-a}z$ for $\sigma$ fixing $S$, with $z\ne0$. Hence $a\in\mathrm{HT}(\rho)$. The last step is a statement of
linear algebra about matrices satisfying the relations of Steps 1--3, which the formal development proves as such.
\end{proof}

\subsubsection{Moving Hodge--Tate weights}

\begin{lemma}[Moving occurrences]\label{pf:steps1:occmove}
Let $a\in\Z$.
\begin{enumerate}[label=\textup{(\roman*)},nosep]
\item If $\rho_1\colon G_K\to M_{d_1}(\Qlbar)$, $\rho_0\colon G_K\to M_{d_0}(\Qlbar)$ and $A\in M_{d_0\times d_1}(\Qlbar)$ is
  injective with $\rho_0(\sigma)A=A\rho_1(\sigma)$, then $a\in\mathrm{HT}(\rho_1)$ implies $a\in\mathrm{HT}(\rho_0)$.
\item If $\rho$ is graded of $\rho_0$ (for the same grading function at every $\sigma$), then $a\in\mathrm{HT}(\rho_0)$
  implies $a\in\mathrm{HT}(\rho)$.
\end{enumerate}
\formallabel{PositivityTheory.AutomorphicGaloisHodgeTate.G8.OccMove_holds}
\end{lemma}

\begin{proof}
(i) If $z$ and $S$ witness $a\in\mathrm{HT}(\rho_1)$, then $Az\ne0$ and, for $\sigma$ fixing $S$ and the entries of $A$,
$\rho_0(\sigma)\sigma(Az)=\rho_0(\sigma)A\sigma(z)=A\rho_1(\sigma)\sigma(z)=\stepsIchar(\sigma)^{-a}Az$.
(ii) Let $c$ be the grading function, $z\ne0$ with $\rho_0(\sigma)\sigma(z)=\stepsIchar(\sigma)^{-a}z$ for $\sigma$ fixing $S$, let
$v=\max\{c(j):z_j\ne0\}$, and let $z'$ be the vector with $z'_j=z_j$ if $c(j)=v$ and $0$ otherwise. For $c(i)=v$,
$(\rho_0(\sigma)\sigma(z))_i=\sum_{c(j)=v}\rho_0(\sigma)_{ij}\sigma(z_j)$, because $z_j=0$ for $c(j)>v$ and
$\rho_0(\sigma)_{ij}=0$ for $c(j)<c(i)$; this equals $(\rho(\sigma)\sigma(z'))_i$. For $c(i)\ne v$ both
$(\rho(\sigma)\sigma(z'))_i$ and $z'_i$ vanish. Hence $\rho(\sigma)\sigma(z')=\stepsIchar(\sigma)^{-a}z'$ with $z'\ne0$.
\end{proof}

\begin{lemma}[Hodge--Tate weights are jumps]\label{pf:steps1:jumps}
Let $P$ be a tied period package, $B_w$ a bridge for $G_w\subseteq G_E$ and $r\colon G_w\to\GL_d(\Qlbar)$.
\begin{enumerate}[label=\textup{(\roman*)},nosep]
\item If $r$ is crystalline for $(P,B_w)$, then for every $\tau\colon E\to\Qlbar$ and every $\tau'\colon K_w\to LL'$
  inducing $\tau$ there is a jump multiset of $r$ at $\tau$.
\item For such $\tau,\tau'$ there are $\gamma\in G_{\Ql}$ and a bridge $B'_w$ for $G_w$ with $j'=\gamma\circ j$,
  $e'(g)=\gamma e(g)\gamma^{-1}$ and $j'|_E=\tau$ (a \emph{conjugate bridge}).
\item If $B'_w$ is conjugate to $B_w$ by $\gamma$, $a\in\mathrm{HT}(r\circ e'^{-1})$ and $J$ is a jump multiset of $r$ at
  $j'|_E$ (for the bridge $B_w$), then $a\in J$.
\end{enumerate}
\formallabel{PositivityTheory.AutomorphicGaloisHodgeTate.G8.jumpsEx_holds}
\formallabel{UpperTheory.Gal.G3.jConjBridge} \formallabel{UpperTheory.Gal.G3.jJumpOfHT2}
\end{lemma}

\begin{proof}
(i) A crystalline representation is de Rham \cite[Exp.~III]{Fontaine1994rss}, and the Hodge filtration on the
$\tau'$-part of $\stepsIDdR$ has a multiset of jumps. (ii) Extend $\tau'$, an embedding of $K_w$ over $\Ql$, to an
automorphism $\gamma$ of $\Qlbar$; since $\tau(x)=\tau'(j(x))$ for $x\in E$, the bridge transported by $\gamma$ has the
required properties. (iii) Transport by $\gamma$ identifies $r\circ e'^{-1}$ with the conjugate of the
$G_{K_w}$-representation, $K_w$ replaced by $\gamma(K_w)$; this does not change the jump multisets (conjugation invariance
of the multiplicities of jumps, \lref{GlobalTheory.Pst.BridgeInd.multV_conj}). For a de Rham representation $V$ there is a basis of
$V\otimes\stepsICl$ on which an open subgroup acts diagonally through powers $\stepsIchar^{-k}$, with the exponents $k$
running through the jump multiset of the relevant embedding with multiplicities
(\lref{GlobalTheory.T.deRham_exists_hodgeTate_eigenbasis}, the Hodge--Tate decomposition of a de Rham
representation \cite{Fontaine1994}). A nonzero vector on which an open subgroup acts through $\stepsIchar^{-a}$ has a
nonzero component in the $\stepsIchar^{-a}$-isotypic part, which is spanned by the basis vectors with $k=a$ (the
characters $\stepsIchar^{-k}$, $k\in\Z$, are distinct on every open subgroup); so $a$ is a jump.
\end{proof}

\subsubsection{The low tower}

\begin{definition}[The low tower]\label{pf:steps1:def:low}
Let $C$ be a good cover of $(H,\theta)$ over $O$ (\cref{pf:steps1:def:cover}). \emph{Low data} $D$ for $C$ choose,
for every $g\in\Gamma$, a finite \'etale correspondence $(s'_g,t'_g)$ of constant degree on $Y$ over $O$ under which
$g$ lies with multiplicity $m_g$. The \emph{low tower} $H_D$ over $k$ has the same index set, $Z_{\mathrm{low}}=Y_k$,
and correspondences the base changes of $(s'_g,t'_g)$ to $k$.
\formal{PositivityTheory.AutomorphicGaloisHodgeTate.G8.LowData}
\formal{PositivityTheory.AutomorphicGaloisHodgeTate.G8.low}
\end{definition}

\begin{lemma}[Properties of the low tower]\label{pf:steps1:low}
Let $C$ be an onto good cover of a nice tower $(H,\theta)$ over $O$ and $D$ low data.
\begin{enumerate}[label=\textup{(\roman*)},nosep]
\item $H_D$ is nice.
\item \textup{(Hodge descent)} Let $\phi\colon k\to F$ and $\psi\colon L\to F$ be maps to a field $F$ of characteristic $0$.
  If $\psi\theta$ occurs in Hodge degree $a$ for $\mathrm{bc}_\phi H$, then $g\mapsto\psi(\theta(g)/m_g)$ occurs in Hodge degree $a$
  for $\mathrm{bc}_\phi H_D$.
\item \textup{(Pieces)} Let $k=E$ be a number field and $L=\Qlbar$. If $\rho_1$ is a matrix of the $(\theta/m)$-piece of $H_D$
  and $\rho_0$ a matrix of the $\theta$-piece of $H$, both for $G_E$, there is an injective $A$ over $\Qlbar$ with
  $\rho_0(\sigma)A=A\rho_1(\sigma)$ for all $\sigma$.
\end{enumerate}
\formallabel{UpperTheory.Gal.G3.jLowNice} \formallabel{UpperTheory.Gal.G3.lowHodge_holds}
\formallabel{UpperTheory.Gal.G3.jLowHodge} \formallabel{UpperTheory.Gal.G3.lowPiece_holds}
\formallabel{UpperTheory.Gal.G3.jLowPiece}
\end{lemma}

\begin{proof}
Let $\pi_k\colon Z\to Y_k$ be the finite \'etale Galois covering of the cover, with group $\Delta$.

(i) $Y_k$ is smooth and projective over $k$ (base change of $Y$), and the base change of a finite \'etale map with a
constant number of geometric lifts has the same property (\lref{UpperTheory.Gal.G3.fEtBasic_holds}).

(ii) We use two facts valid for any cohomology theory with pull-backs and traces along finite \'etale maps and with
coefficients in a field of characteristic $0$, in particular for Hodge cohomology after base change along $\phi$
(\lref{UpperTheory.Gal.G3.coverInv_holds}, \lref{UpperTheory.Gal.G3.liesOp_holds}): $\pi_k^*$ is injective,
and a class on $Z$ fixed by $\mathrm{tr}_{s_g}t_g^*$ for every $g$ with $s_g$ an isomorphism and $\theta(g)=1$ (the deck
transformations of the cover) is a pull-back along $\pi_k$ (Galois descent); and if $g$ lies under $(s',t')$ with
multiplicity $m$, then $\mathrm{tr}_{s_g}t_g^*\pi_k^*=m\,\pi_k^*\,\mathrm{tr}_{s'}t'^*$ (a projection formula for the finite
\'etale map $Y_g\to Z\times_YT$ of degree $m$). Let $x\ne0$ be a $\psi\theta$-eigenclass in some $\stepsIH^b(Z_F,\Omega^a)$.
For a deck transformation $g$, $\mathrm{tr}_{s_g}t_g^*x=\psi(\theta(g))x=x$, so $x=\pi_k^*z$ with $z\ne0$. For every $g$,
$m_g\pi_k^*(\mathrm{tr}_{s'_g}t'^*_gz)=\mathrm{tr}_{s_g}t_g^*x=\psi(\theta(g))\pi_k^*z$, and since $m_g\ne0$ in $F$ and $\pi_k^*$ is
injective, $\mathrm{tr}_{s'_g}t'^*_gz=\psi(\theta(g)/m_g)z$.

(iii) Let $\Pi=\pi_k^*\colon\stepsIHtot(Y_{\Omega})\to\stepsIHtot(Z_{\Omega})$. On \'etale cohomology, $\stepsIhop_H(g)\circ\Pi
=m_g\,\Pi\circ\stepsIhop_{H_D}(g)$ (\lref{UpperTheory.Gal.G3.hopLies_holds}, the \'etale analogue of the projection
formula in (ii)), and $\Pi\otimes\Qlbar$ is injective (\lref{UpperTheory.Gal.G3.pieceLe_holds}). Hence $\Pi$ maps the
$(\theta/m)$-piece of $H_D$ injectively into the $\theta$-piece of $H$, and it commutes with the Galois actions (pull-back
along a morphism defined over $E$). Expressing the images of the basis vectors defining $\rho_1$ in the basis defining
$\rho_0$ gives $A$.
\end{proof}

\begin{remark}
In (iii) it is the $(\theta/m)$-piece of the \emph{low} tower that embeds into the $\theta$-piece of $H$. This is the direction
needed below: Hodge--Tate occurrences are established for the low tower, which has a smooth projective model, and moved to
$H$ by \cref{pf:steps1:occmove}(i).
\end{remark}

\begin{lemma}[Hodge--Tate data for the local low tower]\label{pf:steps1:coverlow}
Assume the conclusion of \cref{pf:steps1:htspec}. Let $E$ be a number field, $w$ a finite place, $C$ a cover of
$(H,\theta)$ over $\stepsIOw$ with low data $D$, and $\tau_K\colon E\to K$ an embedding into a finite extension $K$ of $\Ql$
with $|\tau_K(x)|\le1$ for $x\in\stepsIOw$. Then $\mathrm{bc}_{\tau_K}H_D$ has Hodge--Tate data in every degree.
\formallabel{UpperTheory.Gal.G3.coverLow_of_cov}
\formallabel{PositivityTheory.AutomorphicGaloisHodgeTate.G8.coverOf}
\formallabel{UpperTheory.Gal.G3.lowCov_holds}
\end{lemma}

\begin{proof}
The map $\tau_K$ extends to $\stepsIOw\to\OO_K$. The base change $M_Z=Y\otimes_{\stepsIOw}\OO_K$ is a smooth projective model of
$Z_{\mathrm{low}}\otimes K$, hence a smooth good model (\lref{UpperTheory.Gal.Seam.goodOf}), of pure relative dimension.
For each $g$, the map $s'_g$ is finite \'etale over all of $Y$, so it has a Galois closure: a finite \'etale Galois cover
$W\to Y$ with sections splitting $s'_g$ (\lref{UpperTheory.Gal.G3.galClosD_holds}); its base change $M_W$ to $\OO_K$ is again
a smooth good model, and the cover, the sections and the compositions $t'_g\sigma_i$ are morphisms of models. The base change
of all these data to $K$ consists of Hodge--Tate cover data for the tower $\mathrm{bc}_{\tau_K}H_D$, defined over models
(\lref{UpperTheory.Gal.G3.jLowCov}, deferred for the routine verifications on base change of Galois covers). Now choose by
\cref{pf:steps1:htspec}(i) frames $F_Z,F_W$ with $\sum e_a=d$ and put $\Phi=\Phi(M_Z,F_Z)$, $\Phi'=\Phi(M_W,F_W)$; these are
Hodge--Tate, and the relations required in \cref{pf:steps1:def:htcover} are instances of the naturality in
\cref{pf:steps1:htspec}(ii) for the morphisms of models $M_W\to M_Z$ inducing $c$ and $t'_g\sigma_i$.
\end{proof}

\begin{lemma}[Base change of towers and of matrices]\label{pf:steps1:localmat}
Let $H$ be a nice tower over a number field $E$, $\theta\colon\Gamma\to\Qlbar$, and $B_w$ a bridge, with $\tau_K\colon E\to
K_w$ the embedding induced by $j$.
\begin{enumerate}[label=\textup{(\roman*)},nosep]
\item There is a matrix of the $\theta$-piece of $H$ for $G_E$.
\item $\mathrm{bc}_{\tau_K}H$ is nice, every one of its correspondences has an operator on \'etale cohomology over $\Qlbar$, and
  for every matrix $\rho$ of the $\theta$-piece of $H$, $\sigma\mapsto\rho(e^{-1}(\sigma))$ is the matrix (in the sense of
  \cref{pf:steps1:def:tower}, without the homomorphism property) of the $\theta$-piece of $\mathrm{bc}_{\tau_K}H$ for $G_{K_w}$.
\item For a further field map $K\to F$ the base change of $\mathrm{bc}_{\tau_K}H$ to $F$ is (isomorphic to) the base change of $H$
  along the composite, compatibly with the correspondences.
\item If $B_w$ induces a place over $w$, then $|\gamma(j(x))|\le1$ for $x\in\stepsIOw$ and $\gamma\in G_{\Ql}$.
\end{enumerate}
\formallabel{UpperTheory.Gal.G3.matEx_of} \formallabel{UpperTheory.Gal.G3.jLocalMat}
\formallabel{UpperTheory.Gal.G3.jBCTower} \formallabel{UpperTheory.Gal.G3.jBridgeInt}
\end{lemma}

\begin{proof}
(i) follows from the existence of semisimple representations on pieces,
\lref{UpperTheory.Gal.Int.Geo.repEx_holds} (for every nice tower there are $d$ and a continuous semisimple $\rho$ attached
to the $\theta$-piece; its underlying matrix of the piece is the required matrix). (ii) Niceness is preserved by base change
(\lref{UpperTheory.Gal.Seam.niceBc}); operators exist by \lref{UpperTheory.Gal.Seam.hopEx_holds} (for nice towers
over a field and an algebraically closed field of characteristic $0$, $\stepsIhop(g)$ is an operator of $(s_g,t_g)$); and the
statement on matrices is the local tower lemma \lref{UpperTheory.Gal.Seam.locTower_holds} applied to $j\colon\Omega\to\Qlbar$,
composed with $e$. (iii) is the transitivity of fibre products. (iv) Elements of $\stepsIOw$ have $|j(x)|\le1$ because $B_w$
induces a place over $w$, and $G_{\Ql}$ preserves absolute values on $\Qlbar$.
\end{proof}

\subsubsection{Proof of $\sfname{JumpStep}$}

\begin{theorem}[The core of $\sfname{JumpStep}$]\label{pf:steps1:jumpcore}
Let $E$, $\Omega$, $\ell$, $\iota$ be as in \cref{sec:pf:steps1:conv:loc}, $H$ a nice tower over $E$, $\theta\colon\Gamma\to\Qlbar$,
$\rho_0$ a matrix of the $\theta$-piece and $\rho$ graded of $\rho_0$, and $w$ a finite place such that $(H,\theta)$ has an onto good
cover over $\stepsIOw$. Let $P$ be a tied period package, $A$ a place over $w$, $B_w$ an inducing bridge, and assume that
$\rho|_{D_A}$ is crystalline for $(P,B_w)$. Then for every $\tau\colon E\to\Qlbar$ and $\tau'$ inducing $\tau$ there is a jump
multiset $J$ of $\rho|_{D_A}$ at $\tau$ containing every $b\in\N$ such that $\iota\theta$ occurs in Hodge degree $b$ for
$\mathrm{bc}_{\iota\tau}H$. \formallabel{PositivityTheory.AutomorphicGaloisHodgeTate.G8.jumpCore_of}
\formallabel{PositivityTheory.AutomorphicGaloisHodgeTate.G8.occ_low}
\end{theorem}

\begin{proof}
By \cref{pf:steps1:jumps}(i) there is a jump multiset $J$ of $\rho|_{D_A}$ at $\tau$; we show that it contains every such $b$.
Fix $b$. By \cref{pf:steps1:jumps}(ii) choose a bridge $B'_w$ conjugate to $B_w$ by some $\gamma$ with $j'|_E=\tau$, and let
$\tau_K\colon E\to K'=K'_w$ be the embedding induced by $j'$; then $\iota\tau=\iota\circ\tau_K$ (viewing $K'\subset\Qlbar$).
Let $C$ be the onto good cover, $D$ low data for $C$ and $H_D$ the low tower; it is nice (\cref{pf:steps1:low}(i)), and by
\cref{pf:steps1:localmat}(i) it has a matrix $\rho_1$ of its $(\theta/m)$-piece for $G_E$.

By \cref{pf:steps1:low}(ii) (with $\phi=\iota\tau$, $\psi=\iota$, $F=\C$), $\iota(\theta/m)$ occurs in Hodge degree $b$ for
$\mathrm{bc}_{\iota\tau}H_D=\mathrm{bc}_{\iota}(\mathrm{bc}_{\tau_K}H_D)$. We apply \cref{pf:steps1:loce} to the nice tower
$\mathrm{bc}_{\tau_K}H_D$ over $K'$ (\cref{pf:steps1:localmat}(ii),(iii)) with $\iota_{\C}=\iota$: its correspondences have
operators (\cref{pf:steps1:localmat}(ii)); it has Hodge--Tate data in every degree by \cref{pf:steps1:coverlow}, which applies
because $|\tau_K(x)|=|\gamma(j(x))|\le1$ for $x\in\stepsIOw$ (\cref{pf:steps1:localmat}(iv)); and the matrix
$\sigma\mapsto\rho_1(e'^{-1}(\sigma))$ of its $(\theta/m)$-piece is provided by \cref{pf:steps1:localmat}(ii). Hence
$b\in\mathrm{HT}(\rho_1\circ e'^{-1})$.

By \cref{pf:steps1:low}(iii) there is an injective $A$ with $\rho_0(\sigma)A=A\rho_1(\sigma)$, so $b\in\mathrm{HT}(\rho_0\circ e'^{-1})$
by \cref{pf:steps1:occmove}(i), and $b\in\mathrm{HT}(\rho\circ e'^{-1})$ by \cref{pf:steps1:occmove}(ii), the grading being preserved by
restriction. Finally \cref{pf:steps1:jumps}(iii) (with $r=\rho|_{D_A}$ and the jump multiset $J$ at $j'|_E=\tau$) gives $b\in J$.
\end{proof}

\begin{theorem}[$\sfname{JumpStep}$]\label{pf:steps1:jump}
Let $E$, $\Omega$, $\ell$, $\iota$ be as in \cref{sec:pf:steps1:conv:loc}, $H$ a nice tower over $E$, $\theta\colon\Gamma\to\Qlbar$,
$\rho\colon G_E\to\GL_d(\Qlbar)$ continuous and attached to the $\theta$-piece, and $w\mid\ell$ a finite place at which $(H,\theta)$
has an onto good cover over $\stepsIOw$. Then $\rho$ satisfies the jump condition at $w$ for the Hodge occurrences of $\iota\theta$:
for every tied package $P$, place $A$ over $w$, inducing bridge $B_w$, $\tau\colon E\to\Qlbar$ and $\tau'$ inducing $\tau$, there is a
jump multiset $J$ of $\rho|_{D_A}$ at $\tau$ containing every $b\ge0$ such that $\iota\theta$ occurs in Hodge degree $b$ for
$\mathrm{bc}_{\iota\tau}H$. \formallabel{UpperTheory.Gal.Int.Geo.jumpStep_holds}
\formallabel{UpperTheory.Gal.G3.jumpStep_of}
\formallabel{PositivityTheory.AutomorphicGaloisHodgeTate.G8.jumpStep_of}
\end{theorem}

\begin{proof}
Let $\rho_0$ be a matrix of the piece of which $\rho$ is graded. By $\sfname{CrysStep}$ (\cref{pf:steps1:crys}), applied with the
onto cover, $\rho|_{D_A}$ is crystalline for $(P,B_w)$; \cref{pf:steps1:jumpcore} then gives the multiset $J$. The inputs of
\cref{pf:steps1:jumpcore} are \cref{pf:steps1:htspec} (in \cref{pf:steps1:coverlow}) and the existence of matrices of pieces
(\cref{pf:steps1:localmat}(i)).
\end{proof}

\begin{remark}
$\sfname{JumpStep}$ asserts membership of the Hodge degrees in a jump multiset, not equality of multiplicities, and it concerns all
period packages tied to the concrete period rings. As in \cref{pf:steps1:crys}, neither the continuity of $\rho$ nor the hypothesis
$w\mid\ell$ is used in the proof; at places $w\nmid\ell$ no bridge induces a place over $w$, so the conclusion is vacuous there.
\end{remark}

% ===========================================================================
\subsection{The crystalline Frobenius step}\label{sec:pf:steps1:phi}

$\sfname{PhiStep}$ is the analogue above $\ell$ of the Eichler--Shimura relation: when the integral model of the tower
carries correspondences $T_0,\dots,T_n$ satisfying the pointwise congruence relation, the linearized crystalline
Frobenius $\Phi=\varphi^f$ on $\stepsIDcris$ of the cohomology of the model satisfies
$\sum_j(-1)^jq^{\binom j2}T_j\Phi^{n-j}=0$ (\cref{pf:steps1:finall}); on the image of the Hecke eigen-piece, where
$T_j$ acts by the scalar $\vartheta_j$, this says that the Hecke polynomial $\stepsIHk_{q,n,\vartheta}$ kills $\Phi$, and
the Weil--Deligne representation, on which geometric Frobenius acts through $\Phi$, has its eigenvalues among the
roots. The classical counterpart is the compatibility of the crystalline comparison isomorphism with Frobenius and
with correspondences \cite{Faltings1989,FaltingsChai1990}; the formal proof establishes the relation directly from the
comparison maps of the development.

\begin{definition}[Invariants, annihilation and nilpotence]\label{pf:steps1:def:annd}
Let $G$ be a group, $B$ a $\Ql$-vector space with $\Ql$-linear operators $\sigma_g$ ($g\in G$) and $\psi$, and $\rho$ a
representation of $G$ on a vector space $V$ over a field $L\supseteq\Ql$. Put $D(\rho)=(V\otimes_{\Ql}B)^G$, the
invariants for $\rho(g)\otimes\sigma_g$. For $q,n\in\N$ and endomorphisms $T_j$ of $V$, the \emph{annihilation
relation} $\mathrm{Ann}(q,n,T)$ holds if $\sum_{j=0}^n(-1)^jq^{\binom j2}(T_j\otimes\psi^{n-j})(x)=0$ for all $x\in
D(\rho)$; for a polynomial $Q\in L[X]$, $Q$ is \emph{nilpotent on $D(\rho)$} if $Q(1\otimes\psi)^N$ vanishes on $D(\rho)$
for some $N$. For a period package $P$, $K=K_w$ and $f\ge1$ we use these notions with $B=P.B$, the Galois action of
$P$ and $\psi=\varphi^f$, and with Fontaine's semistable period ring $\stepsIBst$ of $P$ and $\psi=\varphi^f$.
\formal{UpperTheory.Gal.crystallineComparison.Dinv} \formal{UpperTheory.Gal.crystallineComparison.AnnD}
\formal{UpperTheory.Gal.crystallineComparison.NilD}
\end{definition}

For a crystalline representation and $B=P.B\cong\stepsIBcris$, $D(\rho)=\stepsIDcris(V)$ and $1\otimes\varphi^f$ is the
linearized crystalline Frobenius, $f$ being the residue degree.

\subsubsection{The crystalline congruence relation}

\begin{definition}[Operators of correspondences; period ring presentations]\label{pf:steps1:def:ishk}
Let $M$ be a smooth good model over $\OO_K$, $T=(\alpha,\beta)$ an $\ell$-correspondence on $M.X$ with $\beta$ and
$\alpha$ over the same point of $\Spec\OO_K$, and $j\ge0$. An endomorphism $A$ of $V_j(M)$ is \emph{the operator of
$T$} if for every model $W$ with $w\colon W.X\to M.X$ over $\OO_K$ and every $d_T$-tuple of lifts $z_i\colon W.X\to
T_Z$ of $w$ along $\alpha$ which split $\alpha$ over every scheme on which $\ell$ is invertible, one has
$w^*\circ A=\sum_i(\beta z_i)^*$ on $V_j$. A \emph{presentation of $\stepsIBcris$} is a $\Ql$-algebra $C$ with a ring map
$c\colon\stepsIAcris\to C$ such that $c(t)$ is a unit, every $x\in C$ has $c(t)^rx\in c(\stepsIAcris)$ for some $r$, and
$c(a)=0$ iff $t^ra=0$ for some $r$, together with a Frobenius and a $G_K$-action on $C$ compatible with $c$.
\formal{UpperTheory.Ru.crystallineComparison.IsHk} \formal{UpperTheory.Ru.crystallineComparison.PB}
\end{definition}

The proof of the congruence relation compares the Frobenius of the period ring with the Frobenius of the special fibre, acting
on points with values in an auxiliary class of algebras. Fix a compatible system $\ell^\flat$ of $\ell$-power roots of $\ell$ in
$\OO_{\stepsICl}$, viewed in the tilt $\OO^\flat$ of $\OO_{\stepsICl}$.

\begin{definition}[Test algebras and point functors]\label{pf:steps1:def:test}
A \emph{test algebra} is a perfect ring $A$ of characteristic $\ell$ with a ring map $\OO^\flat\to A$ such that $A$ is complete and
torsion free for $\pi=$ the image of $\ell^\flat$, absolutely integrally closed, and such that for $f,g\in A$ with $\pi^N\in(f,g)$ there
is an idempotent $\epsilon$ with $\epsilon f\in(g)$ and $(1-\epsilon)g\in(f)$. Put $A^\sharp=W(A)/(\xi)$, an $\OO_{\stepsICl}$-algebra
($\xi$ the standard generator of the kernel of $\theta$), and $\bar A=A/\pi$. For a scheme $X$ over a ring $O$ and ring maps
$\kappa\colon O\to\OO_{\stepsICl}$ and $g\colon O\to\OO^\flat/\pi$, the \emph{point functors} $\mathrm{PX}(X,\kappa)$ and
$\mathrm{PQ}(X,g)$ send $A$ to the sets of morphisms $\Spec A^\sharp\to X$ over $\kappa$, resp.\ $\Spec\bar A\to X$ over $g$. A
transformation $F$ of $\mathrm{PQ}(X,g)$ is a \emph{Frobenius of level $f$} if $F(y)=y\circ\Spec(x\mapsto x^{\ell^f})$, and
$\rho\colon\mathrm{PX}(X,\kappa)\to\mathrm{PQ}(X,g)$ is a \emph{reduction of level $m$} if $\rho(y)=y\circ\Spec\psi$ for every map
$\psi\colon A^\sharp\to\bar A$ with $\psi(\theta(x))=x_0^{\ell^m}$. A finite family of maps $A\to B_s$ of test algebras is a
\emph{decomposition} if $A\to\prod_sB_s$ is bijective. For functors $Q$, $P$ from test algebras to sets, transformations
$u_i\colon Q\to P$ and integers $a_i$, we say that $\sum_ia_i[u_i]=0$ \emph{locally} if for every test algebra $A$ and $y\in Q(A)$
there is a decomposition $(A\to B_s)_s$ such that for every $s$ either $B_s=0$ or $\sum_ia_i[u_i(y|_{B_s})]=0$ in the free abelian
group on $P(B_s)$. The development attaches to every such functor $P$ cohomology groups $\mathrm{MH}^j_\nu(P)$ (degree $j$, level
$\nu\ge1$) with coefficients in the level-$\nu$ crystalline period rings $A\mapsto A_{\mathrm{cris},\nu}$ of the test algebras,
with pull-backs along transformations and a Frobenius
$\varphi_H$ commuting with them; for a model $M$, the \'etale comparison family $e$ of the development
(\lref{CrystallineSemistableComparison.E.etMap}) maps $T_j(M)$ to $\mathrm{MH}^j_\nu(\mathrm{PX}(M.X))$, compatibly with
Frobenius, pull-backs and the Galois action, and for $v_k\in T_j(M)$, $c_k\in\stepsIAcris$ we write
$\mathrm{ev}_\nu(v,c)=\sum_kc_k\,e_\nu(v_k)$ and $\mathrm{vec}_r(v,c)=\sum_kv_k\otimes c(c_k)c(t)^{-r}\in V_j(M)\otimes C$.
\formal{CrystallineSemistableComparison.PAlg} \formal{CrystallineSemistableComparison.PAlg.IsTest}
\formal{CrystallineSemistableComparison.PX} \formal{UpperTheory.Ru.crystallineComparison.PQ}
\formal{UpperTheory.Ru.crystallineComparison.IsFr} \formal{UpperTheory.Ru.crystallineComparison.IsRd}
\formal{UpperTheory.Ru.crystallineComparison.TZero} \formal{UpperTheory.Ru.crystallineComparison.evE}
\formal{UpperTheory.Ru.crystallineComparison.vec}
\end{definition}

\begin{lemma}[Lattice expansions]\label{pf:steps1:latexp}
Let $M$ be a model, $j\ge0$ and $(C,\varphi)$ a presentation of $\stepsIBcris$.
\begin{enumerate}[label=\textup{(\roman*)},nosep]
\item Every $x\in V_j(M)\otimes C$ equals $\mathrm{vec}_r(v,c)$ for some finite families $v_k\in T_j(M)$, $c_k\in\stepsIAcris$ and
  some $r\ge0$.
\item Let $T_0,\dots,T_n$ be $\ell$-correspondences on $M.X$ with operators $A_i$ on $V_j(M)$, $\pi_W\colon W\to M.X$ a morphism of
  models and $s_{i,k}$ ($k<d_{T_i}$) lifts of $\pi_W$ along $\alpha_{T_i}$ splitting $\alpha_{T_i}$ over every scheme on which $\ell$
  is invertible, and $v_k\in T_j(M)$. Then there are $s,b\ge0$ and $w_{i,k}\in T_j(M)$ with $w_{i,k}=\ell^sA_i(v_k)$ in $V_j(M)$ and
  $\ell^b\pi_W^*w_{i,k}=\ell^{s+b}\sum_{k'}(\beta_{T_i}s_{i,k'})^*v_k$ in $T_j(W)$.
\item In the situation of \textup{(ii)}, let $\kappa\ne0$ be an integer, $\sigma_i\in\Z$, $c_k\in\stepsIAcris$, $r,f,N\ge0$. If
  $\mathrm{vec}_0\bigl((w_{i,k})_{i,k},(\kappa\sigma_i\ell^{fir}\varphi^{f(n-i+N)}(c_k))_{i,k}\bigr)=0$, then
  $\sum_i\sigma_i(A_i\otimes\varphi^{f(n-i+N)})(\mathrm{vec}_r(v,c))=0$.
\end{enumerate}
\formallabel{UpperTheory.exists_crystalline_lattice_expansion}
\formallabel{UpperTheory.hecke_correspondence_product}
\formallabel{UpperTheory.crystalline_expansion_relation_of_scaled_endomorphisms}
\end{lemma}

\begin{proof}
(i) $V_j(M)=T_j(M)\otimes\Q$, and every element of $C$ has the form $c(a)c(t)^{-r}$ (because $c(t)^ry\in c(\stepsIAcris)$ for some
$r$); so every pure tensor, and every finite sum of pure tensors after passing to a common $r$, is of the stated form.
(ii) Since $V_j(M)=T_j(M)\otimes\Q$, there is $s$ with $\ell^sA_i(v_k)$ integral for all $i,k$; let $w_{i,k}$ be integral lifts.
By the definition of the operator, $\pi_W^*A_i(v_k)=\sum_{k'}(\beta_{T_i}s_{i,k'})^*v_k$ in $V_j(W)$, so
$\pi_W^*w_{i,k}-\ell^s\sum_{k'}(\beta_{T_i}s_{i,k'})^*v_k$ is $\ell$-power torsion in $T_j(W)$, hence killed by some $\ell^b$, uniformly in the
finitely many $(i,k)$.
(iii) Since $\varphi(t)=\ell t$, $(A_i\otimes\varphi^{f(n-i+N)})(v_k\otimes c(c_k)c(t)^{-r})=\ell^{-f(n-i+N)r}A_i(v_k)\otimes
c(\varphi^{f(n-i+N)}c_k)c(t)^{-r}$. Hence the left side of the conclusion equals
$(\kappa\ell^{s+f(n+N)r})^{-1}c(t)^{-r}\,\mathrm{vec}_0((w_{i,k}),(\kappa\sigma_i\ell^{fir}\varphi^{f(n-i+N)}(c_k)))=0$.
\end{proof}

\begin{lemma}[The congruence relation on test points]\label{pf:steps1:wzero}
Let $O$ be a discrete valuation ring with uniformizer $\varpi$ in which $\ell\ne0$ is not a unit, $Y$ a smooth quasi-compact
$O$-scheme covered by finitely many affine opens with affine pairwise intersections, $f>0$, and $T_0,\dots,T_n$ $\ell$-correspondences
on $Y$ over $O$ satisfying the pointwise congruence relation for $(\ell^f,n,\varpi)$. Let $\kappa\colon O\to\OO_{\stepsICl}$, $g\colon
O\to\OO^\flat/\pi$ with $g(\varpi)=0$ and $g(x)^{\ell^f}=g(x)$ for all $x$, $F$ a Frobenius of level $f$ on $\mathrm{PQ}(Y,g)$, $\rho$ a
reduction of some level $m$, and $\pi_W\colon W\to Y$ with sections $s_{j,i}$ of $\alpha_{T_j}$ over $\pi_W$ splitting every
$\alpha_{T_j}$ universally over schemes on which $\ell$ is invertible. Then there is $N$ such that, with $\sigma_j=(-1)^j(\ell^f)^{\binom j2}$,
\[
 \sum_{j,i}\sigma_j\,\bigl[\beta_{T_j}s_{j,i}\circ\rho\circ F^{n-j+N}\bigr]=0\quad\text{locally}
\]
as transformations $\mathrm{PX}(W,\kappa)\to\mathrm{PQ}(Y,g)$.
\formallabel{UpperTheory.Ru.crystallineComparison.wZero}
\end{lemma}

\begin{proof}
Let $\mu=\sum_{j\text{ even}}(\ell^f)^{\binom j2}d_{T_j}$, the total multiplicity of either side of the pointwise relation, and choose
$N$ with $(\mu+1)^\mu\le\ell^{fN}$. Let $A$ be a test algebra and $y\in\mathrm{PX}(W,\kappa)(A)$.

\emph{Charts.} Choose finitely many pairs of affine opens $U_\alpha,V_\alpha$ of $Y$ with $\alpha_{T_j}^{-1}U_\alpha\subseteq
\beta_{T_j}^{-1}V_\alpha$ for all $j$, the $U_\alpha$ covering $Y$. \'Etale surjective maps to $\Spec A^\sharp$ have sections and
idempotents of $A^\sharp$ lift uniquely to $A$ (\lref{UpperTheory.Ru.crystallineComparison.testFacts}, deferred), so there is a
decomposition of $A$ after which the image of $\pi_Wy$ on each factor lies in one $U_\alpha$.

\emph{Norms.} Fix a factor, mapping into $U=U_\alpha$, and generators $x_1,\dots,x_r$ of $\Gamma(V_\alpha)$ over $O$. For each $j$, $i$ let
$p_{j,i}$ be the point $\beta_{T_j}s_{j,i}\circ\rho(y)$ composed with the $\ell^{f(n-j)}$-power Frobenius, i.e.\ the tuple of the
$\ell^{f(n-j)}$-th powers of the images of the $x_k$ in $\bar A$. The pointwise congruence relation, which concerns points over
valuation rings at generic points of the special fibre, implies the \emph{norm identity}: the products of the generic linear forms
$X_0-\sum_kp_k X_k$ over the even multiset $\{p_{j,i}\}$ ($j$ even, multiplicity $(\ell^f)^{\binom j2}$) and over the odd multiset
coincide, for points with values in any ring in which $\ell$ is a nonzerodivisor, reduced modulo $\varpi$
(\lref{UpperTheory.Ru.crystallineComparison.normOfPt'}, deferred; it is a specialization argument on charts).

\emph{Equality after Frobenius.} If two multisets of $\mu$ points with values in $\bar A$ have the same norm form, then after
raising all coordinates to the power $\ell^{fN}$, $(\mu+1)^\mu\le\ell^{fN}$, the two multisets coincide on each factor of a
decomposition of $A$ (\lref{UpperTheory.Ru.crystallineComparison.wzLocal}, deferred; the test-algebra axioms provide the
idempotents separating the factors on which the coordinates of the points agree). Since $F^{N}$ raises coordinates to the power
$\ell^{fN}$ and points of the affine $V_\alpha$ are determined by the values of the $x_k$, the multisets of points
$\beta_{T_j}s_{j,i}\circ\rho\circ F^{n-j+N}(y)$, $j$ even, and $j$ odd, counted with multiplicities $(\ell^f)^{\binom j2}$, coincide on
each factor, which is the asserted identity in the free abelian group. Decompositions of decompositions are decompositions
(\lref{UpperTheory.Ru.crystallineComparison.partsEx}, deferred), so the pieces combine.
\end{proof}

\begin{theorem}[The congruence relation from two comparison properties]\label{pf:steps1:top}
Let $K\subset\Qlbar$ be finite over $\Ql$, $\varpi$ a uniformizer of $\OO_K$, $f>0$, $M$ a model with $M.X$ smooth and quasi-compact
over $\OO_K$ and covered by finitely many affine opens with affine pairwise intersections, $n,j\ge0$, $T_0,\dots,T_n$
$\ell$-correspondences on $M.X$ over $\OO_K$ satisfying the pointwise congruence relation for $(\ell^f,n,\varpi)$, $A_i$ the operator
of $T_i$ on $V_j(M)$, and $(C,\varphi)$ a presentation of $\stepsIBcris$. Assume:
\begin{enumerate}[label=\textup{(\alph*)},nosep]
\item \textup{(Injectivity)} if $t^u\,\mathrm{ev}_\nu(v,c)=0$ for all levels $\nu$, then $\mathrm{vec}_0(v,c)=0$;
\item \textup{(Absolute invariance)} there are $m,N_1\ge0$, $g\colon\OO_K\to\OO^\flat/\pi$ with $g(\varpi)=0$ and $g(x)^{\ell^f}=g(x)$, a
  Frobenius $F$ of level $f$ on $\mathrm{PQ}(M.X,g)$ and a reduction $\rho$ of level $m$, such that for every $G_K$-invariant
  $\mathrm{vec}_r(v,c)$ there are $a,u\ge0$ and classes $X_\nu\in\mathrm{MH}^j_\nu(\mathrm{PQ}(M.X,g))$ with $\varphi_H^f(X_\nu)=F^*X_\nu$
  and $\ell^at^u\,\mathrm{ev}_\nu(v,\varphi^{fN_1}(c))=t^{u+r}\rho^*X_\nu$ for all $\nu$.
\end{enumerate}
Then there is $N$ with $\sum_{i=0}^n(-1)^i(\ell^f)^{\binom i2}(A_i\otimes\varphi^{f(n-i+N)})(x)=0$ for every $G_K$-invariant
$x\in V_j(M)\otimes C$.
\formallabel{UpperTheory.Ru.crystallineComparison.top}
\formallabel{UpperTheory.Ru.crystallineComparison.levZero}
\end{theorem}

\begin{proof}
Let $m,N_1,g,F,\rho$ be as in (b) and $\sigma_i=(-1)^i(\ell^f)^{\binom i2}$. There is a model $W$ with $\pi_W\colon W.X\to M.X$ and
sections $s_{i,k}$ of $\alpha_{T_i}$ over $\pi_W$ splitting all $\alpha_{T_i}$ universally over schemes on which $\ell$ is invertible
(\lref{UpperTheory.Ru.crystallineComparison.wEx}, deferred; $W$ is a scheme of ordered splittings). By \cref{pf:steps1:wzero}
there is $N_0$ with $\sum_{i,k}\sigma_i[\beta_{T_i}s_{i,k}\circ\rho\circ F^{n-i+N_0}]=0$ locally. Put $N=N_0+N_1$. We use three facts
of the development. (1) If $\sum_ia_i[u_i]=0$ locally, then $\sum_ia_iu_i^*=0$ on $\mathrm{MH}$; together with $\varphi_H^f=F^*$ on the
classes $X_\nu$ and the compatibility of $\rho^*$ with $\varphi_H$, this gives
$\sum_{i,k}\sigma_i(\beta_{T_i}s_{i,k})^*\varphi_H^{f(n-i+N_0)}\rho^*X_\nu=0$ (\lref{UpperTheory.Ru.crystallineComparison.siteAnn},
deferred). (2) The permutations of the sheets $s_{i,\cdot}$ induce transformations of $\mathrm{PX}(W)$ over $\mathrm{PX}(M.X)$ acting,
locally, transitively on the fibres (\lref{UpperTheory.Ru.crystallineComparison.wTrans}, deferred), so a class on
$\mathrm{PX}(M.X)$ whose pull-back to $\mathrm{PX}(W)$ vanishes is killed by $C_0=\prod_i(d_{T_i})!$
(\lref{UpperTheory.Ru.crystallineComparison.avgMH}, deferred). (3) The classes $e_\nu(v_k)$ are fixed by $\varphi_H$, which is
semilinear with respect to $\varphi$ on the coefficients, so $\mathrm{ev}_\nu(v,\varphi(c))=\varphi_H(\mathrm{ev}_\nu(v,c))$; and
$\varphi(t)=\ell t$.

Let $x$ be invariant. By \cref{pf:steps1:latexp}(i), $x=\mathrm{vec}_r(v,c)$; let $a,u,X_\nu$ be as in (b), and $s,b,w_{i,k}$ as in
\cref{pf:steps1:latexp}(ii) for $W$ and the $s_{i,k}$. Applying $\varphi_H^{f(n-i+N_0)}$ to the identity in (b) and using (3) gives
\[
 \ell^{a+f(n-i+N_0)u}\,t^u\,\mathrm{ev}_\nu\bigl(v,\varphi^{f(n-i+N)}(c)\bigr)=\ell^{f(n-i+N_0)(u+r)}\,t^{u+r}\,
 \varphi_H^{f(n-i+N_0)}\rho^*X_\nu .
\]
Multiply by $\sigma_i\ell^{fi(u+r)}$, which makes the power of $\ell$ on the right side equal to $\ell^{f(n+N_0)(u+r)}$ for all $i$ and
the scalar on the left side equal to $\ell^{a+f(n+N_0)u}\sigma_i\ell^{fir}$, pull back along the sheets $\beta_{T_i}s_{i,k'}$, and sum
over $i$ and $k'$. By (1) the right sides add up to zero. Since $e$ commutes with pull-backs, \cref{pf:steps1:latexp}(ii) gives
$\pi_W^*\mathrm{ev}_\nu(w_{i,\cdot},\ell^b\gamma)=\ell^{s+b}\sum_{k'}(\beta_{T_i}s_{i,k'})^*\mathrm{ev}_\nu(v,\gamma)$ for all coefficient
families $\gamma$; hence the pull-back to $W$ of $t^u\,\mathrm{ev}_\nu(w,c')$ vanishes, where $c'_{i,k}=\ell^{a+f(n+N_0)u+b}\sigma_i
\ell^{fir}\varphi^{f(n-i+N)}(c_k)$. By (2), $t^u\,\mathrm{ev}_\nu(w,C_0c')=0$ for all $\nu$, so by (a) $\mathrm{vec}_0(w,C_0c')=0$, and
\cref{pf:steps1:latexp}(iii) with $\kappa=C_0\ell^{a+f(n+N_0)u+b}\ne0$ gives the relation for $x$.
\end{proof}

\begin{theorem}[Crystalline congruence relation]\label{pf:steps1:finall}
Let $K\subset\Qlbar$ be finite over $\Ql$ with residue field of order $\ell^f$, $f>0$, and uniformizer $\varpi$; let $M$
be a smooth good model, $n,j\ge0$, and $T_0,\dots,T_n$ $\ell$-correspondences on $M.X$ over $\OO_K$ satisfying the
pointwise congruence relation for $(\ell^f,n,\varpi)$ (\cref{pf:steps1:def:cong}, with $\varpi$ pulled back to $M.X$); let
$A_i$ be the operator of $T_i$ on
$V_j(M)$. Then for every presentation $(C,\varphi)$ of $\stepsIBcris$ there is $N\ge0$ such that
\[
 \sum_{i=0}^n(-1)^i(\ell^f)^{\binom i2}\bigl(A_i\otimes\varphi^{f(n-i+N)}\bigr)(x)=0
 \qquad\text{for every $G_K$-invariant }x\in V_j(M)\otimes_{\Ql}C .
\]
\formallabel{UpperTheory.Ru.crystallineComparison.finAll}
\formallabel{UpperTheory.Ru.crystallineComparison.finAll_of}
\formallabel{UpperTheory.Ru.crystallineComparison.final}
\end{theorem}

\begin{proof}
A residue datum $(k,\rho_0,s_0)$ for $K$ (\cref{pf:steps1:cst}(iii)) exists by \lref{CrystallineSemistableComparison.resEx}
(every finite extension of $\Ql$ in $\Qlbar$ has one), and since $\#k=\ell^f$ every $x\in k$ satisfies $x^{\ell^f}=x$. The
comparison maps of the development have the properties collected in \lref{UpperTheory.Ru.crystallineComparison.cIn}
(deferred): the tautological line data and the $P$-map of \cref{pf:steps1:pseam}, a Frobenius factorisation of some level $N'$ with
$f\mid N'$ and the corresponding Teichm\"uller congruence, the \'etale side of smooth good models, the laws of the comparison data, a
comparison map of ``B-type'', and pinning conditions for every smooth good model and degree. (One of the formal statements attached
to the theorem has these properties as hypotheses.) We verify the hypotheses of \cref{pf:steps1:top}.

The model $M.X$ is projective over $\OO_K$, hence quasi-compact and separated, so it is covered by finitely many affine opens with
affine pairwise intersections. Let $e$ be the \'etale comparison family of the development,
which is good (\lref{CrystallineSemistableComparison.E.etMap}). Its restriction to the comparison data is pinned
(\lref{UpperTheory.Ru.crystallineComparison.ePin}, deferred, a repackaging), and the presentation $C$ embeds into the period ring
of the development by an injective map compatible with $c$, with $\Ql$ and with the Galois action
(\lref{UpperTheory.Ru.crystallineComparison.pbIota}, deferred). Hypothesis (a) of \cref{pf:steps1:top} then follows from the
pinning conditions and the B-type map (\lref{UpperTheory.Ru.crystallineComparison.injOf}, deferred), and hypothesis (b) from the
Frobenius factorisation, the Teichm\"uller congruence and the residue datum, which provide $g$ (a Teichm\"uller lift of the residue
map, with $g(\varpi)=0$), $F$ and $\rho$, together with the invariance of $x$ under $G_K$
(\lref{UpperTheory.Ru.crystallineComparison.absInvOf} and \lref{UpperTheory.Ru.crystallineComparison.pbInv}, deferred).
\cref{pf:steps1:top} gives the assertion.
\end{proof}

\begin{corollary}[Annihilation for tied packages]\label{pf:steps1:annmodel}
In the situation of \cref{pf:steps1:finall}, let $P$ be a tied period package, $\rho$ a representation of $G_K$ on $V_0$ and
$i\colon V_j(M)\xrightarrow\sim V_0$ an isomorphism intertwining the Galois actions, and let $T'_a=iA_ai^{-1}$. Then
$\mathrm{Ann}(\ell^f,n,T')$ holds for $\rho$ with $B=P.B$ and $\psi=\varphi^f$.
\formallabel{UpperTheory.Gal.crystallineComparison.Br.annModel}
\formallabel{UpperTheory.Gal.crystallineComparison.Br.annDrop}
\end{corollary}

\begin{proof}
Since $P$ is tied, $P.B$ with its Frobenius and Galois action is isomorphic to a presentation $(C,\varphi)$ of $\stepsIBcris$
with injective Frobenius (\lref{UpperTheory.Gal.crystallineComparison.Br.pBTied}, deferred: the isomorphism
$P.B\cong\stepsIBcris$ of the tie, composed with the presentation $\stepsIBcris=\stepsIAcris[1/t]$). By
\cref{pf:steps1:finall} there is $N$ with $(1\otimes\varphi^{fN})\bigl(\sum_a(-1)^a(\ell^f)^{\binom a2}(A_a\otimes\varphi^{f(n-a)})x\bigr)
=0$ for invariant $x$; as $1\otimes\varphi^{fN}$ is injective on $V_j\otimes C$ ($V_j$ is a free $\Ql$-module), the inner sum
vanishes. Transporting along $i$ and the isomorphism $P.B\cong C$ gives the claim.
\end{proof}

\subsubsection{The geometric side}

\begin{lemma}[Direct sums]\label{pf:steps1:annpi}
Let $P$ be a period package and $V_1,\dots,V_r$ finite-dimensional $\Ql$-representations of $G_{K}$ with endomorphisms
$T^{(i)}_j$. If every $V_i$ is crystalline for $P$, so is $\bigoplus_iV_i$; if $\mathrm{Ann}(q,n,T^{(i)})$ holds for every $i$
(for $(P.B,\varphi^f)$), then $\mathrm{Ann}(q,n,\bigoplus_iT^{(i)})$ holds for $\bigoplus_iV_i$.
\formallabel{UpperTheory.Ru.crystallineComparison.annPi}
\end{lemma}

\begin{proof}
The functor $V\mapsto D(V)=(V\otimes P.B)^{G_K}$ commutes with finite direct sums, compatibly with the operators $T_j\otimes1$ and
$1\otimes\varphi^f$; the dimension condition defining crystallinity and the annihilation relation are therefore checked summand by
summand.
\end{proof}

\begin{definition}[Crystalline embeddings with annihilator]\label{pf:steps1:def:incrys}
Let $P$ be a period package, $B_w$ a bridge for $G_w$, $f,q,n\in\N$, $\vartheta\colon\N\to\Qlbar$ and
$r_0\colon G_w\to\GL_a(\Qlbar)$. We say that $r_0$ has a \emph{crystalline embedding with annihilator
$(f,q,n,\vartheta)$} if there are a finite-dimensional $\Ql$-representation $V$ of $G_{K_w}$ which is crystalline for
$P$, endomorphisms $T_j$ of $V$ with $\mathrm{Ann}(q,n,T)$ for $(P.B,\varphi^f)$, and an injective $\Qlbar$-linear
$b\colon\Qlbar^a\to\Qlbar\otimes V$ with $b(r_0(g)x)=\rho_V(e(g))b(x)$ and $T_jb(x)=\vartheta_jb(x)$ for $j\le n$.
\formal{UpperTheory.Ru.crystallineComparison.InCrys}
\end{definition}

\begin{proposition}[The geometric side]\label{pf:steps1:geoside}
Let $H$ be a nice tower over a number field $E$, $\theta\colon\Gamma\to\Qlbar$, $\rho_0$ a matrix of the $\theta$-piece,
$w\mid\ell$, and suppose that $(H,\theta)$ has a congruence datum at $w$ with parameters $(n,\vartheta)$
(\cref{pf:steps1:def:cong}). Let $f_w$ be the residue degree of $w$. Then for every tied package $P$, place $A$ over $w$ and
inducing bridge $B_w$, $\rho_0|_{D_A}$ has a crystalline embedding with annihilator $(f_w,\stepsINw,n,\vartheta)$.
\formallabel{UpperTheory.Gal.Int.Geo.geoSide_of}
\formallabel{UpperTheory.Ru.crystallineComparison.geoLoc}
\formallabel{UpperTheory.Ru.crystallineComparison.congBC}
\end{proposition}

\begin{proof}
\emph{Base change of the congruence datum.} Let $C$ be the onto cover over $\stepsIOw$ with correspondences $T_0,\dots,T_n$,
indices $\gamma_j$, multiplicities $m_j$, irreducible $\varpi$, prime $p\in w$ and $\stepsINw=p^f$. Then $p=\ell$, since
$\ell\in w$ and two distinct primes generate the unit ideal; so $f=f_w$. By the integral embedding
\lref{UpperTheory.Gal.Seam.intMap} (an inducing bridge gives a ring map $\stepsIOw\to\OO_{K_w}$ compatible with $\tau_w$ and
preserving irreducible elements) we get $\phi\colon\stepsIOw\to\OO_{K_w}$ mapping $\varpi$ to a uniformizer, and the residue field
of $\OO_{K_w}$, which is that of the completion of $E$ at $w$, has $\ell^{f_w}$ elements.
The cover base-changes to an onto cover $C'$ over $\OO_{K_w}$ of $\mathrm{bc}_{\tau_w}H$ (\lref{UpperTheory.Gal.Seam.bcCover}),
the correspondences base-change to $\ell$-correspondences $T'_j$ on its model, still lying under the $\gamma_j$ with
multiplicities $m_j$, and the pointwise congruence relation is preserved, because $\phi(\varpi)$ divides $x^{\ell^f}-x$ for all
$x\in\OO_{K_w}$ and generic points of the special fibre base-change to generic points
(\lref{UpperTheory.Ru.crystallineComparison.ptIdBC}, deferred).

\emph{The local statement.} By the local tower lemma (\lref{UpperTheory.Gal.Seam.locTower_holds}) there is a matrix $\rho_1$
of the $\theta$-piece of $\mathrm{bc}_{\tau_w}H$ for $G_{K_w}$ with $\rho_1(e(g))=\rho_0(g)$, $g\in D_A$. Let $M$ be the model of
$C'$; it is smooth good (\lref{UpperTheory.Gal.Seam.goodOf}), and its $V_i(M)$ are finite-dimensional. For each $j$ let
$A_j$ be the operator of $T'_j$ on $V_i(M)$; it exists and satisfies $\stepsIhop(\gamma_j)\circ\pi^*=m_j\,\pi^*\circ A_j$ on
$V_i(M)$, where $\pi^*$ is pull-back to the tower (\lref{UpperTheory.Gal.crystallineComparison.Br.heckeCl}, deferred). The
piece is contained in the image of the injective equivariant map $\pi^*$ (\lref{UpperTheory.Gal.Seam.pieceIn_holds}), and,
as in the proof of \cref{pf:steps1:matin}, a basis of the piece has preimages in $\bigoplus_{i<N}\Qlbar\otimes V_i(M)$ for some
$N$; this defines an injective $b$ intertwining $\rho_1$ with the Galois action on $V=\bigoplus_{i<N}V_i(M)$ (written in bases).
Each $V_i(M)$ is crystalline for $P$ (\cref{pf:steps1:cst} and the admissibility transfer of \cref{sec:pf:steps1:cris}), and
satisfies $\mathrm{Ann}(\ell^{f},n,A)$ by \cref{pf:steps1:annmodel}; both properties pass to the direct sum $V$ with
$T_j=\bigoplus_iA_j$ by \cref{pf:steps1:annpi}. Finally, for $x$ in
the preimage of the piece, $m_j\pi^*(A_jx)=\stepsIhop(\gamma_j)\pi^*x=\theta(\gamma_j)\pi^*x=m_j\vartheta_j\pi^*x$, and $m_j\ne0$,
$\pi^*$ injective give $A_jx=\vartheta_jx$; hence $T_jb=\vartheta_jb$. Composing with $e$ gives the statement for $\rho_0|_{D_A}$.
\end{proof}

\subsubsection{From the geometric side to the Weil--Deligne representation}

\begin{proposition}[Graded pieces]\label{pf:steps1:joinq}
Let $P$ be a tied package, $B_w$ a bridge, $r_0,r\colon G_w\to\GL_a(\Qlbar)$ with $r$ continuous and graded of $r_0$, and suppose $r_0$
has a crystalline embedding with annihilator $(f,q,n,\vartheta)$. Then there are a finite extension $L\supseteq LL$ of
$\Ql$, a model $r'$ of $r$ over $L$ which is crystalline for $P$, and $\vartheta'\colon\N\to L$ with $\vartheta'_j=\vartheta_j$
for $j\le n$, such that $\stepsIHk_{q,n,\vartheta'}$ is nilpotent on $D(r')$ for the semistable period ring and $\psi=\varphi^f$.
\formallabel{UpperTheory.Ru.crystallineComparison.joinQ}
\formallabel{UpperTheory.Ru.crystallineComparison.JQ.nil_scalar}
\end{proposition}

\begin{proof}
Let $V,T_j,b$ be as in \cref{pf:steps1:def:incrys}, choose a basis of $V$, and let $\mathcal M\in M_{d\times a}(\Qlbar)$ be the
injective matrix of $b$, so that $R(e(g))\mathcal M=\mathcal Mr_0(g)$ for the matrices $R$ of $V$ and
$\mathrm{mat}(T_j)\mathcal M=\vartheta_j\mathcal M$ for $j\le n$. Let $L$ be the finite extension of $LL$ generated by the
entries of $\mathcal M$ and $\vartheta_0,\dots,\vartheta_n$, and $\mathcal N$ a left inverse of $\mathcal M$ over $L$. Then
$r_0(g)=\mathcal NR(e(g))\mathcal M$ has entries in $L$, and so does $r$; let $s_0,s$ be the models over $L$, and $U=L^a$
with $s_0$. Via $\mathcal M$, $U$ is an $L$-subrepresentation of $L\otimes V$ on which $T_j$ acts by $\vartheta_j$.

$V$ is crystalline, i.e.\ admissible for $P.B$; hence admissible for $\stepsIBst$ and $\mathrm{Ann}(q,n,T)$ holds for
$\stepsIBst$ (\lref{UpperTheory.Gal.crystallineComparison.stOfCr_holds}: crystalline representations are semistable and the
relation transfers along $\stepsIBcris\subset\stepsIBst$). Admissibility passes to $L\otimes V$ and to its subrepresentation $U$
(\lref{UpperTheory.Gal.crystallineComparison.admissible_sequence_exact}: subrepresentations and quotients of admissible
representations are admissible, by the dimension bound for invariants). Since $s$ is graded of $s_0$, the representation $W$
given by $s$ is, by the argument of \cref{pf:steps1:grcrys}, the direct sum of subquotients $C_v$ of $U$; these are admissible,
hence so is $W$, i.e.\ $r'=s$ is crystalline. On $D(U)\subseteq D(L\otimes V)$ the relation $\mathrm{Ann}(q,n,T)$ becomes
$\stepsIHk_{q,n,\vartheta'}(1\otimes\varphi^f)=0$, because $T_j=\vartheta_j$ on $U$; it follows that
$\stepsIHk_{q,n,\vartheta'}(1\otimes\varphi^f)$ is nilpotent on $D$ of every subquotient of $U$
(\lref{UpperTheory.Gal.crystallineComparison.joinAbs_holds}, deferred), and therefore on $D(W)=\bigoplus_vD(C_v)$, with the
maximum of the exponents. The continuity of $r$ is not used.
\end{proof}

\begin{proposition}[Frobenius eigenvalues]\label{pf:steps1:niltophi}
Let $P$ be tied, $\iota\colon\Qlbar\cong\C$, $A$ a place over $w$, $B_w$ an inducing bridge, $r\colon D_A\to\GL_a(\Qlbar)$ with a
crystalline model $r'$ over $L$, and $Q\in L[X]$ nilpotent on $D(r')$ for $(\stepsIBst,\varphi^{f_w})$. If
$\iota(Q)=u\,\stepsIrt_S$ with $u\in\C^\times$ and $S$ a multiset, then $r$ is crystalline for $(P,B_w)$, and there are $g\in W_A$
of degree $-1$ and $m\in\N$ such that the characteristic polynomial of $g$ on $\mathrm{WD}(r)$ divides $(\stepsIrt_S)^m$.
\formallabel{UpperTheory.Ru.crystallineComparison.nilToPhi}
\formallabel{UpperTheory.Gal.crystallineComparison.phiKNil_of}
\formallabel{UpperTheory.Gal.crystallineComparison.dvdOfNil_of}
\end{proposition}

\begin{proof}
Crystalline representations are semistable and hence potentially semistable
(\lref{GlobalTheory.CrysSst.isSst_of_isCrys}, \lref{GlobalTheory.GalCut.IwaTwist.isPst_of_isSst}). Let $g$ be the
inverse of a Frobenius element, so $\deg g=-1$. In Fontaine's construction the element $g$ acts on $D_{\mathrm{pst}}$, which
contains $\stepsIDst$, through the linearized Frobenius $\varphi^{f_w}$ (degree $-1$), so the nilpotence of $Q(\varphi^{f_w})$ on
$\stepsIDst$ gives the nilpotence of $Q$ evaluated at the operator of $g$
(\lref{UpperTheory.Gal.crystallineComparison.nilWeil_of}, deferred). This persists after extension of scalars to $\C$ by
$\iota$ (\lref{UpperTheory.Gal.crystallineComparison.nilRho_of}, deferred), and the characteristic polynomial of $g$ on
$\mathrm{WD}(r)$ equals that of $g$ on the Weil--Deligne representation of the model
(\lref{UpperTheory.Gal.crystallineComparison.wdModel_of}). Finally, if an endomorphism $X$ of a finite-dimensional vector space
over an algebraically closed field satisfies that $u\prod_{s\in S}(X-s)$ is nilpotent with $u\ne0$, then every eigenvalue of $X$ lies
in $S$, so its characteristic polynomial $\prod(X-\lambda_i)$ divides $(\stepsIrt_S)^m$ with $m$ the dimension.
\end{proof}

\begin{theorem}[$\sfname{PhiStep}$]\label{pf:steps1:phi}
Let $E,\Omega,\ell,\iota$ be as in \cref{sec:pf:steps1:conv:loc}, $H$ a nice tower over $E$, $\theta\colon\Gamma\to\Qlbar$,
$\rho\colon G_E\to\GL_d(\Qlbar)$ continuous and attached to the $\theta$-piece, $w\mid\ell$, and suppose that $(H,\theta)$ has a congruence
datum at $w$ with parameters $(n,\vartheta)$ and that $\iota(\stepsIHk_{\stepsINw,n,\vartheta})=\iota(\vartheta_0)\stepsIrt_R$ with
$\vartheta_0\ne0$. Then $\rho$ satisfies the Frobenius condition with roots $R$ at $w$: for every tied package $P$, place $A$ over $w$
and inducing bridge $B_w$, $\rho|_{D_A}$ is crystalline for $(P,B_w)$, and there are $g\in W_A$ of degree $-1$ and $m\in\N$ such that
the characteristic polynomial of $g$ on $\mathrm{WD}(\rho|_{D_A})$ divides $(\stepsIrt_R)^m$.
\formallabel{UpperTheory.Ru.crystallineComparison.phiStep}
\formallabel{UpperTheory.Ru.crystallineComparison.phiStep_of}
\end{theorem}

\begin{proof}
Let $\rho_0$ be a matrix of the piece with $\rho$ graded of $\rho_0$. By \cref{pf:steps1:geoside}, $\rho_0|_{D_A}$ has a crystalline embedding
with annihilator $(f_w,\stepsINw,n,\vartheta)$. Restriction to $D_A$ preserves the grading, so \cref{pf:steps1:joinq} gives $L$, a
crystalline model $r'$ of $\rho|_{D_A}$ over $L$ and $\vartheta'$ with $\vartheta'_j=\vartheta_j$ ($j\le n$) such that
$Q=\stepsIHk_{\stepsINw,n,\vartheta'}$ is nilpotent on $D(r')$. Since $Q$ depends only on $\vartheta'_0,\dots,\vartheta'_n$,
$\iota(Q)=\iota(\stepsIHk_{\stepsINw,n,\vartheta})=\iota(\vartheta_0)\stepsIrt_R$ with $\iota(\vartheta_0)\ne0$, and
\cref{pf:steps1:niltophi} gives the conclusion.
\end{proof}

\begin{remark}
The conclusion of $\sfname{PhiStep}$ is weaker than an equality of characteristic polynomials: it only places the eigenvalues of one
element of degree $-1$ among the roots of the Hecke polynomial. As in \cref{pf:steps1:crys}, the continuity of $\rho$ is not used in the
proof (it is a hypothesis of \cref{pf:steps1:joinq} that its proof does not use), and at places $w\nmid\ell$ the Frobenius condition
holds because no bridge induces a place over $w$.
\end{remark}
